% arara: pdflatex: {shell: yes}
% arara: bibtex
% arara: pdflatex: {shell: yes}
% arara: pdflatex: {shell: yes}
% arara: pdflatex: {shell: yes}
\documentclass[10pt]{article}
%   This program is free software: you can redistribute it and/or modify
%   it under the terms of the GNU General Public License as published by
%   the Free Software Foundation, either version 3 of the License, or
%   (at your option) any later version.
%
%   This program is distributed in the hope that it will be useful,
%   but WITHOUT ANY WARRANTY; without even the implied warranty of
%   MERCHANTABILITY or FITNESS FOR A PARTICULAR PURPOSE.  See the
%   GNU General Public License for more details.
%
%   See <http://www.gnu.org/licenses/>.
\usepackage[left=4.5cm,right=2.5cm,showframe=false,
	top=2cm,bottom=1.5cm]{geometry}                      % page setup
\usepackage{lmodern}
\usepackage{parskip}                                 % paragraph skips
\usepackage{booktabs}                                % beautiful tables
\usepackage{listings}                                % nice verbatim environments
\usepackage{titlesec}                                % customize headings
\usepackage{titletoc}                                % customize headings
\usepackage{multicol}
\usepackage{changepage}                              % adjust width of page
\usepackage{fancyhdr}                                % headers & footers
\usepackage{wrapfig}
\usepackage{fontawesome}
\usepackage[sc,format=hang,font=small]{caption}      % captions
\usepackage[backend=bibtex]{biblatex}                % bibliography
\usepackage{tcolorbox}                                % framed environments
\usepackage{xparse}
\usepackage[charter]{mathdesign}                     % changes font
\usepackage[expansion=false,kerning=true]{microtype} % better kerning
\usepackage{enumitem}                                % custom lists
\usepackage{longtable}
\usepackage{array}
% setup gitinfo2, as in the manual, details just above begin{document}
\usepackage[mark,grumpy]{gitinfo2}
% tcolorbox libraries
\tcbuselibrary{breakable,skins,listings,minted,xparse}
%\usepackage{varioref}                                % clever referencing
%\tcbuselibrary{documentation,breakable,skins,minted}
% tikz libraries
\usetikzlibrary{positioning}
\usetikzlibrary{decorations.pathmorphing}
\usetikzlibrary{decorations,shapes}
\usepackage{varioref}                                % clever referencing
\usepackage{hyperref}
\hypersetup{
	pdfauthor={Chris Hughes},
	pdftitle={latexindent.pl package},
	pdfkeywords={perl;beautify;indentation},
	bookmarksnumbered,
	bookmarksopen,
	linktocpage,
	colorlinks=true,
	linkcolor=blue,
	citecolor=black,
}
\usepackage{cleveref}

\addbibresource{latex-indent}
\addbibresource{contributors}

% http://tex.stackexchange.com/questions/122135/how-to-add-a-png-icon-on-the-right-side-of-a-tcolorbox-title
\newtcolorbox{warning}{parbox=false,breakable,enhanced,arc=0mm,colback=red!5,colframe=red,leftrule=12mm,%
	overlay={\node[anchor=north west,outer sep=2pt] at (frame.north west) {\includegraphics[width=8mm]{warning}}; }
}

\definecolor{harvestgold}{cmyk}{0.00, 0.05, 0.51, 0.07}  %EDE275
\definecolor{cmhgold}{cmyk}{0,0.178,0.909,0.008}         %FDD017
\makeatletter
\tcbset{
	addtolol/.style={list entry={\kvtcb@title},add to list={lol}{lstlisting}},
	cmhlistings/.style={
			%	width=\linewidth,
			%breakable,
			colback=blue!5!white,
			colframe=white!25!black,colback=white,
			top=0cm,
			bottom=0cm,
			left=0mm,
			listing only,
			center title,
			listing engine=minted,
			minted style=colorful,
			minted options={obeytabs=true,showtabs=true,tabsize=4,showspaces=true},
			addtolol,
			boxrule=0pt,
			toprule=1pt,bottomrule=1pt,
			titlerule=1pt,
			colframe=white!40!black,
			colback=white,
			sharp corners,
			colbacktitle=white!75!black
		},
	yaml-TCB/.style={
			listing only,
			listing engine=listings,
			left=0cm,
			boxrule=0pt,
			%leftrule=3pt,
			sharp corners,
			center title,
			%colbacktitle=white!75!black,
			colbacktitle=white!75!blue,
			colframe=white!25!blue,
			colback=white!90!blue,
			toprule=2pt,
			titlerule=2pt,
			%bottomrule=1pt,
		},
	MLB-TCB/.style={
			yaml-TCB,
			center title,
			colframe=cmhgold,
			colbacktitle=harvestgold,
			colback=white!60!cmhgold,
			width=0.9\linewidth,
			before=\centering,
			%bottomrule=1pt,
			enhanced,
			overlay={\node[anchor=north east,outer sep=2pt,draw=cmhgold,very thick,double,fill=harvestgold,font =\small] at ([yshift=-3mm]frame.north east) {\texttt{-m}}; }
		}
}

\newtcblisting[use counter=lstlisting]{cmhlistings}[3][]{%
	cmhlistings,
	center title,
	title={\color{black}{\scshape Listing \thetcbcounter}: ~#2},label={#3},
	listing engine=listings,
	listing options={#1},
}

\DeclareTCBInputListing[use counter=lstlisting]{\cmhlistingsfromfile}{O{} m O{} m m}{%
	cmhlistings,
	listing file={#2},
	listing options={#1},
	title={\color{black}{\scshape Listing \thetcbcounter}: ~#4},label={#5},
	#3,
}

% command shell
\newtcblisting{commandshell}{colback=black,colupper=white,colframe=yellow!75!black,
	listing only,listing options={style=tcblatex,language=sh,
			morekeywords={latexindent,pl},
			keywordstyle=\color{blue!35!white}\bfseries,
		},
	listing engine=listings,
	left=0cm,
	every listing line={\textcolor{red}{\small\ttfamily\fontseries{b}\selectfont cmh:$\sim$\$ }}}

% dosprompt
\newtcblisting{dosprompt}{
	colback=black,
	colupper=white,
	colframe=yellow!75!black,
	listing only,
	listing options={
			language=command.com,
			morekeywords={latexindent,pl},
			keywordstyle=\color{blue!35!white}\bfseries,
			basicstyle=\small\color{white}\ttfamily
		},
	listing engine=listings,
	left=0cm,
	every listing line={\textcolor{white}{\small\ttfamily\fontseries{b}\selectfont C:\textbackslash Users\textbackslash cmh$>$}}}

\lstset{%
	basicstyle=\small\ttfamily,language={[LaTeX]TeX},
	numberstyle=\ttfamily%\small,
	breaklines=true,
	keywordstyle=\color{blue},                    % keywords
	commentstyle=\color{purple},    % comments
	tabsize=3,
}%
\DeclareTCBListing[use counter=lstlisting]{yaml}{O{} m O{} m}{
	yaml-TCB,
	listing options={
			style=tcblatex,
			numbers=none,
			numberstyle=\color{red},
			#1,
		},
	title={\color{black}{\scshape Listing \thetcbcounter}: ~#2},label={#4},
	#3,
}

\lstdefinestyle{yaml-LST}{
	style=tcblatex,
	numbers=none,
	numberstyle=\color{red},
}

\lstdefinestyle{demo}{
	numbers=none,
	linewidth=1.25\textwidth,
	columns=fullflexible,
}

\lstdefinestyle{fileExtensionPreference}{
	style=yaml-LST,
	firstnumber=38,linerange={38-42},
	numbers=left,
}

\lstdefinestyle{logFilePreferences}{
	style=yaml-LST,
	firstnumber=79,linerange={79-83},numbers=left
}

\lstdefinestyle{verbatimEnvironments}{
	style=yaml-LST,
	firstnumber=87,linerange={87-90},numbers=left
}

\lstdefinestyle{verbatimCommands}{
	style=yaml-LST,
	firstnumber=93,linerange={93-95},numbers=left
}

\lstdefinestyle{noIndentBlock}{
	style=yaml-LST,
	firstnumber=101,linerange={101-103},numbers=left
}

\lstdefinestyle{removeTrailingWhitespace}{
	style=yaml-LST,
	firstnumber=106,linerange={106-108},numbers=left
}

\lstdefinestyle{fileContentsEnvironments}{
	style=yaml-LST,
	firstnumber=112,linerange={112-114},numbers=left
}

\lstdefinestyle{lookForPreamble}{
	style=yaml-LST,
	firstnumber=120,linerange={120-124},numbers=left
}

\lstdefinestyle{indentAfterItems}{
	style=yaml-LST,
	firstnumber=172,linerange={172-176},numbers=left
}

\lstdefinestyle{itemNames}{
	style=yaml-LST,
	firstnumber=182,linerange={182-184},numbers=left
}

\lstdefinestyle{specialBeginEnd}{
	style=yaml-LST,
	firstnumber=188,linerange={188-200},numbers=left
}

\lstdefinestyle{indentAfterHeadings}{
	style=yaml-LST,
	firstnumber=210,linerange={210-219},numbers=left
}

\lstdefinestyle{noAdditionalIndentGlobalEnv}{
	style=yaml-LST,
	firstnumber=265,linerange={265-266},numbers=left
}

\lstdefinestyle{noAdditionalIndentGlobal}{
	style=yaml-LST,
	firstnumber=265,linerange={265-277},numbers=left
}

\lstdefinestyle{indentRulesGlobalEnv}{
	style=yaml-LST,
	firstnumber=281,linerange={281-282},numbers=left
}

\lstdefinestyle{indentRulesGlobal}{
	style=yaml-LST,
	firstnumber=281,linerange={281-293},numbers=left
}

\lstdefinestyle{commandCodeBlocks}{
	style=yaml-LST,
	firstnumber=296,linerange={296-304},numbers=left
}

\lstdefinestyle{modifylinebreaks}{
	style=yaml-LST,
	firstnumber=374,linerange={374-376},numbers=left
}

\lstdefinestyle{textWrapOptions}{
	style=yaml-LST,
	firstnumber=377,linerange={377-378},numbers=left
}

\lstdefinestyle{textWrapOptionsAll}{
	style=yaml-LST,
	firstnumber=377,linerange={377-379},numbers=left
}

\lstdefinestyle{removeParagraphLineBreaks}{
	style=yaml-LST,
	firstnumber=380,linerange={380-392},numbers=left
}

\lstdefinestyle{paragraphsStopAt}{
	style=yaml-LST,
	firstnumber=393,linerange={393-401},numbers=left
}

\lstdefinestyle{modifylinebreaksEnv}{
	style=yaml-LST,
	firstnumber=402,linerange={402-411},numbers=left
}

% stars around contributors
\pgfdeclaredecoration{stars}{initial}{
	\state{initial}[width=15pt]
	{
		\pgfmathparse{round(rnd*100)}
		\pgfsetfillcolor{yellow!\pgfmathresult!orange}
		\pgfsetstrokecolor{yellow!\pgfmathresult!red}
		\pgfnode{star}{center}{}{}{\pgfusepath{stroke,fill}}
	}
	\state{final}
	{
		\pgfpathmoveto{\pgfpointdecoratedpathlast}
	}
}

\newtcolorbox{stars}{%
	enhanced jigsaw,
	breakable, % allow page breaks
	left=0cm,
	top=0cm,
	before skip=0.2cm,
	boxsep=0cm,
	frame style={draw=none,fill=none}, % hide the default frame
	colback=white,
	overlay={
			\draw[inner sep=0,minimum size=rnd*15pt+2pt]
			decorate[decoration={stars,segment length=2cm}] {
					decorate[decoration={zigzag,segment length=2cm,amplitude=0.3cm}] {
							([xshift=-.5cm,yshift=0.1cm]frame.south west) --  ([xshift=-.5cm,yshift=0.4cm]frame.north west)
						}};
			\draw[inner sep=0,minimum size=rnd*15pt+2pt]
			decorate[decoration={stars,segment length=2cm}] {
					decorate[decoration={zigzag,segment length=2cm,amplitude=0.3cm}] {
							([xshift=.75cm,yshift=0.1cm]frame.south east) --  ([xshift=.75cm,yshift=0.6cm]frame.north east)
						}};
			\node[anchor=north west,outer sep=2pt,opacity=0.25] at ([xshift=-4.25cm]frame.north west) {\resizebox{3cm}{!}{\faGithub}};
		},
	% paragraph skips obeyed within tcolorbox
	parbox=false,
}

% copied from /usr/local/texlive/2013/texmf-dist/tex/latex/biblatex/bbx/numeric.bbx
% the only modification is the \stars and \endstars
\defbibenvironment{specialbib}
{\stars\list
	{\printtext[labelnumberwidth]{%
			\printfield{prefixnumber}%
			\printfield{labelnumber}}}
	{\setlength{\labelwidth}{\labelnumberwidth}%
		\setlength{\leftmargin}{\labelwidth}%
		\setlength{\labelsep}{\biblabelsep}%
		\addtolength{\leftmargin}{\labelsep}%
		\setlength{\itemsep}{\bibitemsep}%
		\setlength{\parsep}{\bibparsep}}%
	\renewcommand*{\makelabel}[1]{\hss##1}}
{\endlist\endstars}
{\item}

\newtcbox{yamltitlebox}[1][]{colframe=black!50!white,boxrule=.5pt,sharp corners,#1}

\newcommand{\flagbox}[1]{%
	\par
	\makebox[30pt][l]{%
		\hspace{-2cm}%
		\ttfamily\fontseries{b}\selectfont #1
	}%
}

\NewDocumentCommand{\yamltitle}{O{} m s m}{%
	\par
	\makebox[30pt][l]{%
		\hspace{-2cm}%
		\yamltitlebox[#1]{%
			{\ttfamily\fontseries{b}\selectfont{\color{blue!80!white}#2}}: %
			\IfBooleanTF{#3}
			{$\langle$\itshape #4$\rangle$}
			{{\bfseries #4}}
		}}
	\par\nobreak%
}

\newcommand{\fixthis}[1]
{%
	\marginpar{\huge \color{red} \framebox{FIX}}%
	\typeout{FIXTHIS: p\thepage : #1^^J}%
}
% custom section
\titleformat{\section}
{\normalfont\Large\bfseries}
{\llap{\thesection\hskip.5cm}}
{0pt}
{}
% custom subsection
\titleformat{\subsection}
{\normalfont\bfseries}
{\llap{\thesubsection\hskip.5cm}}
{0pt}
{}
% custom subsubsection
\titleformat{\subsubsection}
{\normalfont\bfseries}
{\llap{\thesubsubsection\hskip.5cm}}
{0pt}
{}

\titlespacing\section{0pt}{12pt plus 4pt minus 2pt}{-5pt plus 2pt minus 2pt}
\titlespacing\subsection{0pt}{12pt plus 4pt minus 2pt}{-6pt plus 2pt minus 2pt}
\titlespacing\subsubsection{0pt}{12pt plus 4pt minus 2pt}{-6pt plus 2pt minus 2pt}

% list of listings
\contentsuse{lstlisting}{lol}
\titlecontents{lstlisting}[2em]
  {\addvspace{0.25pc}}
  {\textbf{Code \thecontentslabel} }
  {}
  {\titlerule*[0.5em]{$\cdot$}\contentspage}
  []
\AtBeginDocument{\addtocontents{lol}{\protect\begin{widepage}\protect\begin{multicols}{2}}}
\AtEndDocument{\addtocontents{lol}{\protect\end{multicols}\protect\end{widepage}}}

% cleveref settings
\crefname{table}{Table}{Tables}
\Crefname{table}{Table}{Tables}
\crefname{figure}{Figure}{Figures}
\Crefname{figure}{Figure}{Figures}
\crefname{section}{Section}{Sections}
\Crefname{section}{Section}{Sections}
\crefname{listing}{Listing}{Listings}
\Crefname{listing}{Listing}{Listings}

% headers and footers
\fancyhf{} % delete current header and footer
\fancyhead[R]{\bfseries\thepage%
	\tikz[remember picture,overlay] {
		\node at (1,0){\includegraphics{logo-bw}}; }
}
\fancyhead[L]{\rightmark}
\fancyheadoffset[L]{3cm}
\pagestyle{fancy}

% renew plain style
\fancypagestyle{plain}{%
	\fancyhf{} % clear all header and footer fields
	\renewcommand{\headrulewidth}{0pt}
	\renewcommand{\footrulewidth}{0pt}}

% widepage environment
\newenvironment{widepage}{\begin{adjustwidth}{-3cm}{0cm}}{\end{adjustwidth}}

% symbols for the m switch
\newcommand{\BeginStartsOnOwnLine}{\color{red}\spadesuit}
\newcommand{\BodyStartsOnOwnLine}{\color{red}\heartsuit}
\newcommand{\EndStartsOnOwnLine}{\color{red}\diamondsuit}
\newcommand{\EndFinishesWithLineBreak}{\color{red}\clubsuit}
\newcommand{\ElseStartsOnOwnLine}{\color{red}\bigstar}
\newcommand{\ElseFinishesWithLineBreak}{\color{red}\square}
\newcommand{\EqualsStartsOnOwnLine}{\color{red}\bullet}

% table rules
\setlength\heavyrulewidth{0.25ex}
% gitinfo2 settings
\renewcommand{\gitMark}{\gitBranch\,@\,\gitAbbrevHash{}\,\textbullet{}\,\gitAuthorDate\,\textbullet{}\faGithub}

% setting up gitinfo2:
%   copy the file post-xxx-sample.txt from http://mirror.ctan.org/macros/latex/contrib/gitinfo2
%   and put it in .git/hooks/post-checkout
% then
%   cd .git/hooks
%   chmod g+x post-checkout
%   chmod +x post-checkout
%   cp post-checkout post-commit
%   cp post-checkout post-merge
%   cd ../..
%   git checkout master
%   git checkout develop
%   ls .git
% and you should see gitHeadInfo.gin

% http://tex.stackexchange.com/questions/233843/textasteriskcentered-invisible-with-garamondmathdesign
% remove the definition of \textasteriskcentered for TS1 encoding
\UndeclareTextCommand{\textasteriskcentered}{TS1}
% reinstate a default encoding
\DeclareTextSymbolDefault{\textasteriskcentered}{OT1}
% suitably define the command
\DeclareTextCommand{\textasteriskcentered}{OT1}{\raisebox{-.7ex}[1ex][0pt]{*}}

\begin{document}
\renewcommand*{\thefootnote}{\arabic{footnote}}
\title{%
	\begin{tcolorbox}[
			width=5.2cm,
			boxrule=0pt,
			colframe=white!40!black,
			colback=white,
			rightrule=2pt,
			sharp corners,
			enhanced,
			overlay={\node[anchor=north east,outer sep=2pt] at ([xshift=3cm,yshift=4mm]frame.north east) {\includegraphics[width=3cm]{logo}}; }]
		\centering\ttfamily\bfseries latexindent.pl\\[1cm] Version 3.1
	\end{tcolorbox}
}
\author{Chris Hughes \thanks{and contributors! See \vref{sec:contributors}. For
		all communication, please visit \cite{latexindent-home}.}}
\maketitle
\begin{adjustwidth}{1cm}{1cm}
	\small
	\texttt{latexindent.pl} is a \texttt{Perl} script that indents \texttt{.tex} (and other)
	files according to an indentation scheme that the user can modify to suit their
	taste. Environments, including those with alignment delimiters (such as \texttt{tabular}),
	and commands, including those that can split braces and brackets across lines,
	are \emph{usually} handled correctly by the script. Options for \texttt{verbatim}-like
	environments and commands, together with indentation after headings (such as \lstinline!chapter!, \lstinline!section!, etc)
	are also available. The script also has the ability to modifiy line breaks, and add
	comment symbols. All user options are customisable via the switches in the YAML interface.
\end{adjustwidth}
\tableofcontents
{\small
	\lstlistoflistings
}

\section{Introduction}
\subsection{Thanks}
	I first created \texttt{latexindent.pl} to help me format chapter files
	in a big project. After I blogged about it on the
	\TeX{} stack exchange \cite{cmhblog} I received some positive feedback and
	follow-up feature requests. A big thank you to Harish Kumar who
	helped to develop and test the initial versions of the script.

	The \texttt{YAML}-based interface of \texttt{latexindent.pl} was inspired
	by the wonderful \texttt{arara} tool; any similarities are deliberate, and
	I hope that it is perceived as the compliment that it is. Thank you to Paulo Cereda and the
	team for releasing this awesome tool; I initially worried that I was going to
	have to make a GUI for \texttt{latexindent.pl}, but the release of \texttt{arara}
	has meant there is no need.

	There have been several contributors to the project so far (and hopefully more in
	the future!); thank you very much to the people detailed in \vref{sec:contributors}
	for their valued contributions, and thank you to those who report bugs and request features
	at \cite{latexindent-home}.

\subsection{License}
	\texttt{latexindent.pl} is free and open source, and it always will be.
	Before you start using it on any important files, bear in mind that \texttt{latexindent.pl} has the option to overwrite your \texttt{.tex} files.
	It will always make at least one backup (you can choose how many it makes, see \cpageref{page:onlyonebackup})
	but you should still be careful when using it. The script has been tested on many
	files, but there are some known limitations (see \cref{sec:knownlimitations}).
	You, the user, are responsible for ensuring that you maintain backups of your files
	before running \texttt{latexindent.pl} on them. I think it is important at this
	stage to restate an important part of the license here:
	\begin{quote}\itshape
		This program is distributed in the hope that it will be useful,
		but WITHOUT ANY WARRANTY; without even the implied warranty of
		MERCHANTABILITY or FITNESS FOR A PARTICULAR PURPOSE.  See the
		GNU General Public License for more details.
	\end{quote}
	There is certainly no malicious intent in releasing this script, and I do hope
	that it works as you expect it to; if it does not, please first of all
	make sure that you have the correct settings, and then feel free to let me know at \cite{latexindent-home} with a
	complete minimum working example as I would like to improve the code as much as possible.
	\begin{warning}
		Before you try the script on anything important (like your thesis), test it
		out on the sample files in the \texttt{test-case} directory \cite{latexindent-home}.
	\end{warning}

	\emph{If you have used any version 2.* of \texttt{latexindent.pl}, there
		are a few changes to the interface; see \vref{app:differences} and the comments
		throughout this document for details}.

\section{Demonstration: before and after}
 Let's give a demonstration of some before and after code -- after all, you probably
 won't want to try the script if you don't much like the results. You might also
 like to watch the video demonstration I made on youtube \cite{cmh:videodemo}

 As you look at \crefrange{lst:filecontentsbefore}{lst:pstricksafter}, remember
 that \texttt{latexindent.pl} is just following its rules, and there is nothing
 particular about these code snippets. All of the rules can be modified
 so that you can personalize your indentation scheme.

 In each of the samples given in \crefrange{lst:filecontentsbefore}{lst:pstricksafter}
 the `before' case is a `worst case scenario' with no effort to make indentation. The `after'
 result would be the same, regardless of the leading white space at the beginning of
 each line which is stripped by \texttt{latexindent.pl} (unless a \texttt{verbatim}-like
 environment or \texttt{noIndentBlock} is specified -- more on this in \cref{sec:defuseloc}).

 \begin{widepage}
	 \centering
	 \begin{minipage}{.45\linewidth}
		 \cmhlistingsfromfile{demonstrations/filecontents1.tex}{\texttt{filecontents1.tex}}{lst:filecontentsbefore}
	 \end{minipage}\hfill
	 \begin{minipage}{.45\linewidth}
		 \cmhlistingsfromfile{demonstrations/filecontents1-default.tex}{\texttt{filecontents1.tex} default output}{lst:filecontentsafter}
	 \end{minipage}%

	 \begin{minipage}{.45\linewidth}
		 \cmhlistingsfromfile{demonstrations/tikzset.tex}{\texttt{tikzset.tex}}{lst:tikzsetbefore}
	 \end{minipage}\hfill
	 \begin{minipage}{.45\linewidth}
		 \cmhlistingsfromfile{demonstrations/tikzset-default.tex}{\texttt{tikzset.tex} default output}{lst:tikzsetafter}
	 \end{minipage}%

	 \begin{minipage}{.45\linewidth}
		 \cmhlistingsfromfile{demonstrations/pstricks.tex}{\texttt{pstricks.tex}}{lst:pstricksbefore}
	 \end{minipage}\hfill
	 \begin{minipage}{.45\linewidth}
		 \cmhlistingsfromfile{demonstrations/pstricks-default.tex}{\texttt{pstricks.tex} default output}{lst:pstricksafter}
	 \end{minipage}%
 \end{widepage}


% arara: pdflatex: {shell: yes, files: [latexindent]}
\section{How to use the script}
 \texttt{latexindent.pl} ships as part of the \TeX Live distribution for
 Linux and Mac users; \texttt{latexindent.exe} ships as part of the \TeX Live
 and MiK\TeX{} distributions for Windows users. These files are also available
 from github \cite{latexindent-home} should you wish to use them without
 a \TeX{} distribution; in this case, you may like to read \vref{sec:updating-path}
 which details how the \texttt{path} variable can be updated.

 In what follows, we will always refer to \texttt{latexindent.pl}, but depending on
 your operating system and preference, you might substitute \texttt{latexindent.exe} or
 simply \texttt{latexindent}.

 There are two ways to use \texttt{latexindent.pl}: from the command line,
 and using \texttt{arara}; we discuss these in \cref{sec:commandline} and
 \cref{sec:arara} respectively. We will discuss how to change the settings and
 behaviour of the script in \vref{sec:defuseloc}.

 \texttt{latexindent.pl} ships with \texttt{latexindent.exe} for Windows
 users, so that you can use the script with or without a Perl distribution.
 If you plan to use \texttt{latexindent.pl} (i.e, the original Perl script) then you will
 need a few standard Perl modules -- see \vref{sec:requiredmodules} for details.

\subsection{From the command line}\label{sec:commandline}
	\texttt{latexindent.pl} has a number of different switches/flags/options, which
	can be combined in any way that you like, either in short or long form as detailed below.
	\texttt{latexindent.pl}  produces a \texttt{.log} file, \texttt{indent.log}, every time it
	is run; the name of the log file can be customized, but we will
	refer to the log file as \texttt{indent.log} throughout this document.
	There is a base of information that is written to \texttt{indent.log},
	but other additional information will be written depending
	on which of the following options are used.

	\begin{commandshell}
latexindent.pl
      \end{commandshell}

	This will output a welcome message to the terminal, including the version number
	and available options.

\flagbox{-h, --help}

	\begin{commandshell}
latexindent.pl -h
      \end{commandshell}

	As above this will output a welcome message to the terminal, including the version number
	and available options.
	\begin{commandshell}
latexindent.pl myfile.tex
      \end{commandshell}

	This will operate on \texttt{myfile.tex}, but will simply output to your terminal; \texttt{myfile.tex} will	not be changed
	by \texttt{latexindent.pl} in any way using this command.

\flagbox{-w, --overwrite}
	\begin{commandshell}
latexindent.pl -w myfile.tex
latexindent.pl --overwrite myfile.tex
latexindent.pl myfile.tex --overwrite 
      \end{commandshell}

	This \emph{will} overwrite \texttt{myfile.tex}, but it will
	make a copy of \texttt{myfile.tex} first. You can control the name of
	the extension (default is \texttt{.bak}), and how many different backups are made --
	more on this in \cref{sec:defuseloc}, and in particular see \texttt{backupExtension} and \texttt{onlyOneBackUp}.

	Note that if \texttt{latexindent.pl} can not create the backup, then it
	will exit without touching your original file; an error message will be given
	asking you to check the permissions of the backup file.

\flagbox{-o=output.tex,--outputfile=output.tex}
	\begin{commandshell} 
latexindent.pl -o=output.tex myfile.tex
latexindent.pl myfile.tex -o=output.tex 
latexindent.pl --outputfile=output.tex myfile.tex
latexindent.pl --outputfile output.tex myfile.tex
      \end{commandshell}

	This will indent \texttt{myfile.tex} and output it to \texttt{output.tex},
	overwriting it (\texttt{output.tex}) if it already exists\footnote{Users of version 2.* should
		note the subtle change in syntax}. Note that if \texttt{latexindent.pl} is called with both
	the \texttt{-w} and \texttt{-o} switches, then \texttt{-w} will
	be ignored and \texttt{-o} will take priority (this seems safer than the
	other way round).

	Note that using \texttt{-o} is equivalent to using
	\begin{commandshell}
latexindent.pl myfile.tex > output.tex
\end{commandshell}
	See \vref{app:differences} for details of how the interface has changed
	from Version 2.2 to Version 3.0 for this flag.
\flagbox{-s, --silent}
	\begin{commandshell}
latexindent.pl -s myfile.tex
latexindent.pl myfile.tex -s
      \end{commandshell}

	Silent mode: no output will be given to the terminal.

\flagbox{-t, --trace}
	\begin{commandshell}
latexindent.pl -t myfile.tex
latexindent.pl myfile.tex -t
      \end{commandshell}

	\label{page:traceswitch}
	Tracing mode: verbose output will be given to \texttt{indent.log}. This
	is useful if \texttt{latexindent.pl} has made a mistake and you're
	trying to find out where and why. You might also be interested in learning
	about \texttt{latexindent.pl}'s thought process -- if so, this
	switch is for you, although it should be noted that, especially for large files, this does affect
	performance of the script.

\flagbox{-tt, --ttrace}
	\begin{commandshell}
latexindent.pl -tt myfile.tex
latexindent.pl myfile.tex -tt
      \end{commandshell}

	\emph{More detailed} tracing mode: this option gives more details to \texttt{indent.log}
	than the standard \texttt{trace} option (note that, even more so than with \texttt{-t},
	especially for large files, performance of the script will be affected).

\flagbox{-l, --local[=myyaml.yaml,other.yaml,...]}
	\begin{commandshell}
latexindent.pl -l myfile.tex
latexindent.pl -l=myyaml.yaml myfile.tex
latexindent.pl -l myyaml.yaml myfile.tex
latexindent.pl -l first.yaml,second.yaml,third.yaml myfile.tex
latexindent.pl -l=first.yaml,second.yaml,third.yaml myfile.tex
latexindent.pl myfile.tex -l=first.yaml,second.yaml,third.yaml 
      \end{commandshell}

	\label{page:localswitch}
	\texttt{latexindent.pl} will always load \texttt{defaultSettings.yaml} (rhymes with camel)
	and if it is called with the \texttt{-l} switch and it finds \texttt{localSettings.yaml}
	in the same directory as \texttt{myfile.tex} then these settings will be
	added to the indentation scheme. Information will be given in \texttt{indent.log} on
	the success or failure of loading \texttt{localSettings.yaml}.

	The \texttt{-l} flag can take an \emph{optional} parameter which details the name (or names separated by commas) of a YAML file(s)
	that resides in the same directory as \texttt{myfile.tex}; you can use this option if you would
	like to load a settings file in the current working directory that is \emph{not} called \texttt{localSettings.yaml}.
	In fact, you can specify \emph{relative} path names to the current directory, but \emph{not}
	absolute paths -- for absolute paths, see \vref{sec:indentconfig}.
	Explicit demonstrations of how to use the \texttt{-l} switch are given throughout this documentation.

\flagbox{-d, --onlydefault}
	\begin{commandshell}
latexindent.pl -d myfile.tex
      \end{commandshell}

	Only \texttt{defaultSettings.yaml}: you might like to read \cref{sec:defuseloc} before
	using this switch. By default, \texttt{latexindent.pl} will always search for
	\texttt{indentconfig.yaml} or \texttt{.indentconfig.yaml}  in your home directory. If you would prefer it not to do so
	then (instead of deleting or renaming \texttt{indentconfig.yaml}/\texttt{.indentconfig.yaml}) you can simply
	call the script with the \texttt{-d} switch; note that this will also tell
	the script to ignore \texttt{localSettings.yaml} even if it has been called with the
	\texttt{-l} switch.

\flagbox{-c, --cruft=<directory>}
	\begin{commandshell}
latexindent.pl -c=/path/to/directory/ myfile.tex
      \end{commandshell}

	If you wish to have backup files and \texttt{indent.log} written to a directory
	other than the current working directory, then you can send these `cruft' files
	to another directory.
	% this switch was made as a result of http://tex.stackexchange.com/questions/142652/output-latexindent-auxiliary-files-to-a-different-directory

\flagbox{-g, --logfile}
	\begin{commandshell}
latexindent.pl -g=other.log myfile.tex
latexindent.pl -g other.log myfile.tex
latexindent.pl --logfile other.log myfile.tex
latexindent.pl myfile.tex -g other.log 
      \end{commandshell}

	By default, \texttt{latexindent.pl} reports information to \texttt{indent.log}, but if you wish to change the
	name of this file, simply call the script with your chosen name after the \texttt{-g} switch as demonstrated above.

\flagbox{-m, --modifylinebreaks}
	\begin{commandshell}
latexindent.pl -m myfile.tex
latexindent.pl -modifylinebreaks myfile.tex
      \end{commandshell}

	One of the most exciting developments in Version~3.0 is the ability to modify line breaks; for full details
	see \vref{sec:modifylinebreaks}

	\texttt{latexindent.pl} can also be called on a file without the file extension, for
	example
	\begin{commandshell}
latexindent.pl myfile
    \end{commandshell}
	and in which case, you can specify
	the order in which extensions are searched for; see \vref{lst:fileExtensionPreference}
	for full details.

\subsection{From \texttt{arara}}\label{sec:arara}
	Using \texttt{latexindent.pl} from the command line is fine for some folks, but
	others may find it easier to use from \texttt{arara}. \texttt{arara} ships with
	a rule, \texttt{indent.yaml}, but in case you do not have this rule, you can find it at \cite{paulo}.

	You can use the rule in any of the ways described in \cref{lst:arara} (or combinations thereof).
	In fact, \texttt{arara} allows yet greater flexibility -- you can use \texttt{yes/no}, \texttt{true/false}, or \texttt{on/off} to toggle the various options.
	\begin{cmhlistings}[style=demo,escapeinside={(*@}{@*)}]{\texttt{arara} sample usage}{lst:arara}
%(*@@*) arara: indent
%(*@@*) arara: indent: {overwrite: yes}
%(*@@*) arara: indent: {output: myfile.tex}
%(*@@*) arara: indent: {silent: yes}
%(*@@*) arara: indent: {trace: yes}
%(*@@*) arara: indent: {localSettings: yes}
%(*@@*) arara: indent: {onlyDefault: on}
%(*@@*) arara: indent: { cruft: /home/cmhughes/Desktop }
\documentclass{article}
...
\end{cmhlistings}
	%(*@@*) arara: indent: { modifylinebreaks: yes }

	Hopefully the use of these rules is fairly self-explanatory, but for completeness
	\cref{tab:orbsandswitches} shows the relationship between \texttt{arara} directive arguments and the
	switches given in \cref{sec:commandline}.

	\begin{table}[!ht]
		\centering
		\caption{\texttt{arara} directive arguments and corresponding switches}
		\label{tab:orbsandswitches}
		\begin{tabular}{lc}
			\toprule
			\texttt{arara} directive argument & switch      \\
			\midrule
			\texttt{overwrite}                & \texttt{-w} \\
			\texttt{output}                   & \texttt{-o} \\
			\texttt{silent}                   & \texttt{-s} \\
			\texttt{trace}                    & \texttt{-t} \\
			\texttt{localSettings}            & \texttt{-l} \\
			\texttt{onlyDefault}              & \texttt{-d} \\
			\texttt{cruft}                    & \texttt{-c} \\
			\texttt{modifylinebreaks}         & \texttt{-m} \\
			\bottomrule
		\end{tabular}
	\end{table}

	The \texttt{cruft} directive does not work well when used with
	directories that contain spaces.


% arara: pdflatex: {shell: yes, files: [latexindent]}
\section{User, local settings, \texttt{indentconfig.yaml} and \texttt{.indentconfig.yaml}}\label{sec:indentconfig}
 Editing \texttt{defaultSettings.yaml} is not ideal as it may be overwritten when
 updating your distribution--a better way to customize the settings to your liking
 is to set up your own settings file,
 \texttt{mysettings.yaml} (or any name you like, provided it ends with \texttt{.yaml}).
 The only thing you have to do is tell \texttt{latexindent.pl} where to find it.

 \texttt{latexindent.pl} will always check your home directory for \texttt{indentconfig.yaml}
 and  \texttt{.indentconfig.yaml} (unless
 it is called with the \texttt{-d} switch),
 which is a plain text file you can create that contains the \emph{absolute}
 paths for any settings files that you wish \texttt{latexindent.pl} to load. There is no difference
 between \texttt{indentconfig.yaml} and \texttt{.indentconfig.yaml}, other than the
 fact that \texttt{.indentconfig.yaml} is a `hidden' file; thank you to \cite{jacobo-diaz-hidden-config}
 for providing this feature. In what follows, we will use \texttt{indentconfig.yaml}, but it
 is understood that this equally represents \texttt{.indentconfig.yaml} as well. If you
 have both files in existence,  \texttt{indentconfig.yaml} takes priority.

 For Mac and Linux users, their home directory is \texttt{~/username} while
 Windows (Vista onwards) is \lstinline!C:\Users\username!\footnote{If you're not sure
	 where to put \texttt{indentconfig.yaml}, don't
	 worry \texttt{latexindent.pl} will tell you in the log file exactly where to
	 put it assuming it doesn't exist already.}
 \Cref{lst:indentconfig} shows a sample \texttt{indentconfig.yaml} file.

 \begin{yaml}{\texttt{indentconfig.yaml} (sample)}{lst:indentconfig}
	# Paths to user settings for latexindent.pl
	#
	# Note that the settings will be read in the order you
	# specify here- each successive settings file will overwrite
	# the variables that you specify

	paths:
	- /home/cmhughes/Documents/yamlfiles/mysettings.yaml
	- /home/cmhughes/folder/othersettings.yaml
	- /some/other/folder/anynameyouwant.yaml
	- C:\Users\chughes\Documents\mysettings.yaml
	- C:\Users\chughes\Desktop\test spaces\more spaces.yaml
\end{yaml}

 Note that the \texttt{.yaml} files you specify in \texttt{indentconfig.yaml}
 will be loaded in the order that you write them in. Each file doesn't have
 to have every switch from \texttt{defaultSettings.yaml}; in fact, I recommend
 that you only keep the switches that you want to \emph{change} in these
 settings files.

 To get started with your own settings file, you might like to save a copy of
 \texttt{defaultSettings.yaml} in another directory and call it, for
 example, \texttt{mysettings.yaml}. Once you have added the path to \texttt{indentconfig.yaml}
 you can change the switches and add more code-block names to it
 as you see fit -- have a look at \cref{lst:mysettings} for an example
 that uses four tabs for the default indent, adds the \texttt{tabbing}
 environment/command to the list of environments that contains alignment delimiters; you might also like to
 refer to the many YAML files detailed throughout the rest of this documentation.

 \begin{yaml}{\texttt{mysettings.yaml} (example)}{lst:mysettings}
# Default value of indentation
defaultIndent: "\t\t\t\t"

# environments that have tab delimiters, add more
# as needed
lookForAlignDelims:
    tabbing: 1
\end{yaml}

 You can make sure that your settings are loaded by checking \texttt{indent.log}
 for details -- if you have specified a path that \texttt{latexindent.pl} doesn't
 recognize then you'll get a warning, otherwise you'll get confirmation that
 \texttt{latexindent.pl} has read your settings file \footnote{Windows users
	 may find that they have to end \texttt{.yaml} files with a blank line}.

 \begin{warning}
	 When editing \texttt{.yaml} files it is \emph{extremely} important
	 to remember how sensitive they are to spaces. I highly recommend copying
	 and pasting from \texttt{defaultSettings.yaml} when you create your
	 first \texttt{whatevernameyoulike.yaml} file.

	 If \texttt{latexindent.pl} can not read your \texttt{.yaml} file it
	 will tell you so in \texttt{indent.log}.
 \end{warning}

\subsection{\texttt{localSettings.yaml}}\label{sec:localsettings}
	The \texttt{-l} switch tells \texttt{latexindent.pl} to look for \texttt{localSettings.yaml} in the
	\emph{same directory} as \texttt{myfile.tex}.  If you'd prefer to name your \texttt{localSettings.yaml} file something
	different, (say, \texttt{myyaml.yaml}) then
	you can call \texttt{latexindent.pl} using, for example,
	\begin{commandshell}
latexindent.pl -l=myyaml.yaml myfile.tex
\end{commandshell}

	Any settings file(s) specified using the \texttt{-l} switch will be read \emph{after} \texttt{defaultSettings.yaml} and, assuming they exist,
	user settings from \texttt{indentconfig.yaml}.

	Your settings file can contain any switches that you'd
	like to change; a sample is shown in \cref{lst:localSettings}, and you'll find plenty of further examples throughout this manual.

	\begin{yaml}{\texttt{localSettings.yaml} (example)}{lst:localSettings}
#  verbatim environments- environments specified
#  in this hash table will not be changed at all!
verbatimEnvironments:
    cmhenvironment: 0
\end{yaml}

	You can make sure that your settings file has been loaded by checking \texttt{indent.log}
	for details; if it can not be read then you receive a warning, otherwise you'll get confirmation that
	\texttt{latexindent.pl} has read your settings file.

\subsection{Settings load order}\label{sec:loadorder}
	\texttt{latexindent.pl} loads the settings files in the following order:
	\begin{enumerate}
		\item \texttt{defaultSettings.yaml} is always loaded, and can not be renamed;
		\item \texttt{anyUserSettings.yaml} and any other arbitrarily-named files specified in \texttt{indentconfig.yaml};
		\item \texttt{localSettings.yaml} but only if found in the same directory as \texttt{myfile.tex} and called
		      with \texttt{-l} switch; this file can be renamed, provided that the call to \texttt{latexindent.pl} is adjusted
		      accordingly (see \cref{sec:localsettings}). You may specify relative  paths to other
		      YAML files using the \texttt{-l} switch, separating multiple files using commas.
	\end{enumerate}
	A visual representation of this is given in \cref{fig:loadorder}.

	\begin{figure}
		\centering
		\begin{tikzpicture}[
				needed/.style={very thick, draw=blue,fill=blue!20,
						text centered, minimum height=2.5em,rounded corners=1ex},
				optional/.style={draw=black, very thick,scale=0.8,
						text centered, minimum height=2.5em,rounded corners=1ex},
				optionalfill/.style={fill=black!10},
				connections/.style={draw=black!30,dotted,line width=3pt,text=red},
			]
			% Draw diagram elements
			\node (latexindent) [needed,circle]  {\texttt{latexindent.pl}};
			\node (default) [needed,above right=.5cm of latexindent]  {\texttt{defaultSettings.yaml}};
			\node (indentconfig) [optional,right=of latexindent]  {\texttt{indentconfig.yaml}};
			\node (any) [optional,optionalfill,above right=of indentconfig]  {\texttt{any.yaml}};
			\node (name) [optional,optionalfill,right=of indentconfig]  {\texttt{name.yaml}};
			\node (you) [optional,optionalfill,below right=of indentconfig]  {\texttt{you.yaml}};
			\node (want) [optional,optionalfill,below=of indentconfig]  {\texttt{want.yaml}};
			\node (local) [optional,below=of latexindent]  {\texttt{localSettings.yaml}};
			% Draw arrows between elements
			\draw[connections,solid] (latexindent) to[in=-90]node[pos=0.5,anchor=north]{1} (default.south) ;
			\draw[connections,optional] (latexindent) -- node[pos=0.5,anchor=north]{2} (indentconfig) ;
			\draw[connections,optional] (indentconfig) to[in=-90] (any.south) ;
			\draw[connections,optional] (indentconfig) -- (name) ;
			\draw[connections,optional] (indentconfig) to[out=-45,in=90] (you) ;
			\draw[connections,optional] (indentconfig) -- (want) ;
			\draw[connections,optional] (latexindent) -- node[pos=0.5,anchor=west]{3} (local) ;
		\end{tikzpicture}
		\caption{Schematic of the load order described in \cref{sec:loadorder}; solid lines represent
			mandatory files, dotted lines represent optional files. \texttt{indentconfig.yaml} can
			contain as many files as you like. The files will be loaded in order; if you specify
			settings for the same field in more than one file, the most recent takes priority. }
		\label{fig:loadorder}
	\end{figure}

% arara: pdflatex: {shell: yes, files: [latexindent]}
\section{defaultSettings.yaml}\label{sec:defuseloc}
 \texttt{latexindent.pl} loads its settings from \texttt{defaultSettings.yaml}. The idea is to separate the behaviour of the script
 from the internal working -- this is very similar to the way that we separate content
 from form when writing our documents in \LaTeX.

 If you look in \texttt{defaultSettings.yaml} you'll find the switches
 that govern the behaviour of \texttt{latexindent.pl}. If you're not sure where
 \texttt{defaultSettings.yaml} resides on your computer, don't worry as \texttt{indent.log}
 will tell you where to find it.
 \texttt{defaultSettings.yaml} is commented,
 but here is a description of what each switch is designed to do. The default
 value is given in each case; whenever you see \emph{integer} in \emph{this}
 section, assume that it must be greater than or equal to \texttt{0} unless
 otherwise stated.

\yamltitle{fileExtensionPreference}*{fields}
	\texttt{latexindent.pl} can be called to
	act on a file without
	specifying the file extension.  For example we can call
	\begin{commandshell}
latexindent.pl myfile
\end{commandshell}
	\begin{wrapfigure}[8]{r}[0pt]{6cm}
		\cmhlistingsfromfile[style=fileExtensionPreference]{../defaultSettings.yaml}[width=.8\linewidth,before=\centering,yaml-TCB]{\texttt{fileExtensionPreference}}{lst:fileExtensionPreference}
	\end{wrapfigure}

	in which case the script will look for \texttt{myfile} with the extensions
	specified in \texttt{fileExtensionPreference} in their numeric order. If
	no match is found, the script will exit. As with all of the fields, you should
	change and/or add to this as necessary.

	Calling \texttt{latexindent.pl myfile} with the (default) settings specified in \cref{lst:fileExtensionPreference}
	means that the script will first look for \texttt{myfile.tex}, then \texttt{myfile.sty}, \texttt{myfile.cls},
	and finally \texttt{myfile.bib} in order\footnote{Throughout this manual, listings with line numbers represent code
		taken directly from \texttt{defaultSettings.yaml}.}.

\yamltitle{backupExtension}*{extension name}

	If you call \texttt{latexindent.pl} with the \texttt{-w} switch (to overwrite
	\texttt{myfile.tex}) then it will create a backup file before doing
	any indentation; the default extension is \texttt{.bak}, so, for example, \texttt{myfile.bak0}
	would be created when calling \texttt{latexindent.pl myfile.tex}  for the first time.

	By default, every time you subsequently call \texttt{latexindent.pl} with
	the \texttt{-w} to act upon \texttt{myfile.tex}, it will create successive back up files: \texttt{myfile.bak1}, \texttt{myfile.bak2},
	etc.

\yamltitle{onlyOneBackUp}*{integer}
	\label{page:onlyonebackup}
	If you don't want a backup for every time that you call \texttt{latexindent.pl} (so
	you don't want \texttt{myfile.bak1}, \texttt{myfile.bak2}, etc) and you simply
	want \texttt{myfile.bak} (or whatever you chose \texttt{backupExtension} to be)
	then change \texttt{onlyOneBackUp} to \texttt{1}; the default value of
	\texttt{onlyOneBackUp} is \texttt{0}.

\yamltitle{maxNumberOfBackUps}*{integer}
	Some users may only want a finite number of backup files,
	say at most $3$, in which case, they can change this switch.
	The smallest value of \texttt{maxNumberOfBackUps} is $0$ which will \emph{not}
	prevent backup files being made; in this case, the behaviour will be dictated
	entirely by \texttt{onlyOneBackUp}. The default value of \texttt{maxNumberOfBackUps}
	is \texttt{0}.

\yamltitle{cycleThroughBackUps}*{integer}
	Some users may wish to cycle through backup files, by deleting the
	oldest backup file and keeping only the most recent; for example,
	with \texttt{maxNumberOfBackUps: 4}, and \texttt{cycleThroughBackUps}
	set to \texttt{1}  then the \texttt{copy} procedure given below
	would be obeyed.

	\begin{commandshell}
copy myfile.bak1 to myfile.bak0
copy myfile.bak2 to myfile.bak1
copy myfile.bak3 to myfile.bak2
copy myfile.bak4 to myfile.bak3
	\end{commandshell}
	The default value of \texttt{cycleThroughBackUps} is \texttt{0}.

\yamltitle{logFilePreferences}*{fields}
	\begin{wrapfigure}[10]{r}[0pt]{9cm}
		\cmhlistingsfromfile[style=logFilePreferences,]{../defaultSettings.yaml}[width=.85\linewidth,before=\centering,yaml-TCB]{\texttt{logFilePreferences}}{lst:logFilePreferences}
	\end{wrapfigure}
	\texttt{latexindent.pl} writes information to \texttt{indent.log}, some
	of which can be customized by changing \texttt{logFilePreferences}; see \cref{lst:logFilePreferences}.
	If you load your own user settings (see \vref{sec:indentconfig}) then \texttt{latexindent.pl} will
	detail them in \texttt{indent.log}; you can choose not to have the details logged by switching
	\texttt{showEveryYamlRead} to \texttt{0}. Once all of your settings have
	been loaded, you can see the amalgamated settings in the log file by switching \texttt{showAmalgamatedSettings}
	to \texttt{1}, if you wish. The log file will end with the characters
	given in \texttt{endLogFileWith}, and will report the \texttt{GitHub} address
	of \texttt{latexindent.pl} to the log file if \texttt{showGitHubInfoFooter} is set to \texttt{1}.

\yamltitle{verbatimEnvironments}*{fields}

	A field that contains a list of environments
	that you would like left completely alone -- no indentation will be performed
	on environments that you have specified in this field, see \cref{lst:verbatimEnvironments}.

	\begin{minipage}{.45\textwidth}
		\cmhlistingsfromfile[style=verbatimEnvironments]{../defaultSettings.yaml}[width=.8\linewidth,before=\centering,yaml-TCB]{\texttt{verbatimEnvironments}}{lst:verbatimEnvironments}
	\end{minipage}%
	\begin{minipage}{.45\textwidth}
		\cmhlistingsfromfile[style=verbatimCommands]{../defaultSettings.yaml}[width=.8\linewidth,before=\centering,yaml-TCB]{\texttt{verbatimCommands}}{lst:verbatimCommands}
	\end{minipage}%

	Note that if  you put an environment in  \texttt{verbatimEnvironments}
	and in other fields such as \texttt{lookForAlignDelims} or \texttt{noAdditionalIndent}
	then \texttt{latexindent.pl} will \emph{always} prioritize  \texttt{verbatimEnvironments}.

\yamltitle{verbatimCommands}*{fields}
	A field that contains a list of commands that are verbatim commands, for example
	\lstinline|\lstinline|; any commands populated in this field are protected from line breaking
	routines (only relevant if the \texttt{-m} is active, see \vref{sec:modifylinebreaks}).

\yamltitle{noIndentBlock}*{fields}

	\begin{wrapfigure}[8]{r}[0pt]{6cm}
		\cmhlistingsfromfile[style=noIndentBlock]{../defaultSettings.yaml}[width=.8\linewidth,before=\centering,yaml-TCB]{\texttt{noIndentBlock}}{lst:noIndentBlock}
	\end{wrapfigure}
	If you have a block of code that you don't want \texttt{latexindent.pl} to touch (even if it is \emph{not} a verbatim-like
	environment) then you can wrap it in an environment from \texttt{noIndentBlock};
	you can use any name you like for this, provided you populate it as demonstrate in
	\cref{lst:noIndentBlock}.

	Of course, you don't want to have to specify these as null environments
	in your code, so you use them with a comment symbol, \lstinline!%!, followed
	by as many spaces (possibly none) as you like; see \cref{lst:noIndentBlockdemo} for
	example.

	\begin{cmhlistings}[style=demo,escapeinside={(*@}{@*)}]{\texttt{noIndentBlock} demonstration}{lst:noIndentBlockdemo}
%(*@@*) \begin{noindent}
        this code
                won't
     be touched
                    by
             latexindent.pl!
%(*@@*)\end{noindent}
	\end{cmhlistings}

\yamltitle{removeTrailingWhitespace}*{fields}\label{yaml:removeTrailingWhitespace}

	\begin{wrapfigure}[12]{r}[0pt]{6cm}
		\cmhlistingsfromfile[style=removeTrailingWhitespace]{../defaultSettings.yaml}[width=.8\linewidth,before=\centering,yaml-TCB]{removeTrailingWhitespace}{lst:removeTrailingWhitespace}

		\vspace{.2cm}
		\cmhlistingsfromfile[style=fileContentsEnvironments]{../defaultSettings.yaml}[width=.8\linewidth,before=\centering,yaml-TCB]{\texttt{fileContentsEnvironments}}{lst:fileContentsEnvironments}
	\end{wrapfigure}
	Trailing white space can be removed both \emph{before} and \emph{after} processing
	the document, as detailed in \cref{lst:removeTrailingWhitespace}; each of the fields
	can take the values \texttt{0} or \texttt{1}. See \vref{lst:removeTWS-before,lst:env-mlb5-modAll,lst:env-mlb5-modAll-remove-WS}
	for before and after results.  Thanks to \cite{vosskuhle} for providing this feature.

\yamltitle{fileContentsEnvironments}*{field}

	Before \texttt{latexindent.pl} determines the difference between preamble (if any) and the main document,
	it first searches for any of the environments specified in \texttt{fileContentsEnvironments}, see
	\cref{lst:fileContentsEnvironments}.
	The behaviour of \texttt{latexindent.pl} on these environments is determined by their location (preamble or not), and
	the value \texttt{indentPreamble}, discussed next.

\yamltitle{indentPreamble}{0|1}

	The preamble of a document can sometimes contain some trickier code
	for \texttt{latexindent.pl} to operate upon. By default, \texttt{latexindent.pl}
	won't try to operate on the preamble (as \texttt{indentPreamble} is set to \texttt{0},
	by default), but if you'd like \texttt{latexindent.pl} to try then change \texttt{indentPreamble} to \texttt{1}.

\yamltitle{lookForPreamble}*{fields}

	\begin{wrapfigure}[8]{r}[0pt]{5cm}
		\cmhlistingsfromfile[style=lookForPreamble]{../defaultSettings.yaml}[width=.8\linewidth,before=\centering,yaml-TCB]{lookForPreamble}{lst:lookForPreamble}
	\end{wrapfigure}
	Not all files contain preamble; for example, \texttt{sty}, \texttt{cls} and \texttt{bib} files typically do \emph{not}. Referencing
	\cref{lst:lookForPreamble}, if you set, for example, \texttt{.tex} to \texttt{0}, then regardless of the setting of the value of \texttt{indentPreamble}, preamble
	will not be assumed when operating upon \texttt{.tex} files.
\yamltitle{preambleCommandsBeforeEnvironments}{0|1}
	Assuming that \texttt{latexindent.pl} is asked to operate upon the preamble of a document,
	when this switch is set to \texttt{0} then environment code blocks will be sought first,
	and then command code blocks. When this switch is set to \texttt{1}, commands
	will be sought first. The example that first motivated this switch contained the code given in \cref{lst:motivatepreambleCommandsBeforeEnvironments}.

	\begin{cmhlistings}{Motivating \texttt{preambleCommandsBeforeEnvironments}}{lst:motivatepreambleCommandsBeforeEnvironments}
...
preheadhook={\begin{mdframed}[style=myframedstyle]},
postfoothook=\end{mdframed},
...
\end{cmhlistings}

\yamltitle{defaultIndent}*{horizontal space}
	This is the default indentation (\lstinline!\t! means a tab, and is the default value) used in the absence of other details
	for the command or environment we are working with; see \texttt{indentRules} in \vref{sec:noadd-indent-rules}
	for more details.

	If you're interested in experimenting with \texttt{latexindent.pl} then you
	can \emph{remove} all indentation by setting \texttt{defaultIndent: ""}.

\yamltitle{lookForAlignDelims}*{fields}
	\begin{wrapfigure}[12]{r}[0pt]{5cm}
		\begin{yaml}[numbers=none]{\texttt{lookForAlignDelims} (basic)}[width=.8\linewidth,before=\centering]{lst:aligndelims:basic}
lookForAlignDelims:
   tabular: 1
   tabularx: 1
   longtable: 1
   array: 1
   matrix: 1
   ...
	\end{yaml}
	\end{wrapfigure}
	This contains a list of environments and/or commands that
	are operated upon in a special way by \texttt{latexindent.pl} (see \cref{lst:aligndelims:basic}).
	In fact, the fields in \texttt{lookForAlignDelims} can actually
	take two different forms: the \emph{basic} version is shown in \cref{lst:aligndelims:basic}
	and the \emph{advanced} version in \cref{lst:aligndelims:advanced}; we will discuss each in turn.

	The environments specified in this field will be operated on in a special way  by \texttt{latexindent.pl}. In particular, it will try and align each column by its alignment
	tabs. It does have some limitations (discussed further in \cref{sec:knownlimitations}),
	but in many cases it will produce results such as those in \cref{lst:tabularbefore:basic,lst:tabularafter:basic}.

	If you find that \texttt{latexindent.pl} does not perform satisfactorily on such
	environments then you can set the relevant key to \texttt{0}, for example \texttt{tabular: 0}; alternatively, if you just want to ignore \emph{specific}
	instances of the environment, you could wrap them in something from \texttt{noIndentBlock} (see \cref{lst:noIndentBlock}).

	\begin{minipage}{.45\textwidth}
		\cmhlistingsfromfile[columns=fixed]{demonstrations/tabular1.tex}{\texttt{tabular1.tex}}{lst:tabularbefore:basic}
	\end{minipage}%
	\hfill
	\begin{minipage}{.45\textwidth}
		\cmhlistingsfromfile[columns=fixed]{demonstrations/tabular1-default.tex}{\texttt{tabular1.tex} default output}{lst:tabularafter:basic}
	\end{minipage}%

	If you wish to remove the alignment of the \lstinline!\\! within a delimiter-aligned block, then the
	advanced form of \texttt{lookForAlignDelims} shown in \cref{lst:aligndelims:advanced} is for you.

	\cmhlistingsfromfile[style=yaml-LST]{demonstrations/tabular.yaml}[yaml-TCB]{\texttt{tabular.yaml}}{lst:aligndelims:advanced}

	Note that you can use a mixture of the basic and advanced form: in \cref{lst:aligndelims:advanced} \texttt{tabular} and \texttt{tabularx}
	are advanced and \texttt{longtable} is basic. When using the advanced form, each field should receive at least 1 sub-field, and \emph{can} (but does not have to) receive up to 3 fields:
	\begin{itemize}
		\item \texttt{delims}: switch equivalent to simply specifying, for example, \texttt{tabular: 1} in
		      the basic version shown in \cref{lst:aligndelims:basic} (default: 1);
		\item \texttt{alignDoubleBackSlash}: switch to determine if \lstinline!\\! should be aligned (default: 1);
		\item \texttt{spacesBeforeDoubleBackSlash}: optionally, specifies the number of spaces to be inserted
		      before (non-aligned) \lstinline!\\!. In order to use this field, \texttt{alignDoubleBackSlash} needs
		      to be set to 0 (default: 0).
	\end{itemize}

	Assuming that you have the settings in \cref{lst:aligndelims:advanced} saved in \texttt{tabular.yaml}, and the code
	from \cref{lst:tabularbefore:basic} in \texttt{tabular1.tex} and you run
	\begin{commandshell}
latexindent.pl -l tabular.yaml tabular1.tex 
\end{commandshell}
	then you should receive the before-and-after results shown in
	\cref{lst:tabularbefore:advanced,lst:tabularafter:advanced}; note that the ampersands have been aligned, but
	the \lstinline!\\! have not (compare the alignment of \lstinline!\\! in \cref{lst:tabularafter:basic,lst:tabularafter:advanced}).

	\begin{minipage}{.45\textwidth}
		\cmhlistingsfromfile[columns=fixed]{demonstrations/tabular1.tex}{\texttt{tabular1.tex}}{lst:tabularbefore:advanced}
	\end{minipage}%
	\hfill
	\begin{minipage}{.45\textwidth}
		\cmhlistingsfromfile[columns=fixed]{demonstrations/tabular1-advanced.tex}{\texttt{tabular1.tex} using \cref{lst:aligndelims:advanced}}{lst:tabularafter:advanced}
	\end{minipage}%

	Saving \cref{lst:aligndelims:advanced} into \texttt{tabular1.yaml} as in \cref{lst:tabular1YAML}, and running the command
	\begin{commandshell}
latexindent.pl -l tabular1.yaml tabular1.tex 
\end{commandshell}
	gives \cref{lst:tabularafter:spacing}; note the spacing before the \lstinline!\\!.

	\begin{minipage}{.45\textwidth}
		\cmhlistingsfromfile[columns=fixed]{demonstrations/tabular1-advanced-3spaces.tex}{\texttt{tabular1.tex} using \cref{lst:tabular1YAML}}{lst:tabularafter:spacing}
	\end{minipage}%
	\hfill
	\begin{minipage}{.54\textwidth}
		\cmhlistingsfromfile[style=yaml-LST]{demonstrations/tabular1.yaml}[yaml-TCB]{\texttt{tabular1.yaml}}{lst:tabular1YAML}
	\end{minipage}%

	As of Version 3.0, the alignment routine works on mandatory and optional arguments within commands, and also within `special' code blocks
	(see \texttt{specialBeginEnd} on \cpageref{yaml:specialBeginEnd}); for example, assuming that you have a command called \lstinline!\matrix!
	and that it is populated within \texttt{lookForAlignDelims} (which it is, by default), and that you
	run the command
	\begin{commandshell}
latexindent.pl matrix1.tex 
    \end{commandshell}
	then the before-and-after results
	shown in \cref{lst:matrixbefore,lst:matrixafter} are achievable by default.

	\begin{minipage}{.45\textwidth}
		\cmhlistingsfromfile[columns=fixed]{demonstrations/matrix1.tex}{\texttt{matrix1.tex}}{lst:matrixbefore}
	\end{minipage}%
	\hfill
	\begin{minipage}{.45\textwidth}
		\cmhlistingsfromfile[columns=fixed]{demonstrations/matrix1-default.tex}{\texttt{matrix1.tex} default output}{lst:matrixafter}
	\end{minipage}%

	If you have blocks of code that you wish to align at the \&  character that
	are \emph{not} wrapped in, for example, \lstinline!\begin{tabular}! \ldots \lstinline!\end{tabular}!, then you can use the mark up
	illustrated in \cref{lst:alignmentmarkup}; the default output is shown in \cref{lst:alignmentmarkup-default}. Note that the \lstinline!%*! must be next to
	each other, but that there can be any number of spaces (possibly none) between the
	\lstinline!*! and \lstinline!\begin{tabular}!; note also that you may use any
	environment name that you have specified in \texttt{lookForAlignDelims}.

	\begin{minipage}{.45\textwidth}
		\cmhlistingsfromfile[columns=fixed]{demonstrations/align-block.tex}{\texttt{align-block.tex}}{lst:alignmentmarkup}
	\end{minipage}%
	\hfill
	\begin{minipage}{.45\textwidth}
		\cmhlistingsfromfile[columns=fixed]{demonstrations/align-block-default.tex}{\texttt{align-block.tex} default output}{lst:alignmentmarkup-default}
	\end{minipage}%

	With reference to \vref{tab:code-blocks} and the, yet undiscussed, fields of \texttt{noAdditionalIndent} and \texttt{indentRules}
	(see \vref{sec:noadd-indent-rules}), these comment-marked blocks are considered \texttt{environments}.

\yamltitle{indentAfterItems}*{fields}
	\begin{wrapfigure}[5]{r}[0pt]{7cm}
		\cmhlistingsfromfile[style=indentAfterItems]{../defaultSettings.yaml}[width=.8\linewidth,before=\centering,yaml-TCB]{\texttt{indentAfterItems}}{lst:indentafteritems}
	\end{wrapfigure}
	The environment names specified in \texttt{indentAfterItems}  tell
	\texttt{latexindent.pl} to look for \lstinline!\item! commands; if these switches are set to \texttt{1}
	then indentation will be performed so as indent the code after each \texttt{item}.
	A demonstration is given in \cref{lst:itemsbefore,lst:itemsafter}

	\begin{minipage}{.45\textwidth}
		\cmhlistingsfromfile{demonstrations/items1.tex}{\texttt{items1.tex}}{lst:itemsbefore}
	\end{minipage}%
	\hfill
	\begin{minipage}{.45\textwidth}
		\cmhlistingsfromfile{demonstrations/items1-default.tex}{\texttt{items1.tex} default output}{lst:itemsafter}
	\end{minipage}

\yamltitle{itemNames}*{fields}
	\begin{wrapfigure}[5]{r}[0pt]{5cm}
		\cmhlistingsfromfile[style=itemNames]{../defaultSettings.yaml}[width=.8\linewidth,before=\centering,yaml-TCB]{\texttt{itemNames}}{lst:itemNames}
	\end{wrapfigure}
	If you have your own \texttt{item} commands (perhaps you
	prefer to use \texttt{myitem}, for example)
	then you can put populate them in \texttt{itemNames}.
	For example, users of the \texttt{exam} document class might like to add
	\texttt{parts} to \texttt{indentAfterItems} and \texttt{part} to \texttt{itemNames}
	to their user settings (see \vref{sec:indentconfig} for details of how to configure user settings,
	and \vref{lst:mysettings} \\ in particular \label{page:examsettings}.)

\yamltitle{specialBeginEnd}*{fields}\label{yaml:specialBeginEnd}
	The fields specified in \texttt{specialBeginEnd} are, in their default state, focused on math mode begin and end statements, but
	there is no requirement for this to be the case; \cref{lst:specialBeginEnd} shows the
	default settings of \texttt{specialBeginEnd}.

	\cmhlistingsfromfile[style=specialBeginEnd]{../defaultSettings.yaml}[width=.8\linewidth,before=\centering,yaml-TCB]{\texttt{specialBeginEnd}}{lst:specialBeginEnd}

	The field \texttt{displayMath} represents \lstinline!\[...\]!, \texttt{inlineMath} represents
	\lstinline!$...$! and \texttt{displayMathTex} represents \lstinline!$$...$$!. You can, of course,
	rename these in your own YAML files (see \vref{sec:localsettings}); indeed, you
	might like to set up your own special begin and end statements.

	A demonstration of the before-and-after results are shown in \cref{lst:specialbefore,lst:specialafter}.

	\begin{minipage}{.45\textwidth}
		\cmhlistingsfromfile{demonstrations/special1.tex}{\texttt{special1.tex} before}{lst:specialbefore}
	\end{minipage}%
	\hfill
	\begin{minipage}{.45\textwidth}
		\cmhlistingsfromfile{demonstrations/special1-default.tex}{\texttt{special1.tex} default output}{lst:specialafter}
	\end{minipage}

	For each field, \texttt{lookForThis} is set to \texttt{1} by default, which means that \texttt{latexindent.pl}
	will look for this pattern; you can tell \texttt{latexindent.pl} not to look for the pattern, by setting
	\texttt{lookForThis} to \texttt{0}.

\yamltitle{indentAfterHeadings}*{fields}
	\begin{wrapfigure}[17]{r}[0pt]{8cm}
		\cmhlistingsfromfile[style=indentAfterHeadings]{../defaultSettings.yaml}[width=.8\linewidth,before=\centering,yaml-TCB]{\texttt{indentAfterHeadings}}{lst:indentAfterHeadings}
	\end{wrapfigure}
	This field enables the user to specify
	indentation rules that take effect after heading commands such as \lstinline!\part!, \lstinline!\chapter!,
	\lstinline!\section!, \lstinline!\subsection*!, or indeed any user-specified command written in this field.\footnote{There is a slight
		difference in interface for this field when comparing Version 2.2 to Version 3.0; see \vref{app:differences} for details.}

	The default settings do \emph{not} place indentation after a heading, but you
	can easily switch them on by changing \\ \texttt{indentAfterThisHeading: 0} to \\ \texttt{indentAfterThisHeading: 1}.
	The \texttt{level} field tells \texttt{latexindent.pl} the hierarchy of the heading
	structure in your document. You might, for example, like to have both \texttt{section}
	and \texttt{subsection} set with \texttt{level: 3} because you do not want the indentation to go too deep.

	You can add any of your own custom heading commands to this field, specifying the \texttt{level}
	as appropriate.  You can also specify your own indentation in \texttt{indentRules} (see \vref{sec:noadd-indent-rules});
	you will find the default \texttt{indentRules} contains \lstinline!chapter: " "! which
	tells \texttt{latexindent.pl} simply to use a space character after \texttt{\chapter} headings
	(once \texttt{indent} is set to \texttt{1} for \texttt{chapter}).

	For example, assuming that you have the code in \cref{lst:headings1yaml} saved into \texttt{headings1.yaml},
	and that you have the text from \cref{lst:headings1} saved into \texttt{headings1.tex}.

	\begin{minipage}{.45\textwidth}
		\cmhlistingsfromfile[style=yaml-LST]{demonstrations/headings1.yaml}[yaml-TCB]{\texttt{headings1.yaml}}{lst:headings1yaml}
	\end{minipage}%
	\hfill
	\begin{minipage}{.45\textwidth}
		\cmhlistingsfromfile{demonstrations/headings1.tex}{\texttt{headings1.tex}}{lst:headings1}
	\end{minipage}

	If you run the command
	\begin{commandshell}
latexindent.pl headings1.tex -l=headings1.yaml
\end{commandshell}
	then you should receive the output given in \cref{lst:headings1-mod1}.

	\begin{minipage}{.45\textwidth}
		\cmhlistingsfromfile{demonstrations/headings1-mod1.tex}{\texttt{headings1.tex} using \cref{lst:headings1yaml}}{lst:headings1-mod1}
	\end{minipage}%
	\hfill
	\begin{minipage}{.45\textwidth}
		\cmhlistingsfromfile{demonstrations/headings1-mod2.tex}{\texttt{headings1.tex} second modification}{lst:headings1-mod2}
	\end{minipage}

	Now say that you modify the \texttt{YAML} from \cref{lst:headings1yaml} so that the \texttt{paragraph} \texttt{level} is \texttt{1}; after
	running
	\begin{commandshell}
latexindent.pl headings1.tex -l=headings1.yaml
\end{commandshell}
	you should receive the code given in \cref{lst:headings1-mod2}; notice that
	the \texttt{paragraph} and \texttt{subsection} are at the same indentation level.

\subsection{The code blocks known \texttt{latexindent.pl}}\label{subsubsec:code-blocks}
	As of Version 3.0, \texttt{latexindent.pl} processes documents using code blocks; each
	of these are shown in \cref{tab:code-blocks}.

	\begin{longtable}{m{.3\linewidth}@{\hspace{.25cm}}m{.4\linewidth}@{}m{.2\linewidth}}
		\caption{Code blocks known to \texttt{latexindent.pl}}\label{tab:code-blocks}\\
		\toprule
		Code block                    & characters allowed in name                                                                                     & example                                                                                                                                                                                                                                     \\
		\midrule
		environments                  & \lstinline!a-zA-Z@\*0-9_\\!                                                                                        &
		\begin{lstlisting}[,nolol=true,]
\begin{myenv}
body of myenv
\end{myenv}
  \end{lstlisting}
		\\\cmidrule{2-3}
		optionalArguments             & \emph{inherits} name from parent (e.g environment name)                                                        &
		\begin{lstlisting}[,nolol=true,]
[
opt arg text
]
  \end{lstlisting}
		\\\cmidrule{2-3}
		mandatoryArguments            & \emph{inherits} name from parent (e.g environment name)                                                        &
		\begin{lstlisting}[,nolol=true,]
{
mand arg text
}
  \end{lstlisting}
		\\\cmidrule{2-3}
		commands                      & \lstinline!+a-zA-Z@\*0-9_\:!                                                                                        & \lstinline!\mycommand!$\langle$\itshape{arguments}$\rangle$                                                                                                                                                                                \\\cmidrule{2-3}
		keyEqualsValuesBracesBrackets & \lstinline!a-zA-Z@\*0-9_\/.\h\{\}:\#-!                                                                                        & \lstinline!my key/.style=!$\langle$\itshape{arguments}$\rangle$                                                                                                                                                                                \\\cmidrule{2-3}
		namedGroupingBracesBrackets   & \lstinline!a-zA-Z@\*><!                                                                                        & \lstinline!in!$\langle$\itshape{arguments}$\rangle$                                                                                                                                                                                \\\cmidrule{2-3}
		UnNamedGroupingBracesBrackets & \centering\emph{No name!}                                                                                      & \lstinline!{! or \lstinline![! or \lstinline!,! or \lstinline!&! or \lstinline!)! or \lstinline!(! or \lstinline!$! followed by $\langle$\itshape{arguments}$\rangle$ \\\cmidrule{2-3}
		ifElseFi                      & \lstinline!@a-zA-Z! but must begin with either \newline \lstinline!\if! of \lstinline!\@if! &
		\begin{lstlisting}[,nolol=true,]
\ifnum...
...
\else
...
\fi
  \end{lstlisting}\\\cmidrule{2-3}
		items                         & User specified, see \vref{lst:indentafteritems,lst:itemNames}                                                  &
		\begin{lstlisting}[,nolol=true,]
\begin{enumerate}
  \item ...
\end{enumerate}
  \end{lstlisting}\\\cmidrule{2-3}
		specialBeginEnd               & User specified, see \vref{lst:specialBeginEnd}                                                                 &
		\begin{lstlisting}[,nolol=true,]
\[
  ...
\]
  \end{lstlisting}\\\cmidrule{2-3}
		afterHeading                  & User specified, see \vref{lst:indentAfterHeadings}                                                             &
		\begin{lstlisting}[,morekeywords={chapter},nolol=true,]
\chapter{title}
  ...
\section{title}
  \end{lstlisting}\\\cmidrule{2-3}
		filecontents                  & User specified, see \vref{lst:fileContentsEnvironments}                                                        &
		\begin{lstlisting}[,nolol=true,]
\begin{filecontents}
...
\end{filecontents}
  \end{lstlisting}\\
		\bottomrule
	\end{longtable}

	We will refer to these code blocks in what follows.

% arara: pdflatex: {shell: yes, files: [latexindent]}
% arara: pdflatex: {shell: yes, files: [latexindent]}
\subsection{\texttt{noAdditionalIndent} and \texttt{indentRules}}\label{sec:noadd-indent-rules}
	\texttt{latexindent.pl} operates on files by looking for code blocks, as detailed in \vref{subsubsec:code-blocks};
	for each type of code block  in \vref{tab:code-blocks} (which we will call a \emph{$\langle$thing$\rangle$} in what follows)
	it searches YAML fields for information in the following order:
	\begin{enumerate}
		\item \texttt{noAdditionalIndent} for the \emph{name} of the current \emph{$\langle$thing$\rangle$};
		\item \texttt{indentRules} for the \emph{name} of the current \emph{$\langle$thing$\rangle$};
		\item \texttt{noAdditionalIndentGlobal} for the \emph{type} of the current \emph{$\langle$thing$\rangle$};
		\item \texttt{indentRulesGlobal} for the \emph{type} of the current \emph{$\langle$thing$\rangle$}.
	\end{enumerate}

	Using the above list, the first piece of information to be found will be used; failing that,
	the value of \texttt{defaultIndent} is used.
	If information is found in multiple fields, the first one according to the list above will be used; for example,
	if information is present in both \texttt{indentRules} and in \texttt{noAdditionalIndentGlobal}, then the information from \texttt{indentRules}
	takes priority.

	We now present details for the different type of code blocks known to \texttt{latexindent.pl}, as detailed in \vref{tab:code-blocks}; for
	reference, there follows a list of the code blocks covered.

	\startcontents[noAdditionalIndent]
	\printcontents[noAdditionalIndent]{}{0}{}

% arara: pdflatex: {shell: yes, files: [latexindent]}
\subsubsection{Environments and their arguments}\label{subsubsec:env-and-their-args}
	There are a few different YAML switches governing the indentation of environments; let's start
	with the code shown in \cref{lst:myenvtex}.

	\cmhlistingsfromfile{demonstrations/myenvironment-simple.tex}{\texttt{myenv.tex}}{lst:myenvtex}

\yamltitle{noAdditionalIndent}*{fields}
	If we do not wish \texttt{myenv} to receive any additional indentation, we have a few choices available to us,
	as demonstrated in \cref{lst:myenv-noAdd1,lst:myenv-noAdd2}.

	\begin{minipage}{.45\textwidth}
		\cmhlistingsfromfile[style=yaml-LST]{demonstrations/myenv-noAdd1.yaml}[width=.8\linewidth,before=\centering,yaml-TCB]{\texttt{myenv-noAdd1.yaml}}{lst:myenv-noAdd1}
	\end{minipage}
	\hfill
	\begin{minipage}{.45\textwidth}
		\cmhlistingsfromfile[style=yaml-LST]{demonstrations/myenv-noAdd2.yaml}[width=.8\linewidth,before=\centering,yaml-TCB]{\texttt{myenv-noAdd2.yaml}}{lst:myenv-noAdd2}
	\end{minipage}

	On applying either of the following commands,
	\begin{commandshell}
latexindent.pl myenv.tex -l myenv-noAdd1.yaml  
latexindent.pl myenv.tex -l myenv-noAdd2.yaml  
\end{commandshell}
	we obtain the output given in \cref{lst:myenv-output}; note in particular that the environment \texttt{myenv}
	has not received any \emph{additional} indentation, but that the \texttt{outer} environment \emph{has} still
	received indentation.

	\cmhlistingsfromfile{demonstrations/myenvironment-simple-noAdd-body1.tex}{\texttt{myenv.tex output (using either \cref{lst:myenv-noAdd1} or \cref{lst:myenv-noAdd2})}}{lst:myenv-output}

	Upon changing the YAML files to those shown in \cref{lst:myenv-noAdd3,lst:myenv-noAdd4}, and running either
	\begin{commandshell}
latexindent.pl myenv.tex -l myenv-noAdd3.yaml  
latexindent.pl myenv.tex -l myenv-noAdd4.yaml  
\end{commandshell}
	we obtain the output given in \cref{lst:myenv-output-4}.

	\begin{minipage}{.45\textwidth}
		\cmhlistingsfromfile[style=yaml-LST]{demonstrations/myenv-noAdd3.yaml}[width=.8\linewidth,before=\centering,yaml-TCB]{\texttt{myenv-noAdd3.yaml}}{lst:myenv-noAdd3}
	\end{minipage}
	\hfill
	\begin{minipage}{.45\textwidth}
		\cmhlistingsfromfile[style=yaml-LST]{demonstrations/myenv-noAdd4.yaml}[width=.8\linewidth,before=\centering,yaml-TCB]{\texttt{myenv-noAdd4.yaml}}{lst:myenv-noAdd4}
	\end{minipage}

	\cmhlistingsfromfile{demonstrations/myenvironment-simple-noAdd-body4.tex}{\texttt{myenv.tex output} (using either \cref{lst:myenv-noAdd3} or \cref{lst:myenv-noAdd4})}{lst:myenv-output-4}

	Let's now allow \texttt{myenv} to have some optional and mandatory arguments, as in \cref{lst:myenv-args}.
	\cmhlistingsfromfile{demonstrations/myenvironment-args.tex}{\texttt{myenv-args.tex}}{lst:myenv-args}
	Upon running
	\begin{commandshell}
latexindent.pl -l=myenv-noAdd1.yaml myenv-args.tex  
\end{commandshell}
	we obtain the output shown in \cref{lst:myenv-args-noAdd1}; note that the optional argument, mandatory argument and body \emph{all}
	have received no additional indent. This is because, when \texttt{noAdditionalIndent} is specified in `scalar' form (as in \cref{lst:myenv-noAdd1}),
	then \emph{all} parts of the environment (body, optional and mandatory arguments) are assumed to want no additional indent.
	\cmhlistingsfromfile{demonstrations/myenvironment-args-noAdd-body1.tex}{\texttt{myenv-args.tex} using \cref{lst:myenv-noAdd1}}{lst:myenv-args-noAdd1}

	We may customise \texttt{noAdditionalIndent} for optional and mandatory arguments of the \texttt{myenv} environment, as shown in, for example, \cref{lst:myenv-noAdd5,lst:myenv-noAdd6}.

	\begin{minipage}{.49\textwidth}
		\cmhlistingsfromfile[style=yaml-LST]{demonstrations/myenv-noAdd5.yaml}[width=.8\linewidth,before=\centering,yaml-TCB]{\texttt{myenv-noAdd5.yaml}}{lst:myenv-noAdd5}
	\end{minipage}
	\hfill
	\begin{minipage}{.49\textwidth}
		\cmhlistingsfromfile[style=yaml-LST]{demonstrations/myenv-noAdd6.yaml}[width=.8\linewidth,before=\centering,yaml-TCB]{\texttt{myenv-noAdd6.yaml}}{lst:myenv-noAdd6}
	\end{minipage}

	Upon running
	\begin{commandshell}
latexindent.pl myenv.tex -l myenv-noAdd5.yaml  
latexindent.pl myenv.tex -l myenv-noAdd6.yaml  
\end{commandshell}
	we obtain the respective outputs given in \cref{lst:myenv-args-noAdd5,lst:myenv-args-noAdd6}. Note that in \cref{lst:myenv-args-noAdd5}
	the text for the \emph{optional} argument has not received any additional indentation, and that in \cref{lst:myenv-args-noAdd6} the
	\emph{mandatory} argument has not received any additional indentation; in both cases, the \emph{body} has not received any additional indentation.

	\begin{minipage}{.45\textwidth}
		\cmhlistingsfromfile{demonstrations/myenvironment-args-noAdd5.tex}{\texttt{myenv-args.tex} using \cref{lst:myenv-noAdd5}}{lst:myenv-args-noAdd5}
	\end{minipage}
	\hfill
	\begin{minipage}{.45\textwidth}
		\cmhlistingsfromfile{demonstrations/myenvironment-args-noAdd6.tex}{\texttt{myenv-args.tex} using \cref{lst:myenv-noAdd6}}{lst:myenv-args-noAdd6}
	\end{minipage}

\yamltitle{indentRules}*{fields}
	We may also specify indentation rules for environment code blocks using the \texttt{indentRules} field; see, for example,
	\cref{lst:myenv-rules1,lst:myenv-rules2}.

	\begin{minipage}{.45\textwidth}
		\cmhlistingsfromfile[style=yaml-LST]{demonstrations/myenv-rules1.yaml}[width=.8\linewidth,before=\centering,yaml-TCB]{\texttt{myenv-rules1.yaml}}{lst:myenv-rules1}
	\end{minipage}
	\hfill
	\begin{minipage}{.45\textwidth}
		\cmhlistingsfromfile[style=yaml-LST]{demonstrations/myenv-rules2.yaml}[width=.8\linewidth,before=\centering,yaml-TCB]{\texttt{myenv-rules2.yaml}}{lst:myenv-rules2}
	\end{minipage}

	On applying either of the following commands,
	\begin{commandshell}
latexindent.pl myenv.tex -l myenv-rules1.yaml  
latexindent.pl myenv.tex -l myenv-rules2.yaml  
\end{commandshell}
	we obtain the output given in \cref{lst:myenv-rules-output}; note in particular that the environment \texttt{myenv}
	has received one tab (from the \texttt{outer} environment) plus three spaces from \cref{lst:myenv-rules1} or \ref{lst:myenv-rules2}.

	\cmhlistingsfromfile{demonstrations/myenv-rules1.tex}{\texttt{myenv.tex output (using either \cref{lst:myenv-rules1} or \cref{lst:myenv-rules2})}}{lst:myenv-rules-output}

	If you specify a field in \texttt{indentRules} using anything other than horizontal space, it will be ignored.

	Returning to the example in \cref{lst:myenv-args} that contains optional and mandatory arguments. Upon using \cref{lst:myenv-rules1} as in
	\begin{commandshell}
latexindent.pl myenv-args.tex -l=myenv-rules1.yaml  
\end{commandshell}
	we obtain the output in \cref{lst:myenv-args-rules1}; note that the body, optional argument and mandatory argument have \emph{all}
	received the same customised indentation.
	\cmhlistingsfromfile{demonstrations/myenvironment-args-rules1.tex}{\texttt{myenv-args.tex} using \cref{lst:myenv-rules1}}{lst:myenv-args-rules1}

	You can specify different indentation rules for the different features using, for example, \cref{lst:myenv-rules3,lst:myenv-rules4}

	\begin{minipage}{.49\textwidth}
		\cmhlistingsfromfile[style=yaml-LST]{demonstrations/myenv-rules3.yaml}[width=.9\linewidth,before=\centering,yaml-TCB]{\texttt{myenv-rules3.yaml}}{lst:myenv-rules3}
	\end{minipage}
	\hfill
	\begin{minipage}{.49\textwidth}
		\cmhlistingsfromfile[style=yaml-LST]{demonstrations/myenv-rules4.yaml}[width=.9\linewidth,before=\centering,yaml-TCB]{\texttt{myenv-rules4.yaml}}{lst:myenv-rules4}
	\end{minipage}

	After running
	\begin{commandshell}
latexindent.pl myenv-args.tex -l myenv-rules3.yaml  
latexindent.pl myenv-args.tex -l myenv-rules4.yaml  
\end{commandshell}
	then we obtain the respective outputs given in \cref{lst:myenv-args-rules3,lst:myenv-args-rules4}.

	\begin{minipage}{.45\textwidth}
		\cmhlistingsfromfile{demonstrations/myenvironment-args-rules3.tex}{\texttt{myenv-args.tex} using \cref{lst:myenv-rules3}}{lst:myenv-args-rules3}
	\end{minipage}
	\hfill
	\begin{minipage}{.45\textwidth}
		\cmhlistingsfromfile{demonstrations/myenvironment-args-rules4.tex}{\texttt{myenv-args.tex} using \cref{lst:myenv-rules4}}{lst:myenv-args-rules4}
	\end{minipage}

	Note that in \cref{lst:myenv-args-rules3}, the optional argument has only received a single space of indentation, while the mandatory argument
	has received the default (tab) indentation; the environment body has received three spaces of indentation.

	In \cref{lst:myenv-args-rules4}, the optional argument has received the default (tab) indentation, the mandatory argument has received two tabs
	of indentation, and the body has received three spaces of indentation.

\yamltitle{noAdditionalIndentGlobal}*{fields}
	\begin{wrapfigure}[6]{r}[0pt]{7cm}
		\cmhlistingsfromfile[style=noAdditionalIndentGlobalEnv]{../defaultSettings.yaml}[width=.8\linewidth,before=\centering,yaml-TCB]{\texttt{noAdditionalIndentGlobal}}{lst:noAdditionalIndentGlobal:environments}
	\end{wrapfigure}
	Assuming that your environment name is not found within neither \texttt{noAdditionalIndent} nor \texttt{indentRules}, the next
	place that \texttt{latexindent.pl} will look is \texttt{noAdditionalIndentGlobal}, and in particular \emph{for the environments} key
	(see \cref{lst:noAdditionalIndentGlobal:environments}). Let's say that you change
	the value of \texttt{environments} to \texttt{1} in \cref{lst:noAdditionalIndentGlobal:environments}, and that you run

	\begin{widepage}
		\begin{commandshell}
latexindent.pl myenv-args.tex -l env-noAdditionalGlobal.yaml
latexindent.pl myenv-args.tex -l myenv-rules1.yaml,env-noAdditionalGlobal.yaml
\end{commandshell}
	\end{widepage}

	The respective output from these two commands are in \cref{lst:myenv-args-no-add-global1,lst:myenv-args-no-add-global2}; in \cref{lst:myenv-args-no-add-global1} notice that \emph{both}
	environments receive no additional indentation but that the arguments of \texttt{myenv} still \emph{do} receive indentation. In \cref{lst:myenv-args-no-add-global2}
	notice that the \emph{outer} environment does not receive additional indentation, but because of the settings from \texttt{myenv-rules1.yaml} (in \vref{lst:myenv-rules1}), the \texttt{myenv}
	environment still \emph{does} receive indentation.

	\begin{minipage}{.45\textwidth}
		\cmhlistingsfromfile{demonstrations/myenvironment-args-rules1-noAddGlobal1.tex}{\texttt{myenv-args.tex} using \cref{lst:noAdditionalIndentGlobal:environments}}{lst:myenv-args-no-add-global1}
	\end{minipage}
	\hfill
	\begin{minipage}{.45\textwidth}
		\cmhlistingsfromfile{demonstrations/myenvironment-args-rules1-noAddGlobal2.tex}{\texttt{myenv-args.tex} using \cref{lst:noAdditionalIndentGlobal:environments,lst:myenv-rules1}}{lst:myenv-args-no-add-global2}
	\end{minipage}

	In fact, \texttt{noAdditionalIndentGlobal} also contains keys that control the indentation of optional and mandatory
	arguments; on referencing \cref{lst:opt-args-no-add-glob,lst:mand-args-no-add-glob}

	\begin{minipage}{.49\textwidth}
		\cmhlistingsfromfile[style=yaml-LST]{demonstrations/opt-args-no-add-glob.yaml}[width=.8\linewidth,before=\centering,yaml-TCB]{\texttt{opt-args-no-add-glob.yaml}}{lst:opt-args-no-add-glob}
	\end{minipage}
	\hfill
	\begin{minipage}{.49\textwidth}
		\cmhlistingsfromfile[style=yaml-LST]{demonstrations/mand-args-no-add-glob.yaml}[width=.8\linewidth,before=\centering,yaml-TCB]{\texttt{mand-args-no-add-glob.yaml}}{lst:mand-args-no-add-glob}
	\end{minipage}

	we may run the commands
	\begin{commandshell}
latexindent.pl  myenv-args.tex -local opt-args-no-add-glob.yaml
latexindent.pl  myenv-args.tex -local mand-args-no-add-glob.yaml
\end{commandshell}
	which produces the respective outputs given in \cref{lst:myenv-args-no-add-opt,lst:myenv-args-no-add-mand}. Notice that in \cref{lst:myenv-args-no-add-opt}
	the \emph{optional} argument has not received any additional indentation, and in \cref{lst:myenv-args-no-add-mand} the \emph{mandatory} argument
	has not received any additional indentation.

	\begin{minipage}{.45\textwidth}
		\cmhlistingsfromfile{demonstrations/myenvironment-args-rules1-noAddGlobal3.tex}{\texttt{myenv-args.tex} using \cref{lst:opt-args-no-add-glob}}{lst:myenv-args-no-add-opt}
	\end{minipage}
	\hfill
	\begin{minipage}{.45\textwidth}
		\cmhlistingsfromfile{demonstrations/myenvironment-args-rules1-noAddGlobal4.tex}{\texttt{myenv-args.tex} using \cref{lst:mand-args-no-add-glob}}{lst:myenv-args-no-add-mand}
	\end{minipage}

\yamltitle{indentRulesGlobal}*{fields}
	\begin{wrapfigure}[4]{r}[0pt]{7cm}
		\cmhlistingsfromfile[style=indentRulesGlobalEnv]{../defaultSettings.yaml}[width=.8\linewidth,before=\centering,yaml-TCB]{\texttt{indentRulesGlobal}}{lst:indentRulesGlobal:environments}
	\end{wrapfigure}
	The final check that \texttt{latexindent.pl} will make is to look for \texttt{indentRulesGlobal} as detailed in \cref{lst:indentRulesGlobal:environments}; if you change the \texttt{environments}
	field to anything involving horizontal space, say \lstinline!" "!, and then run the following commands

	\begin{commandshell}
latexindent.pl  myenv-args.tex -l env-indentRules.yaml
latexindent.pl  myenv-args.tex -l myenv-rules1.yaml,env-indentRules.yaml
\end{commandshell}
	then the respective output is shown in \cref{lst:myenv-args-indent-rules-global1,lst:myenv-args-indent-rules-global2}. Note that
	in \cref{lst:myenv-args-indent-rules-global1}, both the environment blocks have received a single-space indentation, whereas in
	\cref{lst:myenv-args-indent-rules-global2} the \texttt{outer} environment has received single-space indentation (specified by \texttt{indentRulesGlobal}),
	but \texttt{myenv} has received \lstinline!"   "!, as specified by the particular \texttt{indentRules} for \texttt{myenv} \vref{lst:myenv-rules1}.

	\begin{minipage}{.45\textwidth}
		\cmhlistingsfromfile{demonstrations/myenvironment-args-global-rules1.tex}{\texttt{myenv-args.tex} using \cref{lst:indentRulesGlobal:environments}}{lst:myenv-args-indent-rules-global1}
	\end{minipage}
	\hfill
	\begin{minipage}{.45\textwidth}
		\cmhlistingsfromfile{demonstrations/myenvironment-args-global-rules2.tex}{\texttt{myenv-args.tex} using \cref{lst:myenv-rules1,lst:indentRulesGlobal:environments}}{lst:myenv-args-indent-rules-global2}
	\end{minipage}

	You can specify \texttt{indentRulesGlobal} for both optional and mandatory arguments, as detailed in \cref{lst:opt-args-indent-rules-glob,lst:mand-args-indent-rules-glob}

	\begin{minipage}{.49\textwidth}
		\cmhlistingsfromfile[style=yaml-LST]{demonstrations/opt-args-indent-rules-glob.yaml}[width=.9\linewidth,before=\centering,yaml-TCB]{\texttt{opt-args-indent-rules-glob.yaml}}{lst:opt-args-indent-rules-glob}
	\end{minipage}
	\hfill
	\begin{minipage}{.49\textwidth}
		\cmhlistingsfromfile[style=yaml-LST]{demonstrations/mand-args-indent-rules-glob.yaml}[width=.9\linewidth,before=\centering,yaml-TCB]{\texttt{mand-args-indent-rules-glob.yaml}}{lst:mand-args-indent-rules-glob}
	\end{minipage}

	Upon running the following commands
	\begin{commandshell}
latexindent.pl  myenv-args.tex -local opt-args-indent-rules-glob.yaml
latexindent.pl  myenv-args.tex -local mand-args-indent-rules-glob.yaml
\end{commandshell}
	we obtain the respective outputs in \cref{lst:myenv-args-indent-rules-global3,lst:myenv-args-indent-rules-global4}. Note that the \emph{optional}
	argument in \cref{lst:myenv-args-indent-rules-global3} has received two tabs worth of indentation, while the \emph{mandatory} argument has
	done so in \cref{lst:myenv-args-indent-rules-global4}.

	\begin{minipage}{.45\textwidth}
		\cmhlistingsfromfile{demonstrations/myenvironment-args-global-rules3.tex}{\texttt{myenv-args.tex} using \cref{lst:opt-args-indent-rules-glob}}{lst:myenv-args-indent-rules-global3}
	\end{minipage}
	\hfill
	\begin{minipage}{.45\textwidth}
		\cmhlistingsfromfile{demonstrations/myenvironment-args-global-rules4.tex}{\texttt{myenv-args.tex} using \cref{lst:mand-args-indent-rules-glob}}{lst:myenv-args-indent-rules-global4}
	\end{minipage}

% arara: pdflatex: {shell: yes, files: [latexindent]}
\subsubsection{Environments with items}
	With reference to \vref{lst:indentafteritems,lst:itemNames}, some commands
	may contain \texttt{item} commands; for the purposes of this discussion,
	we will use the code from \vref{lst:itemsbefore}.

	Assuming that you've populated \texttt{itemNames} with the name of your
	\texttt{item}, you can put the item name into \texttt{noAdditionalIndent}
	as in \cref{lst:item-noAdd1}, although a more efficient approach may be
	to change the relevant field in \texttt{itemNames} to \texttt{0}. Similarly,
	you can customise the indentation that your \texttt{item} receives using
	\texttt{indentRules}, as in \cref{lst:item-rules1}

	\begin{minipage}{.45\textwidth}
		\cmhlistingsfromfile[style=yaml-LST]{demonstrations/item-noAdd1.yaml}[yaml-TCB]{\texttt{item-noAdd1.yaml}}{lst:item-noAdd1}
	\end{minipage}%
	\hfill
	\begin{minipage}{.45\textwidth}
		\cmhlistingsfromfile[style=yaml-LST]{demonstrations/item-rules1.yaml}[yaml-TCB]{\texttt{item-rules1.yaml}}{lst:item-rules1}
	\end{minipage}

	Upon running the following commands
	\begin{commandshell}
latexindent.pl items1.tex -local item-noAdd1.yaml  
latexindent.pl items1.tex -local item-rules1.yaml  
\end{commandshell}
	the respective outputs are given in \cref{lst:items1-noAdd1,lst:items1-rules1}; note that in \cref{lst:items1-noAdd1}
	that the text after each \texttt{item} has not received any additional indentation, and in \cref{lst:items1-rules1},
	the text after each \texttt{item} has received a single space of indentation, specified by \cref{lst:item-rules1}.

	\begin{minipage}{.45\textwidth}
		\cmhlistingsfromfile{demonstrations/items1-noAdd1.tex}{\texttt{items1.tex} using \cref{lst:item-noAdd1}}{lst:items1-noAdd1}
	\end{minipage}
	\hfill
	\begin{minipage}{.45\textwidth}
		\cmhlistingsfromfile{demonstrations/items1-rules1.tex}{\texttt{items1.tex} using \cref{lst:item-rules1}}{lst:items1-rules1}
	\end{minipage}

	Alternatively, you might like to populate \texttt{noAdditionalIndentGlobal} or \texttt{indentRulesGlobal} using the \texttt{items}
	key, as demonstrated in \cref{lst:items-noAdditionalGlobal,lst:items-indentRulesGlobal}. Note that there is a need to
	`reset/remove' the \texttt{item} field from \texttt{indentRules} in both cases (see the hierarchy description given on \cpageref{sec:noadd-indent-rules})
	as the \texttt{item} command is a member of \texttt{indentRules} by default.

	\begin{minipage}{.45\textwidth}
		\cmhlistingsfromfile[style=yaml-LST]{demonstrations/items-noAdditionalGlobal.yaml}[yaml-TCB]{\texttt{items-noAdditionalGlobal.yaml}}{lst:items-noAdditionalGlobal}
	\end{minipage}%
	\hfill
	\begin{minipage}{.45\textwidth}
		\cmhlistingsfromfile[style=yaml-LST]{demonstrations/items-indentRulesGlobal.yaml}[yaml-TCB]{\texttt{items-indentRulesGlobal.yaml}}{lst:items-indentRulesGlobal}
	\end{minipage}

	Upon running the following commands,
	\begin{commandshell}
latexindent.pl items1.tex -local items-noAdditionalGlobal.yaml
latexindent.pl items1.tex -local items-indentRulesGlobal.yaml
\end{commandshell}
	the respective outputs from \cref{lst:items1-noAdd1,lst:items1-rules1} are obtained; note, however, that
	\emph{all} such \texttt{item} commands without their own individual \texttt{noAdditionalIndent} or \texttt{indentRules}
	settings would behave as in these listings.

% arara: pdflatex: {shell: yes, files: [latexindent]}
\subsubsection{Commands with arguments}\label{subsubsec:commands-arguments}
	Let's begin with the simple example in \cref{lst:mycommand}; when \texttt{latexindent.pl} operates
	on this file, the default output is shown in \cref{lst:mycommand-default}. \footnote{The command code blocks
		have quite a few subtleties, described in \vref{subsec:commands-string-between}.}

	\begin{minipage}{.45\textwidth}
		\cmhlistingsfromfile{demonstrations/mycommand.tex}{\texttt{mycommand.tex}}{lst:mycommand}
	\end{minipage}%
	\hfill
	\begin{minipage}{.45\textwidth}
		\cmhlistingsfromfile{demonstrations/mycommand-default.tex}{\texttt{mycommand.tex} default output}{lst:mycommand-default}
	\end{minipage}

	As in the environment-based case (see \vref{lst:myenv-noAdd1,lst:myenv-noAdd2}) we may specify \texttt{noAdditionalIndent}
	either in `scalar' form, or in `field' form, as shown in \cref{lst:mycommand-noAdd1,lst:mycommand-noAdd2}

	\begin{minipage}{.45\textwidth}
		\cmhlistingsfromfile[style=yaml-LST]{demonstrations/mycommand-noAdd1.yaml}[width=.8\linewidth,before=\centering,yaml-TCB]{\texttt{mycommand-noAdd1.yaml}}{lst:mycommand-noAdd1}
	\end{minipage}
	\hfill
	\begin{minipage}{.45\textwidth}
		\cmhlistingsfromfile[style=yaml-LST]{demonstrations/mycommand-noAdd2.yaml}[width=.8\linewidth,before=\centering,yaml-TCB]{\texttt{mycommand-noAdd2.yaml}}{lst:mycommand-noAdd2}
	\end{minipage}

	After running the following commands,
	\begin{commandshell}
latexindent.pl mycommand.tex -l mycommand-noAdd1.yaml  
latexindent.pl mycommand.tex -l mycommand-noAdd2.yaml  
\end{commandshell}
	we receive the respective output given in \cref{lst:mycommand-output-noAdd1,lst:mycommand-output-noAdd2}

	\begin{minipage}{.45\textwidth}
		\cmhlistingsfromfile{demonstrations/mycommand-noAdd1.tex}{\texttt{mycommand.tex} using \cref{lst:mycommand-noAdd1}}{lst:mycommand-output-noAdd1}
	\end{minipage}
	\hfill
	\begin{minipage}{.45\textwidth}
		\cmhlistingsfromfile{demonstrations/mycommand-noAdd2.tex}{\texttt{mycommand.tex} using \cref{lst:mycommand-noAdd2}}{lst:mycommand-output-noAdd2}
	\end{minipage}

	Note that in \cref{lst:mycommand-output-noAdd1} that the `body', optional argument \emph{and} mandatory argument have \emph{all} received
	no additional indentation, while in \cref{lst:mycommand-output-noAdd2}, only the `body' has not received any additional indentation. We define
	the `body' of a command as any lines following the command name that include its optional or mandatory arguments.

	We may further customise \texttt{noAdditionalIndent} for \texttt{mycommand} as we did in \vref{lst:myenv-noAdd5,lst:myenv-noAdd6}; explicit examples
	are given in \cref{lst:mycommand-noAdd3,lst:mycommand-noAdd4}.

	\begin{minipage}{.45\textwidth}
		\cmhlistingsfromfile[style=yaml-LST]{demonstrations/mycommand-noAdd3.yaml}[width=.9\linewidth,before=\centering,yaml-TCB]{\texttt{mycommand-noAdd3.yaml}}{lst:mycommand-noAdd3}
	\end{minipage}
	\hfill
	\begin{minipage}{.45\textwidth}
		\cmhlistingsfromfile[style=yaml-LST]{demonstrations/mycommand-noAdd4.yaml}[width=.9\linewidth,before=\centering,yaml-TCB]{\texttt{mycommand-noAdd4.yaml}}{lst:mycommand-noAdd4}
	\end{minipage}

	After running the following commands,
	\begin{commandshell}
latexindent.pl mycommand.tex -l mycommand-noAdd3.yaml  
latexindent.pl mycommand.tex -l mycommand-noAdd4.yaml  
\end{commandshell}
	we receive the respective output given in \cref{lst:mycommand-output-noAdd3,lst:mycommand-output-noAdd4}.

	\begin{minipage}{.45\textwidth}
		\cmhlistingsfromfile{demonstrations/mycommand-noAdd3.tex}{\texttt{mycommand.tex} using \cref{lst:mycommand-noAdd3}}{lst:mycommand-output-noAdd3}
	\end{minipage}
	\hfill
	\begin{minipage}{.45\textwidth}
		\cmhlistingsfromfile{demonstrations/mycommand-noAdd4.tex}{\texttt{mycommand.tex} using \cref{lst:mycommand-noAdd4}}{lst:mycommand-output-noAdd4}
	\end{minipage}

	Attentive readers will note that the body of \texttt{mycommand} in both \cref{lst:mycommand-output-noAdd3,lst:mycommand-output-noAdd4}
	has received no additional indent, even though \texttt{body} is explicitly set to \texttt{0} in both \cref{lst:mycommand-noAdd3,lst:mycommand-noAdd4}.
	This is because, by default, \texttt{noAdditionalIndentGlobal} for \texttt{commands} is set to \texttt{1} by default; this can be easily
	fixed as in \cref{lst:mycommand-noAdd5,lst:mycommand-noAdd6}.\label{page:command:noAddGlobal}

	\begin{minipage}{.45\textwidth}
		\cmhlistingsfromfile[style=yaml-LST]{demonstrations/mycommand-noAdd5.yaml}[width=.9\linewidth,before=\centering,yaml-TCB]{\texttt{mycommand-noAdd5.yaml}}{lst:mycommand-noAdd5}
	\end{minipage}
	\hfill
	\begin{minipage}{.45\textwidth}
		\cmhlistingsfromfile[style=yaml-LST]{demonstrations/mycommand-noAdd6.yaml}[width=.9\linewidth,before=\centering,yaml-TCB]{\texttt{mycommand-noAdd6.yaml}}{lst:mycommand-noAdd6}
	\end{minipage}

	After running the following commands,
	\begin{commandshell}
latexindent.pl mycommand.tex -l mycommand-noAdd5.yaml  
latexindent.pl mycommand.tex -l mycommand-noAdd6.yaml  
\end{commandshell}
	we receive the respective output given in \cref{lst:mycommand-output-noAdd5,lst:mycommand-output-noAdd6}.

	\begin{minipage}{.45\textwidth}
		\cmhlistingsfromfile{demonstrations/mycommand-noAdd5.tex}{\texttt{mycommand.tex} using \cref{lst:mycommand-noAdd5}}{lst:mycommand-output-noAdd5}
	\end{minipage}
	\hfill
	\begin{minipage}{.45\textwidth}
		\cmhlistingsfromfile{demonstrations/mycommand-noAdd6.tex}{\texttt{mycommand.tex} using \cref{lst:mycommand-noAdd6}}{lst:mycommand-output-noAdd6}
	\end{minipage}

	Both \texttt{indentRules} and \texttt{indentRulesGlobal} can be adjusted as they were for \emph{environment} code blocks, as in
	\vref{lst:myenv-rules3,lst:myenv-rules4} and \vref{lst:indentRulesGlobal:environments,lst:opt-args-indent-rules-glob,lst:mand-args-indent-rules-glob}.

% arara: pdflatex: {shell: yes, files: [latexindent]}
\subsubsection{ifelsefi code blocks}
	Let's use the simple example shown in \cref{lst:ifelsefi1}; when
	\texttt{latexindent.pl} operates on this file, the output as in \cref{lst:ifelsefi1-default};
	note that the body of each of the \lstinline!\if! statements have been indented,
	and that the \lstinline!\else! statement has been accounted for correctly.

	\begin{minipage}{.45\textwidth}
		\cmhlistingsfromfile{demonstrations/ifelsefi1.tex}{\texttt{ifelsefi1.tex}}{lst:ifelsefi1}
	\end{minipage}%
	\hfill
	\begin{minipage}{.54\textwidth}
		\cmhlistingsfromfile{demonstrations/ifelsefi1-default.tex}{\texttt{ifelsefi1.tex} default output}{lst:ifelsefi1-default}
	\end{minipage}

	It is recommended to specify \texttt{noAdditionalIndent} and \texttt{indentRules} in the `scalar' form only
	for these type of code blocks, although the `field' form would work, assuming that \texttt{body} was specified.
	Examples are shown in \cref{lst:ifnum-noAdd,lst:ifnum-indent-rules}.

	\begin{minipage}{.45\textwidth}
		\cmhlistingsfromfile[style=yaml-LST]{demonstrations/ifnum-noAdd.yaml}[width=.8\linewidth,before=\centering,yaml-TCB]{\texttt{ifnum-noAdd.yaml}}{lst:ifnum-noAdd}
	\end{minipage}
	\hfill
	\begin{minipage}{.45\textwidth}
		\cmhlistingsfromfile[style=yaml-LST]{demonstrations/ifnum-indent-rules.yaml}[width=.8\linewidth,before=\centering,yaml-TCB]{\texttt{ifnum-indent-rules.yaml}}{lst:ifnum-indent-rules}
	\end{minipage}

	After running the following commands,
	\begin{commandshell}
latexindent.pl ifelsefi1.tex -local ifnum-noAdd.yaml  
latexindent.pl ifelsefi1.tex -l ifnum-indent-rules.yaml  
\end{commandshell}
	we receive the respective output given in \cref{lst:ifelsefi1-output-noAdd,lst:ifelsefi1-output-indent-rules}; note that
	in \cref{lst:ifelsefi1-output-noAdd}, the \texttt{ifnum} code block has \emph{not} received any additional indentation,
	while in \cref{lst:ifelsefi1-output-indent-rules}, the \texttt{ifnum} code block has received one tab and two spaces of indentation.

	\begin{minipage}{.45\textwidth}
		\cmhlistingsfromfile{demonstrations/ifelsefi1-noAdd.tex}{\texttt{ifelsefi1.tex} using \cref{lst:ifnum-noAdd}}{lst:ifelsefi1-output-noAdd}
	\end{minipage}
	\hfill
	\begin{minipage}{.5\textwidth}
		\cmhlistingsfromfile{demonstrations/ifelsefi1-indent-rules.tex}{\texttt{ifelsefi1.tex} using \cref{lst:ifnum-indent-rules}}{lst:ifelsefi1-output-indent-rules}
	\end{minipage}

	We may specify \texttt{noAdditionalIndentGlobal} and \texttt{indentRulesGlobal} as in \cref{lst:ifelsefi-noAdd-glob,lst:ifelsefi-indent-rules-global}.

	\begin{minipage}{.49\textwidth}
		\cmhlistingsfromfile[style=yaml-LST]{demonstrations/ifelsefi-noAdd-glob.yaml}[width=.9\linewidth,before=\centering,yaml-TCB]{\texttt{ifelsefi-noAdd-glob.yaml}}{lst:ifelsefi-noAdd-glob}
	\end{minipage}
	\hfill
	\begin{minipage}{.49\textwidth}
		\cmhlistingsfromfile[style=yaml-LST]{demonstrations/ifelsefi-indent-rules-global.yaml}[width=.9\linewidth,before=\centering,yaml-TCB]{\texttt{ifelsefi-indent-rules-global.yaml}}{lst:ifelsefi-indent-rules-global}
	\end{minipage}

	Upon running the following commands
	\begin{commandshell}
latexindent.pl ifelsefi1.tex -local ifelsefi-noAdd-glob.yaml  
latexindent.pl ifelsefi1.tex -l ifelsefi-indent-rules-global.yaml  
\end{commandshell}
	we receive the outputs in \cref{lst:ifelsefi1-output-noAdd-glob,lst:ifelsefi1-output-indent-rules-global}; notice that  in
	\cref{lst:ifelsefi1-output-noAdd-glob} neither of the \texttt{ifelsefi} code blocks have received indentation, while in
	\cref{lst:ifelsefi1-output-indent-rules-global} both code blocks have received a single space of indentation.

	\begin{minipage}{.45\textwidth}
		\cmhlistingsfromfile{demonstrations/ifelsefi1-noAdd-glob.tex}{\texttt{ifelsefi1.tex} using \cref{lst:ifelsefi-noAdd-glob}}{lst:ifelsefi1-output-noAdd-glob}
	\end{minipage}
	\hfill
	\begin{minipage}{.45\textwidth}
		\cmhlistingsfromfile{demonstrations/ifelsefi1-indent-rules-global.tex}{\texttt{ifelsefi1.tex} using \cref{lst:ifelsefi-indent-rules-global}}{lst:ifelsefi1-output-indent-rules-global}
	\end{minipage}

% arara: pdflatex: {shell: yes, files: [latexindent]}
\subsubsection{\texttt{specialBeginEnd} code blocks}
	Let's use the example from \vref{lst:specialbefore} which has default output shown in
	\vref{lst:specialafter}.

	It is recommended to specify \texttt{noAdditionalIndent} and \texttt{indentRules} in the `scalar' form
	for these type of code blocks, although the `field' form would work, assuming that \texttt{body} was specified.
	Examples are shown in \cref{lst:displayMath-noAdd,lst:displayMath-indent-rules}.

	\begin{minipage}{.49\textwidth}
		\cmhlistingsfromfile[style=yaml-LST]{demonstrations/displayMath-noAdd.yaml}[width=.9\linewidth,before=\centering,yaml-TCB]{\texttt{displayMath-noAdd.yaml}}{lst:displayMath-noAdd}
	\end{minipage}
	\hfill
	\begin{minipage}{.49\textwidth}
		\cmhlistingsfromfile[style=yaml-LST]{demonstrations/displayMath-indent-rules.yaml}[width=.9\linewidth,before=\centering,yaml-TCB]{\texttt{displayMath-indent-rules.yaml}}{lst:displayMath-indent-rules}
	\end{minipage}

	After running the following commands,
	\begin{commandshell}
latexindent.pl special1.tex -local displayMath-noAdd.yaml  
latexindent.pl special1.tex -l displayMath-indent-rules.yaml  
\end{commandshell}
	we receive the respective output given in \cref{lst:special1-output-noAdd,lst:special1-output-indent-rules}; note that
	in \cref{lst:special1-output-noAdd}, the \texttt{displayMath} code block has \emph{not} received any additional indentation,
	while in \cref{lst:special1-output-indent-rules}, the \texttt{displayMath} code block has received three tabs worth of indentation.

	\begin{minipage}{.45\textwidth}
		\cmhlistingsfromfile{demonstrations/special1-noAdd.tex}{\texttt{special1.tex} using \cref{lst:displayMath-noAdd}}{lst:special1-output-noAdd}
	\end{minipage}
	\hfill
	\begin{minipage}{.45\textwidth}
		\cmhlistingsfromfile{demonstrations/special1-indent-rules.tex}{\texttt{special1.tex} using \cref{lst:displayMath-indent-rules}}{lst:special1-output-indent-rules}
	\end{minipage}

	We may specify \texttt{noAdditionalIndentGlobal} and \texttt{indentRulesGlobal} as in \cref{lst:special-noAdd-glob,lst:special-indent-rules-global}.

	\begin{minipage}{.49\textwidth}
		\cmhlistingsfromfile[style=yaml-LST]{demonstrations/special-noAdd-glob.yaml}[width=.9\linewidth,before=\centering,yaml-TCB]{\texttt{special-noAdd-glob.yaml}}{lst:special-noAdd-glob}
	\end{minipage}
	\hfill
	\begin{minipage}{.49\textwidth}
		\cmhlistingsfromfile[style=yaml-LST]{demonstrations/special-indent-rules-global.yaml}[width=.9\linewidth,before=\centering,yaml-TCB]{\texttt{special-indent-rules-global.yaml}}{lst:special-indent-rules-global}
	\end{minipage}

	Upon running the following commands
	\begin{commandshell}
latexindent.pl special1.tex -local special-noAdd-glob.yaml  
latexindent.pl special1.tex -l special-indent-rules-global.yaml  
\end{commandshell}
	we receive the outputs in \cref{lst:special1-output-noAdd-glob,lst:special1-output-indent-rules-global}; notice that  in
	\cref{lst:special1-output-noAdd-glob} neither of the \texttt{special} code blocks have received indentation, while in
	\cref{lst:special1-output-indent-rules-global} both code blocks have received a single space of indentation.

	\begin{minipage}{.45\textwidth}
		\cmhlistingsfromfile{demonstrations/special1-noAdd-glob.tex}{\texttt{special1.tex} using \cref{lst:special-noAdd-glob}}{lst:special1-output-noAdd-glob}
	\end{minipage}
	\hfill
	\begin{minipage}{.45\textwidth}
		\cmhlistingsfromfile{demonstrations/special1-indent-rules-global.tex}{\texttt{special1.tex} using \cref{lst:special-indent-rules-global}}{lst:special1-output-indent-rules-global}
	\end{minipage}

% arara: pdflatex: {shell: yes, files: [latexindent]}
\subsubsection{\texttt{afterHeading} code blocks}\label{subsubsec-headings-no-add-indent-rules}
	Let's use the example \cref{lst:headings2} for demonstration throughout this \namecref{subsubsec-headings-no-add-indent-rules}.
	As discussed on \cpageref{lst:headings1}, by default \texttt{latexindent.pl} will not add indentation after headings.

	\cmhlistingsfromfile{demonstrations/headings2.tex}{\texttt{headings2.tex}}{lst:headings2}

	On using the YAML file in \cref{lst:headings3yaml} by running the command
	\begin{commandshell}
latexindent.pl headings2.tex -l headings3.yaml      
    \end{commandshell}
	we obtain the output in \cref{lst:headings2-mod3}. Note that the argument of \texttt{paragraph} has received (default) indentation,
	and that the body after the heading statement has received (default) indentation.

	\begin{minipage}{.45\textwidth}
		\cmhlistingsfromfile{demonstrations/headings2-mod3.tex}{\texttt{headings2.tex} using \cref{lst:headings3yaml}}{lst:headings2-mod3}
	\end{minipage}%
	\hfill
	\begin{minipage}{.45\textwidth}
		\cmhlistingsfromfile[style=yaml-LST]{demonstrations/headings3.yaml}[yaml-TCB]{\texttt{headings3.yaml}}{lst:headings3yaml}
	\end{minipage}

	If we specify \texttt{noAdditionalIndent} as in \cref{lst:headings4yaml} and run the command
	\begin{commandshell}
latexindent.pl headings2.tex -l headings4.yaml      
    \end{commandshell}
	then we receive the output in \cref{lst:headings2-mod4}. Note that the arguments \emph{and} the body after the heading
	of \texttt{paragraph} has received no additional indentation, because we have specified \texttt{noAdditionalIndent} in scalar form.

	\begin{minipage}{.45\textwidth}
		\cmhlistingsfromfile{demonstrations/headings2-mod4.tex}{\texttt{headings2.tex} using \cref{lst:headings4yaml}}{lst:headings2-mod4}
	\end{minipage}%
	\hfill
	\begin{minipage}{.45\textwidth}
		\cmhlistingsfromfile[style=yaml-LST]{demonstrations/headings4.yaml}[yaml-TCB]{\texttt{headings4.yaml}}{lst:headings4yaml}
	\end{minipage}

	Similarly, if we specify \texttt{indentRules} as in \cref{lst:headings5yaml} and run analogous commands to those above,
	we receive the output in \cref{lst:headings2-mod5}; note that the \emph{body}, \emph{mandatory argument} and content
	\emph{after the heading} of \texttt{paragraph} have \emph{all} received three tabs worth of indentation.

	\begin{minipage}{.55\textwidth}
		\cmhlistingsfromfile{demonstrations/headings2-mod5.tex}{\texttt{headings2.tex} using \cref{lst:headings5yaml}}{lst:headings2-mod5}
	\end{minipage}%
	\hfill
	\begin{minipage}{.42\textwidth}
		\cmhlistingsfromfile[style=yaml-LST]{demonstrations/headings5.yaml}[yaml-TCB]{\texttt{headings5.yaml}}{lst:headings5yaml}
	\end{minipage}

	We may, instead, specify \texttt{noAdditionalIndent} in `field' form, as in \cref{lst:headings6yaml} which gives the output in \cref{lst:headings2-mod6}.

	\begin{minipage}{.45\textwidth}
		\cmhlistingsfromfile{demonstrations/headings2-mod6.tex}{\texttt{headings2.tex} using \cref{lst:headings6yaml}}{lst:headings2-mod6}
	\end{minipage}%
	\hfill
	\begin{minipage}{.45\textwidth}
		\cmhlistingsfromfile[style=yaml-LST]{demonstrations/headings6.yaml}[yaml-TCB]{\texttt{headings6.yaml}}{lst:headings6yaml}
	\end{minipage}

	Analogously, we may specify \texttt{indentRules} as in \cref{lst:headings7yaml} which gives the output in \cref{lst:headings2-mod7};
	note that mandatory argument text has only received a single space of indentation, while the body after the heading has
	received three tabs worth of indentation.

	\begin{minipage}{.45\textwidth}
		\cmhlistingsfromfile{demonstrations/headings2-mod7.tex}{\texttt{headings2.tex} using \cref{lst:headings7yaml}}{lst:headings2-mod7}
	\end{minipage}%
	\hfill
	\begin{minipage}{.45\textwidth}
		\cmhlistingsfromfile[style=yaml-LST]{demonstrations/headings7.yaml}[yaml-TCB]{\texttt{headings7.yaml}}{lst:headings7yaml}
	\end{minipage}

	Finally, let's consider \texttt{noAdditionalIndentGlobal} and \texttt{indentRulesGlobal} shown in \cref{lst:headings8yaml,lst:headings9yaml}
	respectively, with respective output in \cref{lst:headings2-mod8,lst:headings2-mod9}. Note that in \cref{lst:headings8yaml} the
	\emph{mandatory argument} of \texttt{paragraph} has received a (default) tab's worth of indentation, while the body after the
	heading has received \emph{no additional indentation}. Similarly, in \cref{lst:headings2-mod9}, the \emph{argument} has received both a
	(default) tab plus two spaces of indentation (from the global rule specified in \cref{lst:headings9yaml}), and the remaining body
	after \texttt{paragraph} has received just two spaces of indentation.

	\begin{minipage}{.45\textwidth}
		\cmhlistingsfromfile{demonstrations/headings2-mod8.tex}{\texttt{headings2.tex} using \cref{lst:headings8yaml}}{lst:headings2-mod8}
	\end{minipage}%
	\hfill
	\begin{minipage}{.45\textwidth}
		\cmhlistingsfromfile[style=yaml-LST]{demonstrations/headings8.yaml}[yaml-TCB]{\texttt{headings8.yaml}}{lst:headings8yaml}
	\end{minipage}

	\begin{minipage}{.45\textwidth}
		\cmhlistingsfromfile{demonstrations/headings2-mod9.tex}{\texttt{headings2.tex} using \cref{lst:headings9yaml}}{lst:headings2-mod9}
	\end{minipage}%
	\hfill
	\begin{minipage}{.45\textwidth}
		\cmhlistingsfromfile[style=yaml-LST]{demonstrations/headings9.yaml}[yaml-TCB]{\texttt{headings9.yaml}}{lst:headings9yaml}
	\end{minipage}

% arara: pdflatex: {shell: yes, files: [latexindent]}
\subsubsection{The remaining code blocks}
	Referencing the different types of code blocks in \vref{tab:code-blocks}, we have a few
	code blocks yet to cover; these are very similar to the \texttt{commands} code block type
	covered comprehensively in \vref{subsubsec:commands-arguments}, but a small discussion
	defining these remaining code blocks is necessary.

	\paragraph{\texttt{keyEqualsValuesBracesBrackets}}
		\texttt{latexindent.pl} defines this type of code block by the following criteria:
		\begin{itemize}
			\item it must immediately follow either \lstinline!{! OR \lstinline![! OR \lstinline!,! with comments
			      and blank lines allowed;
			\item then it has a name made up of the characters detailed in \vref{tab:code-blocks};
			\item then an $=$ symbol;
			\item then at least one set of curly braces or square brackets (comments and line breaks allowed throughout).
		\end{itemize}

		An example is shown in \cref{lst:pgfkeysbefore}, with the default output given in \cref{lst:pgfkeys1:default}.

		\begin{minipage}{.45\textwidth}
			\cmhlistingsfromfile{demonstrations/pgfkeys1.tex}{\texttt{pgfkeys1.tex}}{lst:pgfkeysbefore}
		\end{minipage}%
		\hfill
		\begin{minipage}{.5\textwidth}
			\cmhlistingsfromfile{demonstrations/pgfkeys1-default.tex}{\texttt{pgfkeys1.tex} default output}{lst:pgfkeys1:default}
		\end{minipage}%

		In \cref{lst:pgfkeys1:default}, note that the maximum indentation is three tabs, and these come from:
		\begin{itemize}
			\item the \lstinline!\pgfkeys! command's mandatory argument;
			\item the \lstinline!start coordinate/.initial! key's mandatory argument;
			\item the \lstinline!start coordinate/.initial! key's body, which is defined as any lines following the name of the
			      key that include its arguments.  This is the part controlled by the \emph{body} field for \texttt{noAdditionalIndent}
			      and friends from \cpageref{sec:noadd-indent-rules}.
		\end{itemize}
	\paragraph{\texttt{namedGroupingBracesBrackets}}
		This type of code block is mostly motivated by tikz-based code; we define this code block as follows:
		\begin{itemize}
			\item it must immediately follow either \emph{horizontal space} OR \emph{one or more line breaks} OR \lstinline!{! OR \lstinline![!
			      OR \lstinline!$!;
			\item the name may contain the characters detailed in \vref{tab:code-blocks};
			\item then at least one set of curly braces or square brackets (comments and line breaks allowed throughout).
		\end{itemize}
		A simple example is given in \cref{lst:child1}, with default output in \cref{lst:child1:default}.

		\begin{minipage}{.45\textwidth}
			\cmhlistingsfromfile{demonstrations/child1.tex}{\texttt{child1.tex}}{lst:child1}
		\end{minipage}%
		\hfill
		\begin{minipage}{.5\textwidth}
			\cmhlistingsfromfile{demonstrations/child1-default.tex}{\texttt{child1.tex} default output}{lst:child1:default}
		\end{minipage}%

		In particular, \texttt{latexindent.pl} considers \texttt{child}, \texttt{parent} and \texttt{node} all to be \texttt{namedGroupingBracesBrackets}\footnote{
			You may like to verify this by using the \texttt{-tt} option and checking \texttt{indent.log}! }.
		Referencing \cref{lst:child1:default},
		note that the maximum indentation is two tabs, and these come from:
		\begin{itemize}
			\item the \lstinline!child!'s mandatory argument;
			\item the \lstinline!child!'s body, which is defined as any lines following the name of the \texttt{namedGroupingBracesBrackets}
			      that include its arguments. This is the part controlled by the \emph{body} field for \texttt{noAdditionalIndent}
			      and friends from \cpageref{sec:noadd-indent-rules}.
		\end{itemize}

	\paragraph{\texttt{UnNamedGroupingBracesBrackets}} occur in a variety of situations; specifically, we define
		this type of code block as satisfying the following criteria:
		\begin{itemize}
			\item it must immediately follow either \lstinline!{! OR \lstinline![! OR \lstinline!,! OR \lstinline!&! OR \lstinline!)! OR \lstinline!(!
			      OR \lstinline!$!;
			\item then at least one set of curly braces or square brackets (comments and line breaks allowed throughout).
		\end{itemize}

		An example is shown in \cref{lst:psforeach1} with default output give in \cref{lst:psforeach:default}.

		\begin{minipage}{.45\textwidth}
			\cmhlistingsfromfile{demonstrations/psforeach1.tex}{\texttt{psforeach1.tex}}{lst:psforeach1}
		\end{minipage}%
		\hfill
		\begin{minipage}{.5\textwidth}
			\cmhlistingsfromfile{demonstrations/psforeach1-default.tex}{\texttt{psforeach1.tex} default output}{lst:psforeach:default}
		\end{minipage}%

		Referencing \cref{lst:psforeach:default}, there are \emph{three} sets of unnamed braces. Note also that the maximum value
		of indentation is three tabs, and these come from:
		\begin{itemize}
			\item the \lstinline!\psforeach! command's mandatory argument;
			\item the \emph{first} un-named braces mandatory argument;
			\item the \emph{first} un-named braces \emph{body}, which we define as any lines following the first opening \lstinline!{! or \lstinline![!
			      that defined the code block.  This is the part controlled by the \emph{body} field for \texttt{noAdditionalIndent}
			      and friends from \cpageref{sec:noadd-indent-rules}.
		\end{itemize}
		Users wishing to customise the mandatory and/or optional arguments on a \emph{per-name} basis for the \texttt{UnNamedGroupingBracesBrackets}
		should use \texttt{always-un-named}.

	\paragraph{\texttt{filecontents}} code blocks behave just as \texttt{environments}, except that neither arguments nor items are sought.

\subsubsection{Summary}
	Having considered all of the different types of code blocks, the functions of the fields given in
	\cref{lst:noAdditionalIndentGlobal,lst:indentRulesGlobal} should now make sense.

	\begin{widepage}
		\begin{minipage}{.47\linewidth}
			\cmhlistingsfromfile[style=noAdditionalIndentGlobal]{../defaultSettings.yaml}[before=\centering,yaml-TCB]{\texttt{noAdditionalIndentGlobal}}{lst:noAdditionalIndentGlobal}
		\end{minipage}%
		\hfill
		\begin{minipage}{.47\linewidth}
			\cmhlistingsfromfile[style=indentRulesGlobal]{../defaultSettings.yaml}[before=\centering,yaml-TCB]{\texttt{indentRulesGlobal}}{lst:indentRulesGlobal}
		\end{minipage}%
	\end{widepage}

\stopcontents[noAdditionalIndent]
% arara: pdflatex: {shell: yes, files: [latexindent]}
\subsection{Commands and the strings between their arguments}\label{subsec:commands-string-between}
	The \texttt{command} code blocks will always look for optional (square bracketed) and
	mandatory (curly braced) arguments which can contain comments, line breaks and
	`beamer' commands \lstinline!<.*?>!  between them. There are switches that can allow them to contain
	other strings, which we discuss next.

\yamltitle{commandCodeBlocks}*{fields}

	The \texttt{commandCodeBlocks} field contains a few switches detailed in \cref{lst:commandCodeBlocks}.

	\cmhlistingsfromfile[style=commandCodeBlocks]{../defaultSettings.yaml}[width=.8\linewidth,before=\centering,yaml-TCB]{\texttt{commandCodeBlocks}}{lst:commandCodeBlocks}

\yamltitle{roundParenthesesAllowed}{0|1}

	The need for this field was mostly motivated by commands found in code used to generate images in \texttt{PSTricks} and \texttt{tikz}; for example,
	let's consider the code given in \cref{lst:pstricks1}.

	\begin{minipage}{.45\textwidth}
		\cmhlistingsfromfile[columns=fixed]{demonstrations/pstricks1.tex}{\texttt{pstricks1.tex}}{lst:pstricks1}
	\end{minipage}
	\hfill
	\begin{minipage}{.45\textwidth}
		\cmhlistingsfromfile[columns=fixed]{demonstrations/pstricks1-default.tex}{\texttt{pstricks1} default output}{lst:pstricks1-default}
	\end{minipage}

	Notice that the \lstinline!\defFunction! command has an optional argument, followed by a
	mandatory argument, followed by a round-parenthesis argument, $(u,v)$.

	By default, because \texttt{roundParenthesesAllowed} is set to $1$ in \cref{lst:commandCodeBlocks}, then \texttt{latexindent.pl}
	will allow round parenthesis between optional and mandatory arguments. In the case of the code in \cref{lst:pstricks1},
	\texttt{latexindent.pl} finds \emph{all} the arguments of \lstinline!defFunction!, both before and after \lstinline!(u,v)!.

	The default output from running \texttt{latexindent.pl} on \cref{lst:pstricks1} actually leaves it unchanged (see \cref{lst:pstricks1-default});
	note in particular, this is because of \texttt{noAdditionalIndentGlobal} as discussed on \cpageref{page:command:noAddGlobal}.

	Upon using the YAML settings in \cref{lst:noRoundParentheses}, and running the command
	\begin{commandshell}
latexindent.pl pstricks1.tex -l noRoundParentheses.yaml
        \end{commandshell}
	we obtain the output given in  \cref{lst:pstricks1-nrp}.

	\begin{minipage}{.45\textwidth}
		\cmhlistingsfromfile{demonstrations/pstricks1-nrp.tex}{\texttt{pstricks1.tex} using \cref{lst:noRoundParentheses}}{lst:pstricks1-nrp}
	\end{minipage}
	\hfill
	\begin{minipage}{.45\textwidth}
		\cmhlistingsfromfile[style=yaml-LST]{demonstrations/noRoundParentheses.yaml}[yaml-TCB]{\texttt{noRoundParentheses.yaml}}{lst:noRoundParentheses}
	\end{minipage}

	Notice the difference between \cref{lst:pstricks1-default} and \cref{lst:pstricks1-nrp}; in particular, in \cref{lst:pstricks1-nrp}, because
	round parentheses are \emph{not} allowed, \texttt{latexindent.pl} finds that the \lstinline!\defFunction! command finishes at the first opening
	round parenthesis. As such, the remaining braced, mandatory, arguments are found to be \texttt{UnNamedGroupingBracesBrackets} (see \vref{tab:code-blocks})
	which, by default, assume indentation for their body, and hence the tabbed indentation in \cref{lst:pstricks1-nrp}.

	Let's explore this using the YAML given in \cref{lst:defFunction} and run the command
	\begin{commandshell}
latexindent.pl pstricks1.tex -l defFunction.yaml
        \end{commandshell}
	then the output is as in \cref{lst:pstricks1-indent-rules}.

	\begin{minipage}{.45\textwidth}
		\cmhlistingsfromfile{demonstrations/pstricks1-indent-rules.tex}{\texttt{pstricks1.tex} using \cref{lst:defFunction}}{lst:pstricks1-indent-rules}
	\end{minipage}
	\hfill
	\begin{minipage}{.45\textwidth}
		\cmhlistingsfromfile[style=yaml-LST]{demonstrations/defFunction.yaml}[yaml-TCB]{\texttt{defFunction.yaml}}{lst:defFunction}
	\end{minipage}

	Notice in \cref{lst:pstricks1-indent-rules} that the \emph{body} of the \lstinline!defFunction! command i.e, the subsequent lines
	containing arguments after the command name, have received the single space of indentation specified by \cref{lst:defFunction}.

\yamltitle{stringsAllowedBetweenArguments}*{fields}
	\texttt{tikz} users may well specify code such as that given in \cref{lst:tikz-node1}; processing this code using
	\texttt{latexindent.pl} gives the default output in \cref{lst:tikz-node1-default}.

	\begin{minipage}{.45\textwidth}
		\cmhlistingsfromfile[columns=fixed]{demonstrations/tikz-node1.tex}{\texttt{tikz-node1.tex}}{lst:tikz-node1}
	\end{minipage}
	\hfill
	\begin{minipage}{.45\textwidth}
		\cmhlistingsfromfile[columns=fixed]{demonstrations/tikz-node1-default.tex}{\texttt{tikz-node1} default output}{lst:tikz-node1-default}
	\end{minipage}

	With reference to \vref{lst:commandCodeBlocks}, we see that the strings
	\begin{quote}
		to, node, ++
	\end{quote}
	are all allowed to appear between arguments, as they are each set to $1$; importantly, you are encouraged to add further names
	to this field as necessary. This means that when \texttt{latexindent.pl}
	processes \cref{lst:tikz-node1}, it consumes:
	\begin{itemize}
		\item the optional argument \lstinline![thin]!
		\item the round-bracketed argument \lstinline!(c)! because \texttt{roundParenthesesAllowed} is $1$ by default
		\item the string \lstinline!to! (specified in \texttt{stringsAllowedBetweenArguments})
		\item the optional argument \lstinline![in=110,out=-90]!
		\item the string \lstinline!++! (specified in \texttt{stringsAllowedBetweenArguments})
		\item the round-bracketed argument \lstinline!(0,-0.5cm)! because \texttt{roundParenthesesAllowed} is $1$ by default
		\item the string \lstinline!node! (specified in \texttt{stringsAllowedBetweenArguments})
		\item the optional argument \lstinline![below,align=left,scale=0.5]!
	\end{itemize}

	We can explore this further, for example using \cref{lst:draw} and running the command
	\begin{commandshell}
latexindent.pl tikz-node1.tex -l draw.yaml  
\end{commandshell}
	we receive the output given in \cref{lst:tikz-node1-draw}.

	\begin{minipage}{.45\textwidth}
		\cmhlistingsfromfile{demonstrations/tikz-node1-draw.tex}{\texttt{tikz-node1.tex} using \cref{lst:draw}}{lst:tikz-node1-draw}
	\end{minipage}
	\hfill
	\begin{minipage}{.45\textwidth}
		\cmhlistingsfromfile[style=yaml-LST]{demonstrations/draw.yaml}[yaml-TCB]{\texttt{draw.yaml}}{lst:draw}
	\end{minipage}

	Notice that each line after the \lstinline!\draw! command (its `body') in \cref{lst:tikz-node1-draw} has been given the
	appropriate two-spaces worth of indentation specified in \cref{lst:draw}.

	Let's compare this with the output from using the YAML settings in \cref{lst:no-to}, and running the command
	\begin{commandshell}
latexindent.pl tikz-node1.tex -l no-to.yaml  
\end{commandshell}
	given in \cref{lst:tikz-node1-no-to}.

	\begin{minipage}{.45\textwidth}
		\cmhlistingsfromfile{demonstrations/tikz-node1-no-to.tex}{\texttt{tikz-node1.tex} using \cref{lst:no-to}}{lst:tikz-node1-no-to}
	\end{minipage}
	\hfill
	\begin{minipage}{.45\textwidth}
		\cmhlistingsfromfile[style=yaml-LST]{demonstrations/no-to.yaml}[yaml-TCB]{\texttt{no-to.yaml}}{lst:no-to}
	\end{minipage}

	In this case, \texttt{latexindent.pl} sees that:
	\begin{itemize}
		\item the \lstinline!\draw! command finishes after the \lstinline!(c)! as (\texttt{stringsAllowedBetweenArguments} has \texttt{to} set to $0$)
		\item it finds a \texttt{namedGroupingBracesBrackets} called \texttt{to} (see \vref{tab:code-blocks}) \emph{with} argument \lstinline![in=110,out=-90]!
		\item it finds another \texttt{namedGroupingBracesBrackets} but this time called \texttt{node} with argument \lstinline![below,align=left,scale=0.5]!
	\end{itemize}

\clearpage
% arara: pdflatex: {shell: yes, files: [latexindent]}
\fancyhead[R]{\bfseries\thepage%
	\tikz[remember picture,overlay] {
		\node at (1,0){\includegraphics{logo}};
	}}
\section{The \texttt{-m} (\texttt{modifylinebreaks}) switch}\label{sec:modifylinebreaks}
 All features described in this section will only be relevant if the \texttt{-m} switch
 is used.

\yamltitle{modifylinebreaks}*{fields}
	\begin{wrapfigure}[7]{r}[0pt]{8cm}
		\cmhlistingsfromfile[style=modifylinebreaks]{../defaultSettings.yaml}[MLB-TCB,width=.85\linewidth,before=\centering]{\texttt{modifyLineBreaks}}{lst:modifylinebreaks}
	\end{wrapfigure}
	\makebox[0pt][r]{%
		\raisebox{-\totalheight}[0pt][0pt]{%
			\tikz\node[opacity=1] at (0,0) {\includegraphics[width=4cm]{logo}};}}%	
	As of Version 3.0, \texttt{latexindent.pl} has the \texttt{-m} switch, which
	permits \texttt{latexindent.pl} to modify line breaks, according to the
	specifications in the \texttt{modifyLineBreaks} field. \emph{The settings
		in this field will only be considered if the \texttt{-m} switch has been used}.
	A snippet of the default settings of this field is shown in \cref{lst:modifylinebreaks}.

	Having read the previous paragraph, it should sound reasonable that, if you call \texttt{latexindent.pl}
	using the \texttt{-m} switch, then you give it permission to modify line breaks in your file,
	but let's be clear:

	\begin{warning}
		If you call \texttt{latexindent.pl} with the \texttt{-m} switch, then you
		are giving it permission to modify line breaks. By default, the only
		thing that will happen is that multiple blank lines will be condensed into
		one blank line; many other settings are possible, discussed next.
	\end{warning}

\yamltitle{preserveBlankLines}{0|1}
	This field is directly related to \emph{poly-switches}, discussed below.
	By default, it is set to \texttt{1}, which means that blank lines will
	be protected from removal; however, regardless of this setting, multiple
	blank lines can be condensed if \texttt{condenseMultipleBlankLinesInto} is
	greater than \texttt{0}, discussed next.

\yamltitle{condenseMultipleBlankLinesInto}*{integer $\geq 0$}
	Assuming that this switch takes an integer value greater than \texttt{0}, \texttt{latexindent.pl} will condense multiple blank lines into
	the number of blank lines illustrated by this switch. As an example, \cref{lst:mlb-bl} shows a sample file
	with blank lines; upon running
	\begin{commandshell}
latexindent.pl myfile.tex -m  
\end{commandshell}
	the output is shown in \cref{lst:mlb-bl-out}; note that the multiple blank lines have been
	condensed into one blank line, and note also that we have used the \texttt{-m} switch!

	\begin{minipage}{.45\textwidth}
		\cmhlistingsfromfile{demonstrations/mlb1.tex}{\texttt{mlb1.tex}}{lst:mlb-bl}
	\end{minipage}%
	\hfill
	\begin{minipage}{.45\textwidth}
		\cmhlistingsfromfile{demonstrations/mlb1-out.tex}{\texttt{mlb1.tex} out output}{lst:mlb-bl-out}
	\end{minipage}

\yamltitle{textWrapOptions}*{fields}
	When the \texttt{-m} switch is active \texttt{latexindent.pl} has the ability to wrap text using the options
	specified in the \texttt{textWrapOptions} field, see \cref{lst:textWrapOptions}. The value of
	\texttt{columns} specifies the column at which the text should be wrapped.
	By default, the value of \texttt{columns} is \texttt{0}, so \texttt{latexindent.pl}
	will \emph{not} wrap text; if you change it to a value of \texttt{2} or more, then
	text will be wrapped after the character in the specified column.

	\cmhlistingsfromfile[style=textWrapOptions]{../defaultSettings.yaml}[MLB-TCB,width=.85\linewidth,before=\centering]{\texttt{textWrapOptions}}{lst:textWrapOptions}

	For example, consider the file give in \cref{lst:textwrap1}.

	\begin{widepage}
		\cmhlistingsfromfile{demonstrations/textwrap1.tex}{\texttt{textwrap1.tex}}{lst:textwrap1}
	\end{widepage}

	Using the file \texttt{textwrap1.yaml} in \cref{lst:textwrap1-yaml}, and running the command
	\begin{commandshell}
latexindent.pl -m textwrap1.tex -o textwrap1-mod1.tex -l textwrap1.yaml
\end{commandshell}
	we obtain the output in \cref{lst:textwrap1-mod1}.

	\begin{minipage}{.45\linewidth}
		\cmhlistingsfromfile{demonstrations/textwrap1-mod1.tex}{\texttt{textwrap1-mod1.tex}}{lst:textwrap1-mod1}
	\end{minipage}
	\hfill
	\begin{minipage}{.45\linewidth}
		\cmhlistingsfromfile{demonstrations/textwrap1.yaml}[MLB-TCB]{\texttt{textwrap1.yaml}}{lst:textwrap1-yaml}
	\end{minipage}

	The text wrapping routine is performed \emph{after} verbatim environments have been stored, so verbatim
	environments and verbatim commands are exempt from the routine. For example, using the file in
	\cref{lst:textwrap2},
	\begin{widepage}
		\cmhlistingsfromfile{demonstrations/textwrap2.tex}{\texttt{textwrap2.tex}}{lst:textwrap2}
	\end{widepage}
	and running the following command and continuing to use \texttt{textwrap1.yaml} from \cref{lst:textwrap1-yaml},
	\begin{commandshell}
latexindent.pl -m textwrap2.tex -o textwrap2-mod1.tex -l textwrap1.yaml
\end{commandshell}
	then the output is as in \cref{lst:textwrap2-mod1}.
	\begin{widepage}
		\cmhlistingsfromfile{demonstrations/textwrap2-mod1.tex}{\texttt{textwrap2-mod1.tex}}{lst:textwrap2-mod1}
	\end{widepage}
	Furthermore, the text wrapping routine is performed after the trailing comments have been
	stored, and they are also exempt from text wrapping. For example, using the file in \cref{lst:textwrap3}
	\begin{widepage}
		\cmhlistingsfromfile{demonstrations/textwrap3.tex}{\texttt{textwrap3.tex}}{lst:textwrap3}
	\end{widepage}
	and running the following command and continuing to use \texttt{textwrap1.yaml} from \cref{lst:textwrap1-yaml},
	\begin{commandshell}
latexindent.pl -m textwrap3.tex -o textwrap3-mod1.tex -l textwrap1.yaml
\end{commandshell}
	then the output is as in \cref{lst:textwrap3-mod1}.

	\cmhlistingsfromfile{demonstrations/textwrap3-mod1.tex}{\texttt{textwrap3-mod1.tex}}{lst:textwrap3-mod1}

	\begin{wrapfigure}[6]{r}[0pt]{8cm}
		\cmhlistingsfromfile[style=textWrapOptionsAll]{../defaultSettings.yaml}[MLB-TCB,width=.85\linewidth,before=\centering]{\texttt{textWrapOptions}}{lst:textWrapOptionsAll}
	\end{wrapfigure}
	The text wrapping routine of \texttt{latexindent.pl} is performed by the \texttt{Text::Wrap} module, which provides a
	\texttt{separator} feature to separate lines with characters other than a new line (see \cite{textwrap}). By default,
	the separator is empty (see \cref{lst:textWrapOptionsAll}) which means that a new line token will be used, but you can change it as you see fit.

	For example starting with the file in \cref{lst:textwrap4}

	\cmhlistingsfromfile{demonstrations/textwrap4.tex}{\texttt{textwrap4.tex}}{lst:textwrap4}
	and using \texttt{textwrap2.yaml} from \cref{lst:textwrap2-yaml} with the following command
	\begin{commandshell}
latexindent.pl -m textwrap4.tex -o textwrap4-mod2.tex -l textwrap2.yaml
\end{commandshell}
	then we obtain the output in \cref{lst:textwrap4-mod2}.

	\begin{minipage}{.45\linewidth}
		\cmhlistingsfromfile{demonstrations/textwrap4-mod2.tex}{\texttt{textwrap4-mod2.tex}}{lst:textwrap4-mod2}
	\end{minipage}
	\hfill
	\begin{minipage}{.45\linewidth}
		\cmhlistingsfromfile{demonstrations/textwrap2.yaml}[MLB-TCB]{\texttt{textwrap2.yaml}}{lst:textwrap2-yaml}
	\end{minipage}

	\paragraph{Summary of text wrapping}
		It is important to note the following:
		\begin{itemize}
			\item Verbatim environments (\vref{lst:verbatimEnvironments}) and verbatim commands (\vref{lst:verbatimCommands}) will \emph{not} be affected by the text wrapping routine (see \vref{lst:textwrap2-mod1});
			\item comments will \emph{not} be affected by the text wrapping routine (see \vref{lst:textwrap3-mod1});
			\item indentation is performed \emph{after} the text wrapping routine; as such, indented code
			      will likely exceed any maximum value set in the \texttt{columns} field.
		\end{itemize}

\yamltitle{removeParagraphLineBreaks}*{fields}
	When the \texttt{-m} switch is active \texttt{latexindent.pl} has the ability to remove line breaks
	from within paragraphs; the behaviour is controlled by the \texttt{removeParagraphLineBreaks} field, detailed in
	\cref{lst:removeParagraphLineBreaks}. Thank you to \cite{jowens} for shaping and assisting with the testing of this feature.

	\cmhlistingsfromfile[style=removeParagraphLineBreaks]{../defaultSettings.yaml}[MLB-TCB,width=.85\linewidth,before=\centering]{\texttt{removeParagraphLineBreaks}}{lst:removeParagraphLineBreaks}

	This routine can be turned on \emph{globally} for \emph{every} code block type known to \texttt{latexindent.pl}
	(see \vref{tab:code-blocks}) by using the \texttt{all} switch; by default, this switch is \emph{off}. Assuming
	that the \texttt{all} switch is off, then the routine can be controlled on a per-code-block-type basis, and
	within that, on a per-name basis. We will consider examples of each of these in turn, but
	before we do, let's specify what \texttt{latexindent.pl} considers as a paragraph:
	\begin{itemize}
		\item it must begin on its own line with either an alphabetic or numeric character, and not with any of the
		      code-block types detailed in \vref{tab:code-blocks};
		\item it can include line breaks, but finishes when it meets either a blank line, a \lstinline!\par!
		      command, or any of the user-specified settings in the \texttt{paragraphsStopAt} field,
		      detailed in \vref{lst:paragraphsStopAt}.
	\end{itemize}

	Let's start with the \texttt{.tex} file in \cref{lst:shortlines}, together with the YAML settings in
	\cref{lst:remove-para1-yaml}.

	\begin{minipage}{.45\linewidth}
		\cmhlistingsfromfile{demonstrations/shortlines.tex}{\texttt{shortlines.tex}}{lst:shortlines}
	\end{minipage}
	\hfill
	\begin{minipage}{.49\linewidth}
		\cmhlistingsfromfile{demonstrations/remove-para1.yaml}[MLB-TCB]{\texttt{remove-para1.yaml}}{lst:remove-para1-yaml}
	\end{minipage}

	Upon running the command
	\begin{commandshell}
latexindent.pl -m shortlines.tex -o shortlines1.tex -l remove-para1.yaml
\end{commandshell}
	then we obtain the output given in \cref{lst:shortlines1}.

	\cmhlistingsfromfile{demonstrations/shortlines1.tex}{\texttt{shortlines1.tex}}{lst:shortlines1}

	Keen readers may notice that some trailing white space must be present in the file in \cref{lst:shortlines} which
	has crept in to the output in \cref{lst:shortlines1}. This can be fixed using the YAML file in
	\vref{lst:removeTWS-before} and running, for example,
	\begin{commandshell}
latexindent.pl -m shortlines.tex -o shortlines1-tws.tex -l remove-para1.yaml,removeTWS-before.yaml  
    \end{commandshell}
	in which case the output is as in \cref{lst:shortlines1-tws}; notice that the double spaces present in \cref{lst:shortlines1}
	have been addressed.

	\cmhlistingsfromfile{demonstrations/shortlines1-tws.tex}{\texttt{shortlines1-tws.tex}}{lst:shortlines1-tws}

	Keeping with the settings in \cref{lst:remove-para1-yaml}, we note that the \texttt{all} switch applies
	to \emph{all} code block types. So, for example, let's consider the files in \cref{lst:shortlines-mand,lst:shortlines-opt}

	\begin{minipage}{.45\linewidth}
		\cmhlistingsfromfile{demonstrations/shortlines-mand.tex}{\texttt{shortlines-mand.tex}}{lst:shortlines-mand}
	\end{minipage}
	\hfill
	\begin{minipage}{.45\linewidth}
		\cmhlistingsfromfile{demonstrations/shortlines-opt.tex}{\texttt{shortlines-opt.tex}}{lst:shortlines-opt}
	\end{minipage}

	Upon running the commands
	\begin{widepage}
		\begin{commandshell}
latexindent.pl -m shortlines-mand.tex -o shortlines-mand1.tex -l remove-para1.yaml
latexindent.pl -m shortlines-opt.tex -o shortlines-opt1.tex -l remove-para1.yaml
\end{commandshell}
	\end{widepage}

	then we obtain the respective output given in \cref{lst:shortlines-mand1,lst:shortlines-opt1}.

	\cmhlistingsfromfile{demonstrations/shortlines-mand1.tex}{\texttt{shortlines-mand1.tex}}{lst:shortlines-mand1}
	\cmhlistingsfromfile{demonstrations/shortlines-opt1.tex}{\texttt{shortlines-opt1.tex}}{lst:shortlines-opt1}

	Assuming that we turn \emph{off} the \texttt{all} switch (by setting it to \texttt{0}), then
	we can control the behaviour of \texttt{removeParagraphLineBreaks} either on a per-code-block-type basis,
	or on a per-name basis.

	For example, let's use the code in \cref{lst:shortlines-envs}, and consider the settings in \cref{lst:remove-para2-yaml,lst:remove-para3-yaml};
	note that in \cref{lst:remove-para2-yaml} we specify that \emph{every} environment should receive
	treatment from the routine, while in \cref{lst:remove-para3-yaml} we specify that \emph{only} the
	\texttt{one} environment should receive the treatment.

	\begin{minipage}{.45\linewidth}
		\cmhlistingsfromfile{demonstrations/shortlines-envs.tex}{\texttt{shortlines-envs.tex}}{lst:shortlines-envs}
	\end{minipage}
	\hfill
	\begin{minipage}{.49\linewidth}
		\cmhlistingsfromfile{demonstrations/remove-para2.yaml}[MLB-TCB]{\texttt{remove-para2.yaml}}{lst:remove-para2-yaml}
		\cmhlistingsfromfile{demonstrations/remove-para3.yaml}[MLB-TCB]{\texttt{remove-para3.yaml}}{lst:remove-para3-yaml}
	\end{minipage}

	Upon running the commands
	\begin{widepage}
		\begin{commandshell}
latexindent.pl -m shortlines-envs.tex -o shortlines-envs2.tex -l remove-para2.yaml
latexindent.pl -m shortlines-envs.tex -o shortlines-envs3.tex -l remove-para3.yaml
\end{commandshell}
	\end{widepage}
	then we obtain the respective output given in \cref{lst:shortlines-envs2,lst:shortlines-envs3}.

	\cmhlistingsfromfile{demonstrations/shortlines-envs2.tex}{\texttt{shortlines-envs2.tex}}{lst:shortlines-envs2}
	\cmhlistingsfromfile{demonstrations/shortlines-envs3.tex}{\texttt{shortlines-envs3.tex}}{lst:shortlines-envs3}

	The remaining code-block types can be customized in analogous ways, although note that \texttt{commands},
	\texttt{keyEqualsValuesBracesBrackets}, \texttt{namedGroupingBracesBrackets}, \texttt{UnNamedGroupingBracesBrackets}
	are controlled by the \texttt{optionalArguments} and the \texttt{mandatoryArguments}.

	The only special case is the \texttt{masterDocument} field; this is designed for `chapter'-type files that
	may contain paragraphs that are not within any other code-blocks. For example, consider the file in
	\cref{lst:shortlines-md}, with the YAML settings in \cref{lst:remove-para4-yaml}.

	\begin{minipage}{.45\linewidth}
		\cmhlistingsfromfile{demonstrations/shortlines-md.tex}{\texttt{shortlines-md.tex}}{lst:shortlines-md}
	\end{minipage}
	\hfill
	\begin{minipage}{.49\linewidth}
		\cmhlistingsfromfile{demonstrations/remove-para4.yaml}[MLB-TCB]{\texttt{remove-para4.yaml}}{lst:remove-para4-yaml}
	\end{minipage}

	Upon running the following command
	\begin{widepage}
		\begin{commandshell}
latexindent.pl -m shortlines-md.tex -o shortlines-md4.tex -l remove-para4.yaml
\end{commandshell}
	\end{widepage}
	then we obtain the output in \cref{lst:shortlines-md4}.
	\cmhlistingsfromfile{demonstrations/shortlines-md4.tex}{\texttt{shortlines-md4.tex}}{lst:shortlines-md4}

\yamltitle{paragraphsStopAt}*{fields}
	The paragraph line break routine considers blank lines and the \lstinline|\par| command to be the end of a paragraph;
	you can fine tune the behaviour of the routine further by using the \texttt{paragraphsStopAt} fields, shown in \cref{lst:paragraphsStopAt}.

	\cmhlistingsfromfile[style=paragraphsStopAt]{../defaultSettings.yaml}[MLB-TCB,width=.85\linewidth,before=\centering]{\texttt{paragraphsStopAt}}{lst:paragraphsStopAt}

	The fields specified in \texttt{paragraphsStopAt} tell \texttt{latexindent.pl} to stop the current paragraph
	when it reaches a line that \emph{begins} with any of the code-block types specified as \texttt{1} in \cref{lst:paragraphsStopAt}.
	By default, you'll see that the paragraph line break routine will stop when it reaches an environment at the
	beginning of a line. It is \emph{not} possible to specify these fields on a per-name basis.

	Let's use the \texttt{.tex} file in \cref{lst:sl-stop}; we will, in turn, consider the settings in
	\cref{lst:stop-command-yaml,lst:stop-comment-yaml}.

	\begin{minipage}{.45\linewidth}
		\cmhlistingsfromfile{demonstrations/sl-stop.tex}{\texttt{sl-stop.tex}}{lst:sl-stop}
	\end{minipage}
	\hfill
	\begin{minipage}{.49\linewidth}
		\cmhlistingsfromfile{demonstrations/stop-command.yaml}[MLB-TCB]{\texttt{stop-command.yaml}}{lst:stop-command-yaml}

		\cmhlistingsfromfile{demonstrations/stop-comment.yaml}[MLB-TCB]{\texttt{stop-comment.yaml}}{lst:stop-comment-yaml}
	\end{minipage}

	Upon using the settings from \vref{lst:remove-para4-yaml} and running the commands
	\begin{widepage}
		\begin{commandshell}
latexindent.pl -m sl-stop.tex -o sl-stop4.tex -l remove-para4.yaml
latexindent.pl -m sl-stop.tex -o sl-stop4-command.tex -l=remove-para4.yaml,stop-command.yaml
latexindent.pl -m sl-stop.tex -o sl-stop4-comment.tex -l=remove-para4.yaml,stop-comment.yaml
    \end{commandshell}
	\end{widepage}
	we obtain the respective outputs in \crefrange{lst:sl-stop4}{lst:sl-stop4-comment}; notice in particular that:
	\begin{itemize}
		\item in \cref{lst:sl-stop4} the paragraph line break routine has included commands and comments;
		\item in \cref{lst:sl-stop4-command} the paragraph line break routine has \emph{stopped} at the
		      \texttt{emph} command, because in \cref{lst:stop-command-yaml} we have specified \texttt{commands}
		      to be \texttt{1}, and \texttt{emph} is at the beginning of a line;
		\item in \cref{lst:sl-stop4-comment} the paragraph line break routine has \emph{stopped}
		      at the comments, because in \cref{lst:stop-comment-yaml} we have specified \texttt{comments}
		      to be \texttt{1}, and the comment is at the beginning of a line.
	\end{itemize}
	In all outputs in \crefrange{lst:sl-stop4}{lst:sl-stop4-comment} we notice that the paragraph line break
	routine has stopped at \lstinline!\begin{myenv}! because, by default, \texttt{environments}
	is set to \texttt{1} in \vref{lst:paragraphsStopAt}.

	\cmhlistingsfromfile{demonstrations/sl-stop4.tex}{\texttt{sl-stop4.tex}}{lst:sl-stop4}
	\cmhlistingsfromfile{demonstrations/sl-stop4-command.tex}{\texttt{sl-stop4-command.tex}}{lst:sl-stop4-command}
	\cmhlistingsfromfile{demonstrations/sl-stop4-comment.tex}{\texttt{sl-stop4-comment.tex}}{lst:sl-stop4-comment}

\subsection{Poly-switches}
	Every other field in the \texttt{modifyLineBreaks} field uses poly-switches, and can take
	one of four integer values\footnote{You might like to associate one of the four circles in the logo with one of the four given values}:
	\begin{itemize}[font=\bfseries]
		\item[$-1$] \emph{remove mode}: line breaks before or after the \emph{<part of thing>} can be removed (assuming that \texttt{preserveBlankLines} is set to \texttt{0});
		\item[0] \emph{off mode}: line breaks will not be modified for the \emph{<part of thing>} under consideration;
		\item[1] \emph{add mode}: a line break will be added before or after the \emph{<part of thing>} under consideration, assuming that
		      there is not already a line break before or after the \emph{<part of thing>};
		\item[2] \emph{comment then add mode}: a comment symbol will be added, followed by a line break before or after the \emph{<part of thing>} under consideration, assuming that
		      there is not already a comment and line break before or after the \emph{<part of thing>}.
	\end{itemize}
	All poly-switches are \emph{off} by default; \texttt{latexindent.pl} searches first of all for per-name settings, and then followed by global per-thing settings.

\subsection{modifyLineBreaks for environments}\label{sec:modifylinebreaks-environments}
	We start by viewing a snippet of \texttt{defaultSettings.yaml} in \cref{lst:environments-mlb}; note that it contains \emph{global} settings (immediately
	after the \texttt{environments} field) and that \emph{per-name} settings are also allowed -- in the case of \cref{lst:environments-mlb}, settings
	for \texttt{equation*} have been specified. Note that all poly-switches are \emph{off} by default.

	\cmhlistingsfromfile[style=modifylinebreaksEnv]{../defaultSettings.yaml}[width=.8\linewidth,before=\centering,MLB-TCB]{\texttt{environments}}{lst:environments-mlb}

\subsubsection{Adding line breaks (poly-switches set to $1$ or $2$)}
	Let's begin with the simple example given in \cref{lst:env-mlb1-tex}; note that we have annotated key parts of the file using $\BeginStartsOnOwnLine$,
	$\BodyStartsOnOwnLine$, $\EndStartsOnOwnLine$ and $\EndFinishesWithLineBreak$, these will be related to fields specified in \cref{lst:environments-mlb}.

	\begin{cmhlistings}[escapeinside={(*@}{@*)}]{\texttt{env-mlb1.tex}}{lst:env-mlb1-tex}
before words(*@$\BeginStartsOnOwnLine$@*) \begin{myenv}(*@$\BodyStartsOnOwnLine$@*)body of myenv(*@$\EndStartsOnOwnLine$@*)\end{myenv}(*@$\EndFinishesWithLineBreak$@*) after words
\end{cmhlistings}

	Let's explore \texttt{BeginStartsOnOwnLine} and \texttt{BodyStartsOnOwnLine} in \cref{lst:env-mlb1,lst:env-mlb2}, and in particular,
	let's allow each of them in turn to take a value of $1$.

	\begin{minipage}{.45\textwidth}
		\cmhlistingsfromfile[style=yaml-LST]{demonstrations/env-mlb1.yaml}[MLB-TCB]{\texttt{env-mlb1.yaml}}{lst:env-mlb1}
	\end{minipage}
	\hfill
	\begin{minipage}{.45\textwidth}
		\cmhlistingsfromfile[style=yaml-LST]{demonstrations/env-mlb2.yaml}[MLB-TCB]{\texttt{env-mlb2.yaml}}{lst:env-mlb2}
	\end{minipage}

	After running the following commands,
	\begin{commandshell}
latexindent.pl -m env-mlb.tex -l env-mlb1.yaml
latexindent.pl -m env-mlb.tex -l env-mlb2.yaml
\end{commandshell}
	the output is as in \cref{lst:env-mlb-mod1,lst:env-mlb-mod2} respectively.

	\begin{widepage}
		\begin{minipage}{.57\linewidth}
			\cmhlistingsfromfile{demonstrations/env-mlb-mod1.tex}{\texttt{env-mlb.tex} using \cref{lst:env-mlb1}}{lst:env-mlb-mod1}
		\end{minipage}
		\hfill
		\begin{minipage}{.42\linewidth}
			\cmhlistingsfromfile{demonstrations/env-mlb-mod2.tex}{\texttt{env-mlb.tex} using \cref{lst:env-mlb2}}{lst:env-mlb-mod2}
		\end{minipage}
	\end{widepage}

	There are a couple of points to note:
	\begin{itemize}
		\item in \cref{lst:env-mlb-mod1} a line break has been added at the point denoted by $\BeginStartsOnOwnLine$ in \cref{lst:env-mlb1-tex}; no
		      other line breaks have been changed;
		\item in \cref{lst:env-mlb-mod2} a line break has been added at the point denoted by $\BodyStartsOnOwnLine$ in \cref{lst:env-mlb1-tex};
		      furthermore, note that the \emph{body} of \texttt{myenv} has received the appropriate (default) indentation.
	\end{itemize}

	Let's now change each of the \texttt{1} values in \cref{lst:env-mlb1,lst:env-mlb2} so that they are $2$ and
	save them into \texttt{env-mlb3.yaml} and \texttt{env-mlb4.yaml} respectively (see \cref{lst:env-mlb3,lst:env-mlb4}).

	\begin{minipage}{.45\textwidth}
		\cmhlistingsfromfile[style=yaml-LST]{demonstrations/env-mlb3.yaml}[MLB-TCB]{\texttt{env-mlb3.yaml}}{lst:env-mlb3}
	\end{minipage}
	\hfill
	\begin{minipage}{.45\textwidth}
		\cmhlistingsfromfile[style=yaml-LST]{demonstrations/env-mlb4.yaml}[MLB-TCB]{\texttt{env-mlb4.yaml}}{lst:env-mlb4}
	\end{minipage}

	Upon running  commands analogous to the above, we obtain \cref{lst:env-mlb-mod3,lst:env-mlb-mod4}.

	\begin{widepage}
		\begin{minipage}{.57\linewidth}
			\cmhlistingsfromfile{demonstrations/env-mlb-mod3.tex}{\texttt{env-mlb.tex} using \cref{lst:env-mlb3}}{lst:env-mlb-mod3}
		\end{minipage}
		\hfill
		\begin{minipage}{.42\linewidth}
			\cmhlistingsfromfile{demonstrations/env-mlb-mod4.tex}{\texttt{env-mlb.tex} using \cref{lst:env-mlb4}}{lst:env-mlb-mod4}
		\end{minipage}
	\end{widepage}

	Note that line breaks have been added as in \cref{lst:env-mlb-mod1,lst:env-mlb-mod2}, but this time a comment symbol
	has been added before adding the line break; in both cases, trailing horizontal
	space has been stripped before doing so.

	Let's explore \texttt{EndStartsOnOwnLine} and \texttt{EndFinishesWithLineBreak} in \cref{lst:env-mlb5,lst:env-mlb6},
	and in particular, let's allow each of them in turn to take a value of $1$.

	\begin{minipage}{.49\textwidth}
		\cmhlistingsfromfile[style=yaml-LST]{demonstrations/env-mlb5.yaml}[MLB-TCB]{\texttt{env-mlb5.yaml}}{lst:env-mlb5}
	\end{minipage}
	\hfill
	\begin{minipage}{.49\textwidth}
		\cmhlistingsfromfile[style=yaml-LST]{demonstrations/env-mlb6.yaml}[MLB-TCB]{\texttt{env-mlb6.yaml}}{lst:env-mlb6}
	\end{minipage}

	After running the following commands,
	\begin{commandshell}
latexindent.pl -m env-mlb.tex -l env-mlb5.yaml
latexindent.pl -m env-mlb.tex -l env-mlb6.yaml
\end{commandshell}
	the output is as in \cref{lst:env-mlb-mod5,lst:env-mlb-mod6}.

	\begin{widepage}
		\begin{minipage}{.42\linewidth}
			\cmhlistingsfromfile{demonstrations/env-mlb-mod5.tex}{\texttt{env-mlb.tex} using \cref{lst:env-mlb5}}{lst:env-mlb-mod5}
		\end{minipage}
		\hfill
		\begin{minipage}{.57\linewidth}
			\cmhlistingsfromfile{demonstrations/env-mlb-mod6.tex}{\texttt{env-mlb.tex} using \cref{lst:env-mlb6}}{lst:env-mlb-mod6}
		\end{minipage}
	\end{widepage}

	There are a couple of points to note:
	\begin{itemize}
		\item in \cref{lst:env-mlb-mod5} a line break has been added at the point denoted by $\EndStartsOnOwnLine$ in \vref{lst:env-mlb1-tex}; no
		      other line breaks have been changed and the \lstinline!\end{myenv}! statement has \emph{not} received indentation (as intended);
		\item in \cref{lst:env-mlb-mod6} a line break has been added at the point denoted by $\EndFinishesWithLineBreak$ in \vref{lst:env-mlb1-tex}.
	\end{itemize}

	Let's now change each of the \texttt{1} values in \cref{lst:env-mlb5,lst:env-mlb6} so that they are $2$ and
	save them into \texttt{env-mlb7.yaml} and \texttt{env-mlb8.yaml} respectively (see \cref{lst:env-mlb7,lst:env-mlb8}).

	\begin{minipage}{.49\textwidth}
		\cmhlistingsfromfile[style=yaml-LST]{demonstrations/env-mlb7.yaml}[MLB-TCB]{\texttt{env-mlb7.yaml}}{lst:env-mlb7}
	\end{minipage}
	\hfill
	\begin{minipage}{.49\textwidth}
		\cmhlistingsfromfile[style=yaml-LST]{demonstrations/env-mlb8.yaml}[MLB-TCB]{\texttt{env-mlb8.yaml}}{lst:env-mlb8}
	\end{minipage}

	Upon running  commands analogous to the above, we obtain \cref{lst:env-mlb-mod7,lst:env-mlb-mod8}.

	\begin{widepage}
		\begin{minipage}{.42\linewidth}
			\cmhlistingsfromfile{demonstrations/env-mlb-mod7.tex}{\texttt{env-mlb.tex} using \cref{lst:env-mlb7}}{lst:env-mlb-mod7}
		\end{minipage}
		\hfill
		\begin{minipage}{.57\linewidth}
			\cmhlistingsfromfile{demonstrations/env-mlb-mod8.tex}{\texttt{env-mlb.tex} using \cref{lst:env-mlb8}}{lst:env-mlb-mod8}
		\end{minipage}
	\end{widepage}

	Note that line breaks have been added as in \cref{lst:env-mlb-mod5,lst:env-mlb-mod6}, but this time a comment symbol
	has been added before adding the line break; in both cases, trailing horizontal
	space has been stripped before doing so.

	If you ask \texttt{latexindent.pl} to add a line break (possibly with a comment) using a poly-switch value of $1$ (or $2$),
	it will only do so if necessary. For example, if you process the file in \vref{lst:mlb2} using any of the YAML
	files presented so far in this section, it will be left unchanged.

	\begin{minipage}{.45\linewidth}
		\cmhlistingsfromfile{demonstrations/env-mlb2.tex}{\texttt{env-mlb2.tex}}{lst:mlb2}
	\end{minipage}
	\hfill
	\begin{minipage}{.45\linewidth}
		\cmhlistingsfromfile{demonstrations/env-mlb3.tex}{\texttt{env-mlb3.tex}}{lst:mlb3}
	\end{minipage}

	In contrast, the output from processing the file in \cref{lst:mlb3} will vary depending
	on the poly-switches used; in \cref{lst:env-mlb3-mod2} you'll see that the comment symbol after
	the \lstinline!\begin{myenv}! has been moved to the next line, as \texttt{BodyStartsOnOwnLine}
	is set to \texttt{1}. In \cref{lst:env-mlb3-mod4} you'll see that the comment has been accounted
	for correctly because \texttt{BodyStartsOnOwnLine} has been set to \texttt{2},
	and the comment symbol has \emph{not} been moved to its own line. You're encouraged to experiment
	with \cref{lst:mlb3} and by setting the other poly-switches considered so far to \texttt{2} in turn.

	\begin{minipage}{.45\linewidth}
		\cmhlistingsfromfile{demonstrations/env-mlb3-mod2.tex}{\texttt{env-mlb3.tex} using \vref{lst:env-mlb2}}{lst:env-mlb3-mod2}
	\end{minipage}
	\hfill
	\begin{minipage}{.45\linewidth}
		\cmhlistingsfromfile{demonstrations/env-mlb3-mod4.tex}{\texttt{env-mlb3.tex} using \vref{lst:env-mlb4}}{lst:env-mlb3-mod4}
	\end{minipage}

	The details of the discussion in this section have concerned \emph{global} poly-switches in the \texttt{environments} field;
	each switch can also be specified on a \emph{per-name} basis, which would take priority over the global values; with
	reference to \vref{lst:environments-mlb}, an example is shown for the \texttt{equation*} environment.

\subsubsection{Removing line breaks (poly-switches set to $-1$)}
	Setting poly-switches to $-1$ tells \texttt{latexindent.pl} to remove line breaks of the \emph{<part of the thing>}, if necessary. We will consider the
	example code given in \cref{lst:mlb4}, noting in particular the positions of
	the line break highlighters, $\BeginStartsOnOwnLine$, $\BodyStartsOnOwnLine$, $\EndStartsOnOwnLine$
	and $\EndFinishesWithLineBreak$, together with the associated YAML files in \crefrange{lst:env-mlb9}{lst:env-mlb12}.

	\begin{minipage}{.45\linewidth}
		\begin{cmhlistings}[escapeinside={(*@}{@*)}]{\texttt{env-mlb4.tex}}{lst:mlb4}
before words(*@$\BeginStartsOnOwnLine$@*)
\begin{myenv}(*@$\BodyStartsOnOwnLine$@*)
body of myenv(*@$\EndStartsOnOwnLine$@*)
\end{myenv}(*@$\EndFinishesWithLineBreak$@*)
after words
\end{cmhlistings}
	\end{minipage}%
	\hfill
	\begin{minipage}{.51\textwidth}
		\cmhlistingsfromfile[style=yaml-LST]{demonstrations/env-mlb9.yaml}[MLB-TCB]{\texttt{env-mlb9.yaml}}{lst:env-mlb9}

		\cmhlistingsfromfile[style=yaml-LST]{demonstrations/env-mlb10.yaml}[MLB-TCB]{\texttt{env-mlb10.yaml}}{lst:env-mlb10}

		\cmhlistingsfromfile[style=yaml-LST]{demonstrations/env-mlb11.yaml}[MLB-TCB]{\texttt{env-mlb11.yaml}}{lst:env-mlb11}

		\cmhlistingsfromfile[style=yaml-LST]{demonstrations/env-mlb12.yaml}[MLB-TCB]{\texttt{env-mlb12.yaml}}{lst:env-mlb12}
	\end{minipage}

	After running the commands
	\begin{commandshell}
latexindent.pl -m env-mlb4.tex -l env-mlb9.yaml
latexindent.pl -m env-mlb4.tex -l env-mlb10.yaml
latexindent.pl -m env-mlb4.tex -l env-mlb11.yaml
latexindent.pl -m env-mlb4.tex -l env-mlb12.yaml
\end{commandshell}

	we obtain the respective output in \crefrange{lst:env-mlb4-mod9}{lst:env-mlb4-mod12}.

	\begin{minipage}{.45\linewidth}
		\cmhlistingsfromfile{demonstrations/env-mlb4-mod9.tex}{\texttt{env-mlb4.tex} using \cref{lst:env-mlb9}}{lst:env-mlb4-mod9}
	\end{minipage}
	\hfill
	\begin{minipage}{.45\linewidth}
		\cmhlistingsfromfile{demonstrations/env-mlb4-mod10.tex}{\texttt{env-mlb4.tex} using \cref{lst:env-mlb10}}{lst:env-mlb4-mod10}
	\end{minipage}

	\begin{minipage}{.45\linewidth}
		\cmhlistingsfromfile{demonstrations/env-mlb4-mod11.tex}{\texttt{env-mlb4.tex} using \cref{lst:env-mlb11}}{lst:env-mlb4-mod11}
	\end{minipage}
	\hfill
	\begin{minipage}{.45\linewidth}
		\cmhlistingsfromfile{demonstrations/env-mlb4-mod12.tex}{\texttt{env-mlb4.tex} using \cref{lst:env-mlb12}}{lst:env-mlb4-mod12}
	\end{minipage}

	Notice that in
	\begin{itemize}
		\item \cref{lst:env-mlb4-mod9} the line break denoted by $\BeginStartsOnOwnLine$ in \cref{lst:mlb4} has been removed;
		\item \cref{lst:env-mlb4-mod10} the line break denoted by $\BodyStartsOnOwnLine$ in \cref{lst:mlb4} has been removed;
		\item \cref{lst:env-mlb4-mod11} the line break denoted by $\EndStartsOnOwnLine$ in \cref{lst:mlb4} has been removed;
		\item \cref{lst:env-mlb4-mod12} the line break denoted by $\EndFinishesWithLineBreak$ in \cref{lst:mlb4} has been removed.
	\end{itemize}
	We examined each of these cases separately for clarity of explanation, but you can combine all of the YAML
	settings in \crefrange{lst:env-mlb9}{lst:env-mlb12} into one file; alternatively, you could tell \texttt{latexindent.pl}
	to load them all by using the following command, for example
	\begin{widepage}
		\begin{commandshell}
latexindent.pl -m env-mlb4.tex -l env-mlb9.yaml,env-mlb10.yaml,env-mlb11.yaml,env-mlb12.yaml
\end{commandshell}
	\end{widepage}
	which gives the output in \vref{lst:env-mlb1-tex}.

	\paragraph{About trailing horizontal space}
		Recall that on \cpageref{yaml:removeTrailingWhitespace} we discussed the YAML field \texttt{removeTrailingWhitespace},
		and that it has two (binary) switches to determine if horizontal space should be removed \texttt{beforeProcessing} and \texttt{afterProcessing}.
		The \texttt{beforeProcessing} is particularly relevant when considering the \texttt{-m} switch; let's consider the
		file shown in \cref{lst:mlb5}, which highlights trailing spaces.

		\begin{minipage}{.45\linewidth}
			\begin{cmhlistings}[showspaces=true,escapeinside={(*@}{@*)}]{\texttt{env-mlb5.tex}}{lst:mlb5}
before words   (*@$\BeginStartsOnOwnLine$@*) 
\begin{myenv}           (*@$\BodyStartsOnOwnLine$@*)
body of myenv      (*@$\EndStartsOnOwnLine$@*) 
\end{myenv}     (*@$\EndFinishesWithLineBreak$@*)
after words
\end{cmhlistings}
		\end{minipage}
		\hfill
		\begin{minipage}{.45\linewidth}
			\cmhlistingsfromfile[style=yaml-LST]{demonstrations/removeTWS-before.yaml}[yaml-TCB]{\texttt{removeTWS-before.yaml}}{lst:removeTWS-before}
		\end{minipage}

		The output from the following commands
		\begin{widepage}
			\begin{commandshell}
latexindent.pl -m env-mlb5.tex -l env-mlb9.yaml,env-mlb10.yaml,env-mlb11.yaml,env-mlb12.yaml
latexindent.pl -m env-mlb5.tex -l env-mlb9.yaml,env-mlb10.yaml,env-mlb11.yaml,env-mlb12.yaml,removeTWS-before.yaml
\end{commandshell}
		\end{widepage}
		is shown, respectively, in \cref{lst:env-mlb5-modAll,lst:env-mlb5-modAll-remove-WS}; note that
		the trailing horizontal white space has been preserved (by default) in \cref{lst:env-mlb5-modAll}, while
		in \cref{lst:env-mlb5-modAll-remove-WS}, it has been removed using the switch specified in \cref{lst:removeTWS-before}.

		\begin{widepage}
			\cmhlistingsfromfile{demonstrations/env-mlb5-modAll.tex}{\texttt{env-mlb5.tex} using \crefrange{lst:env-mlb4-mod9}{lst:env-mlb4-mod12}}{lst:env-mlb5-modAll}

			\cmhlistingsfromfile{demonstrations/env-mlb5-modAll-remove-WS.tex}{\texttt{env-mlb5.tex} using \crefrange{lst:env-mlb4-mod9}{lst:env-mlb4-mod12} \emph{and} \cref{lst:removeTWS-before}}{lst:env-mlb5-modAll-remove-WS}
		\end{widepage}

	\paragraph{Blank lines}
		Now let's consider the file in \cref{lst:mlb6}, which contains blank lines.

		\begin{minipage}{.45\linewidth}
			\begin{cmhlistings}[escapeinside={(*@}{@*)}]{\texttt{env-mlb6.tex}}{lst:mlb6}
before words(*@$\BeginStartsOnOwnLine$@*)


\begin{myenv}(*@$\BodyStartsOnOwnLine$@*)


body of myenv(*@$\EndStartsOnOwnLine$@*)


\end{myenv}(*@$\EndFinishesWithLineBreak$@*)

after words
\end{cmhlistings}
		\end{minipage}%
		\hfill
		\begin{minipage}{.45\linewidth}
			\cmhlistingsfromfile[style=yaml-LST]{demonstrations/UnpreserveBlankLines.yaml}[MLB-TCB]{\texttt{UnpreserveBlankLines.yaml}}{lst:UnpreserveBlankLines}
		\end{minipage}

		Upon running the following commands
		\begin{widepage}
			\begin{commandshell}
latexindent.pl -m env-mlb6.tex -l env-mlb9.yaml,env-mlb10.yaml,env-mlb11.yaml,env-mlb12.yaml
latexindent.pl -m env-mlb6.tex -l env-mlb9.yaml,env-mlb10.yaml,env-mlb11.yaml,env-mlb12.yaml,UnpreserveBlankLines.yaml
\end{commandshell}
		\end{widepage}
		we receive the respective outputs in \cref{lst:env-mlb6-modAll,lst:env-mlb6-modAll-un-Preserve-Blank-Lines}. In
		\cref{lst:env-mlb6-modAll} we see that the multiple blank lines have each been condensed into one blank line,
		but that blank lines have \emph{not} been removed by the poly-switches -- this is because, by default, \texttt{preserveBlankLines}
		is set to \texttt{1}. By contrast, in \cref{lst:env-mlb6-modAll-un-Preserve-Blank-Lines}, we have allowed
		the poly-switches to remove blank lines because, in \cref{lst:UnpreserveBlankLines}, we have set \texttt{preserveBlankLines} to \texttt{0}.

		\begin{widepage}
			\begin{minipage}{.30\linewidth}
				\cmhlistingsfromfile{demonstrations/env-mlb6-modAll.tex}{\texttt{env-mlb6.tex} using \crefrange{lst:env-mlb4-mod9}{lst:env-mlb4-mod12}}{lst:env-mlb6-modAll}
			\end{minipage}
			\hfill
			\begin{minipage}{.65\linewidth}
				\cmhlistingsfromfile{demonstrations/env-mlb6-modAll-un-Preserve-Blank-Lines.tex}{\texttt{env-mlb6.tex} using \crefrange{lst:env-mlb4-mod9}{lst:env-mlb4-mod12} \emph{and} \cref{lst:UnpreserveBlankLines}}{lst:env-mlb6-modAll-un-Preserve-Blank-Lines}
			\end{minipage}
		\end{widepage}

\subsection{Poly-switches for other code blocks}
	Rather than repeat the examples shown for the environment code blocks (in \vref{sec:modifylinebreaks-environments}), we choose to detail the poly-switches for
	all other code blocks in \cref{tab:poly-switch-mapping}; note that each and every one of these poly-switches is \emph{off by default}, i.e, set to \texttt{0}. Note also that,
	by design, line breaks involving \texttt{verbatim}, \texttt{filecontents} and `comment-marked' code blocks (\vref{lst:alignmentmarkup}) can \emph{not} be
	modified using \texttt{latexindent.pl}.

	\begin{longtable}{m{.2\textwidth}@{\hspace{.75cm}}m{.35\textwidth}@{}m{.4\textwidth}}
		\caption{Poly-switch mappings for all code-block types}\label{tab:poly-switch-mapping}\\
		\toprule
		Code block & Sample & Poly-switch mapping \\
		\midrule
		environment &
		\begin{lstlisting}[escapeinside={(*@}{@*)},nolol=true]
before words(*@$\BeginStartsOnOwnLine$@*)
\begin{myenv}(*@$\BodyStartsOnOwnLine$@*)
body of myenv(*@$\EndStartsOnOwnLine$@*)
\end{myenv}(*@$\EndFinishesWithLineBreak$@*)
after words
  \end{lstlisting}
		&
		\begin{tabular}[t]{c@{~}l@{}}
			$\BeginStartsOnOwnLine$     & BeginStartsOnOwnLine     \\
			$\BodyStartsOnOwnLine$      & BodyStartsOnOwnLine      \\
			$\EndStartsOnOwnLine$       & EndStartsOnOwnLine       \\
			$\EndFinishesWithLineBreak$ & EndFinishesWithLineBreak \\
		\end{tabular}
		\\
		\cmidrule{2-3}
		ifelsefi &
		\begin{lstlisting}[escapeinside={(*@}{@*)},nolol=true]
before words(*@$\BeginStartsOnOwnLine$@*)
\if...(*@$\BodyStartsOnOwnLine$@*)
body of if statement(*@$\ElseStartsOnOwnLine$@*)
\else(*@$\ElseFinishesWithLineBreak$@*)
body of else statement(*@$\EndStartsOnOwnLine$@*)
\fi(*@$\EndFinishesWithLineBreak$@*)
after words
  \end{lstlisting}
		&
		\begin{tabular}[t]{c@{~}l@{}}
			$\BeginStartsOnOwnLine$      & IfStartsOnOwnLine         \\
			$\BodyStartsOnOwnLine$       & BodyStartsOnOwnLine       \\
			$\ElseStartsOnOwnLine$       & ElseStartsOnOwnLine       \\
			$\ElseFinishesWithLineBreak$ & ElseFinishesWithLineBreak \\
			$\EndStartsOnOwnLine$        & FiStartsOnOwnLine         \\
			$\EndFinishesWithLineBreak$  & FiFinishesWithLineBreak   \\
		\end{tabular}
		\\
		\cmidrule{2-3}
		optionalArguments &
		\begin{lstlisting}[escapeinside={(*@}{@*)},nolol=true]
...(*@$\BeginStartsOnOwnLine$@*)
[(*@$\BodyStartsOnOwnLine$@*)
body of opt arg(*@$\EndStartsOnOwnLine$@*)
](*@$\EndFinishesWithLineBreak$@*)
...
  \end{lstlisting}
		&
		\begin{tabular}[t]{c@{~}l@{}}
			$\BeginStartsOnOwnLine$     & LSqBStartsOnOwnLine\footnote{LSqB stands for Left Square Bracket} \\
			$\BodyStartsOnOwnLine$      & OptArgBodyStartsOnOwnLine                                         \\
			$\EndStartsOnOwnLine$       & RSqBStartsOnOwnLine                                               \\
			$\EndFinishesWithLineBreak$ & RSqBFinishesWithLineBreak                                         \\
		\end{tabular}
		\\
		\cmidrule{2-3}
		mandatoryArguments &
		\begin{lstlisting}[escapeinside={(*@}{@*)},nolol=true]
...(*@$\BeginStartsOnOwnLine$@*)
{(*@$\BodyStartsOnOwnLine$@*)
body of mand arg(*@$\EndStartsOnOwnLine$@*)
}(*@$\EndFinishesWithLineBreak$@*)
...
  \end{lstlisting}
		&
		\begin{tabular}[t]{c@{~}l@{}}
			$\BeginStartsOnOwnLine$     & LCuBStartsOnOwnLine\footnote{LCuB stands for Left Curly Brace} \\
			$\BodyStartsOnOwnLine$      & MandArgBodyStartsOnOwnLine                                     \\
			$\EndStartsOnOwnLine$       & RCuBStartsOnOwnLine                                            \\
			$\EndFinishesWithLineBreak$ & RCuBFinishesWithLineBreak                                      \\
		\end{tabular}
		\\
		\cmidrule{2-3}
		commands &
		\begin{lstlisting}[escapeinside={(*@}{@*)},morekeywords={mycommand},nolol=true,]
before words(*@$\BeginStartsOnOwnLine$@*)
\mycommand(*@$\BodyStartsOnOwnLine$@*)
(*@$\langle$\itshape{arguments}$\rangle$@*)
  \end{lstlisting}
		&
		\begin{tabular}[t]{c@{~}l@{}}
			$\BeginStartsOnOwnLine$ & CommandStartsOnOwnLine           \\
			$\BodyStartsOnOwnLine$  & CommandNameFinishesWithLineBreak \\
		\end{tabular}
		\\
		\cmidrule{2-3}
		namedGroupingBraces Brackets &
		\begin{lstlisting}[escapeinside={(*@}{@*)},morekeywords={myname},nolol=true,]
before words(*@$\BeginStartsOnOwnLine$@*)
myname(*@$\BodyStartsOnOwnLine$@*)
(*@$\langle$\itshape{braces/brackets}$\rangle$@*)
  \end{lstlisting}
		&
		\begin{tabular}[t]{c@{~}l@{}}
			$\BeginStartsOnOwnLine$ & NameStartsOnOwnLine       \\
			$\BodyStartsOnOwnLine$  & NameFinishesWithLineBreak \\
		\end{tabular}
		\\
		\cmidrule{2-3}
		keyEqualsValuesBraces\newline Brackets &
		\begin{lstlisting}[escapeinside={(*@}{@*)},morekeywords={key},nolol=true,]
before words(*@$\BeginStartsOnOwnLine$@*)
key(*@$\EqualsStartsOnOwnLine$@*)=(*@$\BodyStartsOnOwnLine$@*)
(*@$\langle$\itshape{braces/brackets}$\rangle$@*)
  \end{lstlisting}
		&
		\begin{tabular}[t]{c@{~}l@{}}
			$\BeginStartsOnOwnLine$  & KeyStartsOnOwnLine          \\
			$\EqualsStartsOnOwnLine$ & EqualsStartsOnOwnLine       \\
			$\BodyStartsOnOwnLine$   & EqualsFinishesWithLineBreak \\
		\end{tabular}
		\\
		\cmidrule{2-3}
		items &
		\begin{lstlisting}[escapeinside={(*@}{@*)},nolol=true]
before words(*@$\BeginStartsOnOwnLine$@*)
\item(*@$\BodyStartsOnOwnLine$@*)
...
  \end{lstlisting}
		&
		\begin{tabular}[t]{c@{~}l@{}}
			$\BeginStartsOnOwnLine$ & ItemStartsOnOwnLine       \\
			$\BodyStartsOnOwnLine$  & ItemFinishesWithLineBreak \\
		\end{tabular}
		\\
		\cmidrule{2-3}
		specialBeginEnd &
		\begin{lstlisting}[escapeinside={(*@}{@*)},nolol=true]
before words(*@$\BeginStartsOnOwnLine$@*)
\[(*@$\BodyStartsOnOwnLine$@*)
body of special(*@$\EndStartsOnOwnLine$@*)
\](*@$\EndFinishesWithLineBreak$@*)
after words
  \end{lstlisting}
		&
		\begin{tabular}[t]{c@{~}l@{}}
			$\BeginStartsOnOwnLine$     & SpecialBeginStartsOnOwnLine     \\
			$\BodyStartsOnOwnLine$      & SpecialBodyStartsOnOwnLine      \\
			$\EndStartsOnOwnLine$       & SpecialEndStartsOnOwnLine       \\
			$\EndFinishesWithLineBreak$ & SpecialEndFinishesWithLineBreak \\
		\end{tabular}
		\\
		\bottomrule
	\end{longtable}

% arara: pdflatex: {shell: yes, files: [latexindent]}
\subsection{Partnering \texttt{BodyStartsOnOwnLine} with argument-based poly-switches}
	Some poly-switches need to be partnered together; in particular, when line breaks involving the \emph{first} argument
	of a code block need to be accounted for using both \texttt{BodyStartsOnOwnLine} (or its equivalent, see \vref{tab:poly-switch-mapping})
	and \texttt{LCuBStartsOnOwnLine} for mandatory arguments, and \texttt{LSqBStartsOnOwnLine} for optional arguments.

	Let's begin with the code in \cref{lst:mycommand1} and the YAML settings in \cref{lst:mycom-mlb1}; with reference
	to \vref{tab:poly-switch-mapping}, the key \texttt{CommandNameFinishesWithLineBreak} is an alias for \texttt{BodyStartsOnOwnLine}.

	\cmhlistingsfromfile{demonstrations/mycommand1.tex}{\texttt{mycommand1.tex}}{lst:mycommand1}

	Upon running the command
	\begin{commandshell}
latexindent.pl -m -l=mycom-mlb1.yaml mycommand1.tex
\end{commandshell}
	we obtain \cref{lst:mycommand1-mlb1}; note that the \emph{second} mandatory argument beginning brace \lstinline!{! has had
	its leading line break removed, but that the \emph{first} brace has not.

	\begin{minipage}{.4\linewidth}
		\cmhlistingsfromfile{demonstrations/mycommand1-mlb1.tex}{\texttt{mycommand1.tex} using \cref{lst:mycom-mlb1}}{lst:mycommand1-mlb1}
	\end{minipage}
	\hfill
	\begin{minipage}{.55\linewidth}
		\cmhlistingsfromfile[style=yaml-LST]{demonstrations/mycom-mlb1.yaml}[MLB-TCB,width=\linewidth]{\texttt{mycom-mlb1.yaml}}{lst:mycom-mlb1}
	\end{minipage}

	Now let's change the YAML file so that it is as in \cref{lst:mycom-mlb2}; upon running the analogous command to that given above,
	we obtain \cref{lst:mycommand1-mlb2}; both beginning braces \lstinline!{! have had their leading line breaks removed.

	\begin{minipage}{.4\linewidth}
		\cmhlistingsfromfile{demonstrations/mycommand1-mlb2.tex}{\texttt{mycommand1.tex} using \cref{lst:mycom-mlb2}}{lst:mycommand1-mlb2}
	\end{minipage}
	\hfill
	\begin{minipage}{.55\linewidth}
		\cmhlistingsfromfile[style=yaml-LST]{demonstrations/mycom-mlb2.yaml}[MLB-TCB,width=\linewidth]{\texttt{mycom-mlb2.yaml}}{lst:mycom-mlb2}
	\end{minipage}

	Now let's change the YAML file so that it is as in \cref{lst:mycom-mlb3}; upon running the analogous command to that given above,
	we obtain \cref{lst:mycommand1-mlb3}.

	\begin{minipage}{.4\linewidth}
		\cmhlistingsfromfile{demonstrations/mycommand1-mlb3.tex}{\texttt{mycommand1.tex} using \cref{lst:mycom-mlb3}}{lst:mycommand1-mlb3}
	\end{minipage}
	\hfill
	\begin{minipage}{.55\linewidth}
		\cmhlistingsfromfile[style=yaml-LST]{demonstrations/mycom-mlb3.yaml}[MLB-TCB,width=\linewidth]{\texttt{mycom-mlb3.yaml}}{lst:mycom-mlb3}
	\end{minipage}

% arara: pdflatex: {shell: yes, files: [latexindent]}
\subsection{Conflicting poly-switches: sequential code blocks}
	It is very easy to have conflicting poly-switches; if we use the example from \vref{lst:mycommand1},
	and consider the YAML settings given in \cref{lst:mycom-mlb4}. The output from running
	\begin{commandshell}
latexindent.pl -m -l=mycom-mlb4.yaml mycommand1.tex
\end{commandshell}
	is given in \cref{lst:mycom-mlb4}.

	\begin{minipage}{.4\linewidth}
		\cmhlistingsfromfile{demonstrations/mycommand1-mlb4.tex}{\texttt{mycommand1.tex} using \cref{lst:mycom-mlb4}}{lst:mycommand1-mlb4}
	\end{minipage}
	\hfill
	\begin{minipage}{.55\linewidth}
		\cmhlistingsfromfile[style=yaml-LST]{demonstrations/mycom-mlb4.yaml}[MLB-TCB,width=\linewidth]{\texttt{mycom-mlb4.yaml}}{lst:mycom-mlb4}
	\end{minipage}

	Studying \cref{lst:mycom-mlb4}, we see that the two poly-switches are at opposition with one another:
	\begin{itemize}
		\item on the one hand, \texttt{LCuBStartsOnOwnLine} should \emph{not} start on its own line (as poly-switch is set to $-1$);
		\item on the other hand, \texttt{RCuBFinishesWithLineBreak} \emph{should} finish with a line break.
	\end{itemize}
	So, which should win the conflict? As demonstrated in \cref{lst:mycommand1-mlb4}, it is clear that \texttt{LCuBStartsOnOwnLine} won
	this conflict, and the reason is that \emph{the second argument was processed after the first} -- in general, the most recently-processed
	code block and associated poly-switch takes priority.

	We can explore this further by considering the YAML settings in \cref{lst:mycom-mlb5}; upon running the command
	\begin{commandshell}
latexindent.pl -m -l=mycom-mlb5.yaml mycommand1.tex
\end{commandshell}
	we obtain the output given in \cref{lst:mycommand1-mlb5}.

	\begin{minipage}{.4\linewidth}
		\cmhlistingsfromfile{demonstrations/mycommand1-mlb5.tex}{\texttt{mycommand1.tex} using \cref{lst:mycom-mlb5}}{lst:mycommand1-mlb5}
	\end{minipage}
	\hfill
	\begin{minipage}{.55\linewidth}
		\cmhlistingsfromfile[style=yaml-LST]{demonstrations/mycom-mlb5.yaml}[MLB-TCB,width=\linewidth]{\texttt{mycom-mlb5.yaml}}{lst:mycom-mlb5}
	\end{minipage}

	As previously, the most-recently-processed code block takes priority -- as before, the second (i.e, \emph{last}) argument. Exploring this
	further, we consider the YAML settings in \cref{lst:mycom-mlb6}, which give associated output in \cref{lst:mycommand1-mlb6}.

	\begin{minipage}{.4\linewidth}
		\cmhlistingsfromfile{demonstrations/mycommand1-mlb6.tex}{\texttt{mycommand1.tex} using \cref{lst:mycom-mlb6}}{lst:mycommand1-mlb6}
	\end{minipage}
	\hfill
	\begin{minipage}{.55\linewidth}
		\cmhlistingsfromfile[style=yaml-LST]{demonstrations/mycom-mlb6.yaml}[MLB-TCB,width=\linewidth]{\texttt{mycom-mlb6.yaml}}{lst:mycom-mlb6}
	\end{minipage}

	Note that a \lstinline!%! \emph{has} been added to the trailing first \lstinline!}!; this is because:
	\begin{itemize}
		\item while processing the \emph{first} argument, the trailing line break has been removed (\texttt{RCuBFinishesWithLineBreak} set to $-1$);
		\item while processing the \emph{second} argument, \texttt{latexindent.pl} finds that it does \emph{not} begin on its own line, and so
		      because \texttt{LCuBStartsOnOwnLine} is set to $2$, it adds a comment, followed by a line break.
	\end{itemize}

\subsection{Conflicting poly-switches: nested code blocks}
	Now let's consider an example when nested code blocks have conflicting poly-switches; we'll use the code in \cref{lst:nested-env},
	noting that it contains nested environments.

	\cmhlistingsfromfile{demonstrations/nested-env.tex}{\texttt{nested-env.tex}}{lst:nested-env}

	Let's use the YAML settings given in \cref{lst:nested-env-mlb1-yaml}, which upon running the command
	\begin{commandshell}
latexindent.pl -m -l=nested-env-mlb1.yaml nested-env.tex
        \end{commandshell}
	gives the output in \cref{lst:nested-env-mlb1}.

	\begin{minipage}{.45\linewidth}
		\cmhlistingsfromfile{demonstrations/nested-env-mlb1.tex}{\texttt{nested-env.tex} using \cref{lst:nested-env-mlb1}}{lst:nested-env-mlb1}
	\end{minipage}
	\hfill
	\begin{minipage}{.55\linewidth}
		\cmhlistingsfromfile[style=yaml-LST]{demonstrations/nested-env-mlb1.yaml}[MLB-TCB,width=\linewidth]{\texttt{nested-env-mlb1.yaml}}{lst:nested-env-mlb1-yaml}
	\end{minipage}

	In \cref{lst:nested-env-mlb1}, let's first of all note that both environments have received the appropriate (default) indentation; secondly,
	note that the poly-switch \texttt{EndStartsOnOwnLine} appears to have won the conflict, as \lstinline!\end{one}! has had its leading line break removed.

	To understand it, let's talk about the three basic phases of \texttt{latexindent.pl}:
	\begin{enumerate}
		\item Phase 1: packing, in which code blocks are replaced with unique ids, working from \emph{the inside to the outside}, and then sequentially -- for example, in
		      \cref{lst:nested-env}, the \texttt{two} environment is found \emph{before} the \texttt{one} environment; if the -m switch is active, then during this phase:
		      \begin{itemize}
			      \item line breaks at the beginning of the \texttt{body} can be added (if \texttt{BodyStartsOnOwnLine} is $1$ or $2$) or removed (if \texttt{BodyStartsOnOwnLine} is $-1$);
			      \item line breaks at the end of the body can be added (if \texttt{EndStartsOnOwnLine} is $1$ or $2$) or removed (if \texttt{EndStartsOnOwnLine} is $-1$);
			      \item line breaks after the end statement can be added (if \texttt{EndFinishesWithLineBreak} is $1$ or $2$).
		      \end{itemize}
		\item Phase 2: indentation, in which white space is added to the begin, body, and end statements;
		\item Phase 3: unpacking, in which unique ids are replaced by their \emph{indented} code blocks; if the -m switch is active, then during this phase,
		      \begin{itemize}
			      \item line breaks before \texttt{begin} statements can be added or removed (depending upon \texttt{BeginStartsOnOwnLine});
			      \item line breaks after \emph{end} statements can be removed but \emph{NOT} added (see \texttt{EndFinishesWithLineBreak}).
		      \end{itemize}
	\end{enumerate}

	With reference to \cref{lst:nested-env-mlb1}, this means that during Phase 1:
	\begin{itemize}
		\item the \texttt{two} environment is found first, and
		      the line break ahead of the \lstinline!\end{two}! statement is removed because \texttt{EndStartsOnOwnLine} is set to $-1$. Importantly,
		      because, \emph{at this stage}, \lstinline!\end{two}! \emph{does} finish with a line break, \texttt{EndFinishesWithLineBreak} causes
		      no action.
		\item next, the \texttt{one} environment is found; the line break ahead of \lstinline!\end{one}! is removed because \texttt{EndStartsOnOwnLine}
		      is set to $-1$.
	\end{itemize}
	The indentation is done in Phase 2; in Phase 3 \emph{there is no option to add a line break after the \lstinline!end! statements}.
	We can justify this by remembering that during Phase 3, the \texttt{one} environment will be found and processed first, followed
	by the \texttt{two} environment. If the \texttt{two} environment were to add a line break after the \lstinline!\end{two}! statement, then
	\texttt{latexindent.pl} would have no way of knowing how much indentation to add to the subsequent text (in this case, \lstinline!\end{one}!).

	We can explore this further using the poly-switches in \cref{lst:nested-env-mlb2}; upon running the command
	\begin{commandshell}
latexindent.pl -m -l=nested-env-mlb2.yaml nested-env.tex
        \end{commandshell}
	we obtain the output given in \cref{lst:nested-env-mlb2-output}.

	\begin{minipage}{.45\linewidth}
		\cmhlistingsfromfile{demonstrations/nested-env-mlb2.tex}{\texttt{nested-env.tex} using \cref{lst:nested-env-mlb2}}{lst:nested-env-mlb2-output}
	\end{minipage}
	\hfill
	\begin{minipage}{.55\linewidth}
		\cmhlistingsfromfile[style=yaml-LST]{demonstrations/nested-env-mlb2.yaml}[MLB-TCB,width=\linewidth]{\texttt{nested-env-mlb2.yaml}}{lst:nested-env-mlb2}
	\end{minipage}

	During Phase 1:
	\begin{itemize}
		\item the \texttt{two} environment is found first, and
		      the line break ahead of the \lstinline!\end{two}! statement is not changed because \texttt{EndStartsOnOwnLine} is set to $1$.
		      Importantly, because, \emph{at this stage}, \lstinline!\end{two}! \emph{does} finish with a line break, \texttt{EndFinishesWithLineBreak} causes
		      no action.
		\item next, the \texttt{one} environment is found; the line break ahead of \lstinline!\end{one}! is already present, and no action is needed.
	\end{itemize}
	The indentation is done in Phase 2, and then in Phase 3, the \texttt{one} environment is found and processed first, followed by
	the \texttt{two} environment. \emph{At this stage}, the \texttt{two} environment finds \texttt{EndFinishesWithLineBreak} is $-1$, so it removes
	the trailing line break; remember, at this point, \texttt{latexindent.pl} has completely finished with the \texttt{one} environment.

% arara: pdflatex: {shell: yes, files: [latexindent]}
\section{Conclusions and known limitations}\label{sec:knownlimitations}
 There are a number of known limitations of the script, and almost certainly quite a
 few that are \emph{unknown}!

 For example, \texttt{latexindent.pl} will not indent the following code correctly,
 because of the unmatched \lstinline![!. I'm hopeful to be able to resolve this
 issue in a future version.

 \begin{lstlisting}[,nolol=true,]
\parbox{
\@ifnextchar[{\@assignmentwithcutoff}{\@assignmentnocutoff}
}
\end{lstlisting}

 The main other limitation is to do with the alignment routine of environments/commands that contain
 delimiters which are specified in \texttt{lookForAlignDelims}.

 The routine works well for `standard' blocks of code that have the same number of \lstinline!&!
 per line, but it will not do anything for lines that do not -- such examples
 include \texttt{tabular} environments that use \lstinline!\multicolumn! or
 perhaps spread cell contents across multiple lines.  For each alignment block (\texttt{tabular},
 \texttt{align}, etc) \texttt{latexindent.pl} first of all makes a record
 of the maximum number of \lstinline!&!; if each row does not have that
 number of \lstinline!&! then it will not try to format that row. Details
 will be given in \texttt{indent.log} assuming that \texttt{trace} mode
 is active.

 You can run \texttt{latexindent} on \texttt{.sty}, \texttt{.cls} and any file types
 that you specify in \lstinline[breaklines=true]!fileExtensionPreference! (see \vref{lst:fileExtensionPreference});
 if you find a case in which the script struggles, please feel free
 to report it at \cite{latexindent-home}, and
 in the meantime, consider using a \texttt{noIndentBlock} (see \cpageref{lst:noIndentBlockdemo}).

 I hope that this script is useful to some; if you find an example where the
 script does not behave as you think it should, the best way to contact me is to
 report an issue on \cite{latexindent-home}; otherwise, feel free to find me on
 the \url{http://tex.stackexchange.com/users/6621/cmhughes}.

\nocite{*}
\section{References}
 \printbibliography[heading=subbibnumbered,title={External links},notkeyword=contributor]
 \printbibliography[env=specialbib,heading=subbibnumbered,title={Contributors\label{sec:contributors}},keyword=contributor]

% arara: pdflatex: {shell: yes, files: [latexindent]}
\appendix
	\section{Required \texttt{Perl} modules}\label{sec:requiredmodules}
	 If you intend to use \texttt{latexindent.pl} and \emph{not} one of the supplied standalone executable files, then you will need a few standard Perl modules -- if you can run the
	 minimum code in \cref{lst:helloworld} (\texttt{perl helloworld.pl}) then you will be able to run \texttt{latexindent.pl}, otherwise you may
	 need to install the missing modules.

	 \begin{cmhlistings}[language=Perl]{\texttt{helloworld.pl}}{lst:helloworld}
#!/usr/bin/perl

use strict;
use warnings;
use utf8;
use PerlIO::encoding;
use Unicode::GCString;
use open ':std', ':encoding(UTF-8)';
use Text::Wrap;
use FindBin;
use YAML::Tiny;
use File::Copy;
use File::Basename;
use File::HomeDir;
use Getopt::Long;
use Data::Dumper;

print "hello world";
exit;
\end{cmhlistings}
	 Installing the modules given in \cref{lst:helloworld} will vary depending on your
	 operating system and \texttt{Perl} distribution. For example, Ubuntu users
	 might visit the software center, or else run
	 \begin{commandshell}
sudo perl -MCPAN -e 'install "File::HomeDir"'
 \end{commandshell}

	 Linux users may be interested in exploring Perlbrew \cite{perlbrew}; possible installation and setup
	 options follow for Ubuntu (other distributions will need slightly different commands).
	 \begin{commandshell}
sudo apt-get install perlbrew
perlbrew install perl-5.20.1
perlbrew switch perl-5.20.1
sudo apt-get install curl
curl -L http://cpanmin.us | perl - App::cpanminus
cpanm YAML::Tiny
cpanm File::HomeDir
\end{commandshell}

	 Strawberry Perl users on Windows might use
	 \texttt{CPAN client}. All of the modules are readily available on CPAN \cite{cpan}.

	 \texttt{indent.log} will contain details of the location
	 of the Perl modules on your system.  \texttt{latexindent.exe} is a standalone
	 executable for Windows (and therefore does not require a Perl distribution) and caches copies of the Perl modules onto your system; if you
	 wish to see where they are cached, use the  \texttt{trace} option, e.g
	 \begin{dosprompt}
latexindent.exe -t myfile.tex
 \end{dosprompt}

	\section{Updating the \texttt{path} variable}\label{sec:updating-path}
	 \texttt{latexindent.pl} has a few scripts (available at \cite{latexindent-home}) that can update the \texttt{path} variables\footnote{Thanks to \cite{jasjuang} for this feature!}. If you're
	 on a Linux or Mac machine, then you'll want \texttt{CMakeLists.txt} from \cite{latexindent-home}.
	\subsection{Add to path for Linux}
		To add \texttt{latexindent.pl} to the path for Linux, follow these steps:
		\begin{enumerate}
			\item download  \texttt{latexindent.pl} and its associated modules, \texttt{defaultSettings.yaml},
			      to your chosen directory from \cite{latexindent-home} ;
			\item within your directory, create a directory called \texttt{path-helper-files} and
			      download \texttt{CMakeLists.txt} and \lstinline!cmake_uninstall.cmake.in!
			      from \cite{latexindent-home}/path-helper-files to this directory;
			\item run
			      \begin{commandshell}
ls /usr/local/bin
          \end{commandshell}
			      to see what is \emph{currently} in there;
			\item run the following commands
			      \begin{commandshell}
sudo apt-get install cmake
sudo apt-get update && sudo apt-get install build-essential
mkdir build && cd build
cmake ../path-helper-files
sudo make install
\end{commandshell}
			\item run
			      \begin{commandshell}
ls /usr/local/bin
          \end{commandshell}
			      again to check that \texttt{latexindent.pl}, its modules and \texttt{defaultSettings.yaml} have been added.
		\end{enumerate}
		To \emph{remove} the files, run
		\begin{commandshell}
sudo make uninstall}.
    \end{commandshell}
	\subsection{Add to path for Windows}
		To add \texttt{latexindent.exe} to the path for Windows, follow these steps:
		\begin{enumerate}
			\item download  \texttt{latexindent.exe}, \texttt{defaultSettings.yaml},  \texttt{add-to-path.bat}
			      from \cite{latexindent-home} to your chosen directory;
			\item open a command prompt and run the following command to see what is \emph{currently} in your \lstinline!%path%! variable;
			      \begin{dosprompt}
echo %path%
          \end{dosprompt}
			\item right click on \texttt{add-to-path.bat} and \emph{Run as administrator};
			\item log out, and log back in;
			\item open a command prompt and run
			      \begin{dosprompt}
echo %path%
          \end{dosprompt}
			      to check that the appropriate directory has been added to your \lstinline!%path%!.
		\end{enumerate}
		To \emph{remove} the directory from your \lstinline!%path%!, run \texttt{remove-from-path.bat} as administrator.

	\section{Differences from Version 2.2 to 3.0}\label{app:differences}
	 There are a few (small) changes to the interface when comparing Version 2.2 to Version 3.0.
	 Explicitly, in previous versions you might have run, for example,
	 \begin{commandshell}
latexindent.pl -o myfile.tex outputfile.tex
 \end{commandshell}
	 whereas in Version 3.0 you would run any of the following, for example,
	 \begin{commandshell}
latexindent.pl -o=outputfile.tex myfile.tex
latexindent.pl -o outputfile.tex myfile.tex
latexindent.pl myfile.tex -o outputfile.tex 
latexindent.pl myfile.tex -o=outputfile.tex 
latexindent.pl myfile.tex -outputfile=outputfile.tex 
latexindent.pl myfile.tex -outputfile outputfile.tex 
 \end{commandshell}
	 noting that the \emph{output} file is given \emph{next to} the \texttt{-o} switch.

	 The fields given in \cref{lst:obsoleteYaml} are \emph{obsolete} from Version 3.0
	 onwards.
	 \begin{yaml}[style=yaml-LST,numbers=none]{Obsolete YAML fields from Version 3.0}[colframe=white!25!red,colbacktitle=white!75!red,colback=white!90!red,]{lst:obsoleteYaml}
alwaysLookforSplitBrackets
alwaysLookforSplitBrackets
checkunmatched
checkunmatchedELSE
checkunmatchedbracket
constructIfElseFi
\end{yaml}

	 There is a slight difference when specifying indentation after headings; specifically,
	 we now write \texttt{indentAfterThisHeading} instead of \texttt{indent}. See \cref{lst:indentAfterThisHeadingOld,lst:indentAfterThisHeadingNew}

	 \begin{minipage}{.45\textwidth}
		 \begin{yaml}[style=yaml-LST,numbers=none]{\texttt{indentAfterThisHeading} in Version 2.2}{lst:indentAfterThisHeadingOld}
indentAfterHeadings:
    part:
       indent: 0
       level: 1
\end{yaml}
	 \end{minipage}%
	 \hfill
	 \begin{minipage}{.45\textwidth}
		 \begin{yaml}[style=yaml-LST,numbers=none]{\texttt{indentAfterThisHeading} in Version 3.0}{lst:indentAfterThisHeadingNew}
indentAfterHeadings:
    part:
       indentAfterThisHeading: 0
       level: 1
\end{yaml}
	 \end{minipage}%

	 To specify \texttt{noAdditionalIndent} for display-math environments in Version 2.2, you would write YAML
	 as in \cref{lst:noAdditionalIndentOld}; as of Version 3.0, you would write YAML as in \cref{lst:indentAfterThisHeadingNew1}
	 or, if you're using \texttt{-m} switch, \cref{lst:indentAfterThisHeadingNew2}.

	 \begin{minipage}{.45\textwidth}
		 \begin{yaml}[style=yaml-LST,numbers=none]{\texttt{noAdditionalIndent} in Version 2.2}{lst:noAdditionalIndentOld}
noAdditionalIndent:
    \[: 0
    \]: 0
\end{yaml}
	 \end{minipage}%
	 \hfill
	 \begin{minipage}{.45\textwidth}
		 \begin{yaml}[style=yaml-LST,numbers=none]{\texttt{noAdditionalIndent} for \texttt{displayMath} in Version 3.0}{lst:indentAfterThisHeadingNew1}
specialBeginEnd:
    displayMath:
        begin: '\\\['
        end: '\\\]'
        lookForThis: 0
\end{yaml}

		 \begin{yaml}[style=yaml-LST,numbers=none]{\texttt{noAdditionalIndent} for \texttt{displayMath}  in Version 3.0}{lst:indentAfterThisHeadingNew2}
noAdditionalIndent:
    displayMath: 1
\end{yaml}
	 \end{minipage}%

	 \mbox{}\hfill     \begin{minipage}{.25\textwidth}
		 \hrule

		 \hfill\itshape End

	 \end{minipage}

\end{document}

\subsection{The phases of \texttt{latexindent.pl}}
	With these rules in mind, let's study a few test cases:

	latexindent.pl environments-line-break-conflict.tex -s -t -m -o environments-line-break-conflict-mod1.tex -l=env-conflicts-mod1.yaml
	latexindent.pl environments-line-break-conflict-nested.tex -s -t -m -o environments-line-break-conflict-nested-mod-2.tex -l=env-conflicts-mod2.yaml
	latexindent.pl environments-line-break-conflict-nested.tex -s -t -m -o environments-line-break-conflict-nested-mod-3.tex -l=env-conflicts-mod3.yaml
	environments-first-opt-args.tex, see all of the different examples in test-cases.sh
	environments-second-opt-args.tex provides some interesting cases too

	The \lstinline!\fi! command knows to insert a space, so as to give, for example, \lstinline!\fi! text, and avoid things such as \lstinline!\fitext!

	from yaml

